% !TeX program = pdflatex
% If you use biblatex/biber (recommended), compile with:
% pdflatex -> biber -> pdflatex -> pdflatex

\documentclass[11pt,a4paper,twoside,openany]{report}

% -------------------------------------------------
% Page layout / typography
% -------------------------------------------------
\usepackage[a4paper,margin=30mm,headheight=14pt]{geometry}
\usepackage{microtype}              % better spacing/kerning
\usepackage[T1]{fontenc}
\usepackage[utf8]{inputenc}         % ok for pdflatex; remove if using lualatex/xelatex
\usepackage{xcolor}
\usepackage{lmodern}                % modern Latin Modern fonts
\usepackage{setspace}               % optional: \onehalfspacing etc.

% -------------------------------------------------
% Graphics / tables / lists
% -------------------------------------------------
\usepackage{graphicx}
\usepackage{booktabs}               % professional tables
\usepackage{array}
\usepackage{multirow}
\usepackage{caption}
\usepackage{subcaption}
\usepackage{enumitem}

% -------------------------------------------------
% Math (standard praxis)
% -------------------------------------------------
\usepackage{amsmath,amssymb,amsthm} % AMS core
\usepackage{mathtools}             % extends/fixes amsmath
\usepackage{stmaryrd}              % \llbracket, \rrbracket for stochastic intervals

% Upright differential (use as \int f(x)\,\diff x)
\newcommand{\diff}{\mathop{}\!\mathrm{d}}

% Operators
\DeclareMathOperator*{\arginf}{arginf}
\DeclareMathOperator*{\argmin}{arg\,min}

\usepackage{bm}                    % bold math symbols
\usepackage{dsfont}                % \mathds{1} indicator
\usepackage{mathrsfs}              % script font \mathscr
\usepackage{bbm}                   % alternative blackboard bold, incl. \mathbbm{1}
\usepackage{siunitx}               % units/numbers (also nice in finance tables)

% -------------------------------------------------
% References / hyperlinks (standard modern setup)
% -------------------------------------------------
\usepackage{xcolor}
\usepackage{hyperref}
\hypersetup{
  colorlinks=true,
  linkcolor=blue,
  citecolor=blue,
  urlcolor=blue,
  pdftitle={Master Project},
  pdfauthor={Your Name},
  pdfsubject={Mathematical Finance},
  pdfkeywords={Mathematical Finance, Stochastic Calculus, ...}
}

% cleveref + shared theorem counters:
% If environments share the "theorem" counter, \cref may otherwise call everything "Theorem".
% aliascnt creates alias counters so cleveref can distinguish Remark/Lemma/etc.
\usepackage{aliascnt}

% cleveref must be loaded AFTER hyperref
\usepackage[nameinlink,noabbrev]{cleveref}

% -------------------------------------------------
% Theorems / definitions (standard mathematical setup)
% -------------------------------------------------
% Numbering: Theorem 2.3, Lemma 2.4, ... within each chapter.
\theoremstyle{plain} % italic body, bold heading
\newtheorem{theorem}{Theorem}[chapter]

\newaliascnt{proposition}{theorem}
\newtheorem{proposition}[proposition]{Proposition}
\aliascntresetthe{proposition}

\newaliascnt{lemma}{theorem}
\newtheorem{lemma}[lemma]{Lemma}
\aliascntresetthe{lemma}

\newaliascnt{corollary}{theorem}
\newtheorem{corollary}[corollary]{Corollary}
\aliascntresetthe{corollary}

\theoremstyle{definition} % upright body
\newaliascnt{definition}{theorem}
\newtheorem{definition}[definition]{Definition}
\aliascntresetthe{definition}

\newaliascnt{assumption}{theorem}
\newtheorem{assumption}[assumption]{Assumption}
\aliascntresetthe{assumption}

\theoremstyle{remark} % upright body, italic heading
\newaliascnt{remark}{theorem}
\newtheorem{remark}[remark]{Remark}
\aliascntresetthe{remark}

% Unnumbered variants (use sparingly)
\newtheorem*{theorem*}{Theorem}
\newtheorem*{lemma*}{Lemma}
\newtheorem*{definition*}{Definition}

% Cleveref names for theorem-like environments
\crefname{theorem}{Theorem}{Theorems}
\Crefname{theorem}{Theorem}{Theorems}
\crefname{proposition}{Proposition}{Propositions}
\Crefname{proposition}{Proposition}{Propositions}
\crefname{lemma}{Lemma}{Lemmas}
\Crefname{lemma}{Lemma}{Lemmas}
\crefname{corollary}{Corollary}{Corollaries}
\Crefname{corollary}{Corollary}{Corollaries}
\crefname{definition}{Definition}{Definitions}
\Crefname{definition}{Definition}{Definitions}
\crefname{assumption}{Assumption}{Assumptions}
\Crefname{assumption}{Assumption}{Assumptions}
\crefname{remark}{Remark}{Remarks}
\Crefname{remark}{Remark}{Remarks}

% -------------------------------------------------
% Bibliography (biblatex + biber)
% -------------------------------------------------
\usepackage{csquotes} % recommended with biblatex
\usepackage[
  backend=biber,
  style=authoryear,
  dashed=false, % don't replace repeated author/editor with an em-dash
  maxcitenames=2,
  maxbibnames=99
]{biblatex}
\addbibresource{references.bib}

% -------------------------------------------------
% Appendix + PDF inclusion (optional)
% -------------------------------------------------
\usepackage{appendix}
\usepackage{pdfpages}

% -------------------------------------------------
% Nomenclature / symbols list (optional but common in math finance)
% -------------------------------------------------
\usepackage[intoc]{nomencl}
\makenomenclature
% Example: \nomenclature{$W_t$}{Standard Brownian motion}

% -------------------------------------------------
% Header / footer (simple, thesis-friendly)
% -------------------------------------------------
\usepackage{fancyhdr}
\pagestyle{fancy}
\fancyhf{}
\fancyhead[LE,RO]{\thepage}
\fancyhead[LO]{\nouppercase{\rightmark}}
\fancyhead[RE]{\nouppercase{\leftmark}}
\renewcommand{\headrulewidth}{0.4pt}

% -------------------------------------------------
% Table of contents depth (DETAILED ToC)
% -------------------------------------------------
% tocdepth: 0=part, 1=chapter, 2=section, 3=subsection, 4=subsubsection
\setcounter{tocdepth}{3}  % show down to \subsection in ToC
\setcounter{secnumdepth}{3} % number down to \subsection (or 4 for subsubsection)

% -------------------------------------------------
% Title metadata
% -------------------------------------------------
\newcommand{\thesistitle}{Aspects of Neural Networks in Continuous-Time Mathematical Finance}
\newcommand{\thesisauthor}{Joël Wäspi}
\newcommand{\thesisadvisor}{Martin Schweizer}
\newcommand{\thesisaffil}{ETH Zürich}
\newcommand{\thesisdate}{\today}

\begin{document}

% =================================================
% Front matter (Roman numerals)
% =================================================
\pagenumbering{roman}

\begin{titlepage}
  \centering
  \vspace*{2cm}
  {\Large \thesisaffil\par}
  \vspace{1.5cm}
  {\huge\bfseries \thesistitle\par}
  \vspace{1.0cm}
  {\Large Master Project\par}
  \vspace{2.0cm}
  {\Large \thesisauthor\par}
  \vspace{0.5cm}
  {\large Advisor: \thesisadvisor\par}
  \vfill
  {\large \thesisdate\par}
\end{titlepage}

\chapter*{Abstract}
\addcontentsline{toc}{chapter}{Abstract}
% Write later. Keep short (often 150--300 words).
\noindent
\textit{Placeholder abstract.}

\chapter*{Acknowledgements}
\addcontentsline{toc}{chapter}{Acknowledgements}
% Optional
\noindent
\textit{Placeholder acknowledgements.}

\tableofcontents


\printnomenclature
\cleardoublepage

% =================================================
% Main matter (Arabic numerals)
% =================================================
\pagenumbering{arabic}

% =================================================
% Detailed scaffold for a Mathematical Finance project
% (Edit freely before you start writing)
% =================================================

\chapter{Introduction}
\section{Motivation and context}
\section{Outline of the thesis}

\chapter{Hedging}\label{sec:hedging}
In order to explain the concept of hedging, we first convey a conceptual understanding. This will motivate the probabilistic set-up that leads us to a well defined mathematical problem that can be studied analytically.
We formulate hedging from the perspective of a seller of a European call option. The latter gives the buyer the option to buy the underlying asset at the time of maturity at the strike price \(K\). the strike price is known at the time the option was acquired. Crucially, the owner of the option can only exercise the contract at maturity, not before. At maturity, the owner of the option will exercise his right to buy if the price of the underlying asset \(X\) at maturity is higher than the strike price \(K\). If the price of the underlying asset at maturity is below the strike price, the owner of the option doesn't exercise his right to buy. Thus, the payoff of the option for the owner is given by 
\begin{align}
    H = \max(X_T-K,0)\;.
\end{align}
Note that the price \(S\) at maturity of the underlying asset is modeled to be stochastic whereas the strike price is deterministic. Hence, for a given state of the world \(\omega\), the random payoff of the European call option is \(H(\omega)=\max(X_T(\omega)-K,0)\).
\bigbreak
From the perspective of the seller of the European call option, the gain of the buyer of the European call option is the amount of money he is obliged to pay. We call this a \textit{liability}. Intuitively, the seller will ask for a price for the option equal to the expected payoff of the option. In order for the seller to be able to pay out the payoff of the call option, he might be interested in investing in the stock market himself and hoping to make at least a profit that enables him to compensate for his liability. This is what we call \textit{Hedging}. The questions arising are hence: Is there a trading strategy that will give me, almost surely, a gain at least equal to the payoff of the liability? If not, how close can I get? What does it mean to be close? What is the right price I should ask for the option? In order to give satisfactory answers, we will introduce a more technical framework, allowing us to give a precise definition of the problem at hand.
\section{Self-financing strategies and trading gains}
In order to model the stochastic nature of the hedging problem, we naturally want to start with a filtered probability space \((\Omega,\mathcal{F},\mathbb{F},\mathbb{P})\) with a filtration \(\mathbb{F}=(\mathcal{F}_t)_{t\in I}\), for \(I=[0,T]\) subset of \(\mathbb{R}_+=[0,\infty)\) with \(T\in\mathbb{R}_+\) or a subset of \(\mathbb{N}\) with \(T\in\mathbb{N}\). If it makes sense, we might take, \(T=\{\infty\}\), in which case \(I=\overline{\mathbb{R}}=\mathbb{R}_+\cup\{\infty\}\) or \(\overline{\mathbb{N}}=\mathbb{N}\cup\{\infty\}\). In either case, we call \(T\) the final time.\\
Hence, \(\mathcal{F}_t\subset\mathcal{F}\) is non-decreasing in \(t\in I\) as a sequence of sub-\(\sigma\) fields. \(\mathcal{F}_t\) models the knowledge available at time \(t\); hence, the filtration \(\mathbb{F}\) models the flow of available information through time. More precisely, \(\mathbb{F}\) shall not be any filtration but one generated by market information. The latter is modeled as a \(\mathbb{R}^{r}\)-valued random process \(Y=(Y_t)_{t\in I}\), where \(r\) is a natural number. Hence, we define the raw filtration generated by the information process \(Y\) as \(\mathcal{F}_t^0=\sigma(Y_s:0\leq s\leq t)\). For the canonical completion of the latter, we denote \begin{equation}
    \overline{\mathcal{F}}_t=\sigma(Y_s:0\leq s\leq t)\vee \mathcal{N}_{\mathbb{P}}\;,
\end{equation} where 
\begin{equation}\label{eq:nullsets}
    \mathcal{N}_{\mathbb{P}}=\left\{A\in2^{\Omega}:\exists B\in\mathcal{F}, A\subset B, \mathbb{P(B)}=0\right\}\;.
\end{equation} 
Then \(\mathbb{F}^0\) and  \(\mathbb{F}\) denote the raw and completed filtration respectively generated by the market information \(Y\). We give a useful charaterisation of the completed \(\sigma\)-algebra as stated in \parencite[Thm.~1.9]{Folland1999RealAnalysis}
\begin{lemma}\label{lemma:characterisation of completed sigma-algebra}
    Let \(\overline{\mathcal{F}}=\sigma(\mathcal{F}\vee\mathcal{N}_{\mathbb{P}})\). Then we have that 
    \(\mathcal{C}=\left\{A\cup E:A\in\mathcal{F},E\in\mathcal{N}_{\mathbb{P}}\right\}\) is a \(\sigma\)-algebra and equal to \(\overline{\mathcal{F}}\).
\end{lemma}
\begin{proof}
    The fact that \(\mathcal{C}\) is a \(\sigma\)-algebra can be found in \parencite[Thm.~1.9]{Folland1999RealAnalysis}. The equality follows from \(\overline{\mathcal{F}}=\sigma(\mathcal{C})\) by definition and the fact that \(\mathcal{C}\) is a \(\sigma\)-algebra.
\end{proof}
The price of \(d+1\) tradable assets will be modeled by a \(\mathbb{R}_+^{d+1}\)-valued random process \(S = (S_t)_{t\in\mathbb{R}_+}\), with \(S_t(\omega) = (S_t^0(\omega),\dots,S_t^d(\omega))\in \mathbb{R}_+^d\) for every \(t\in I\). Naturally, the stochastic process \(S\) is assumed to be \(\mathbb{F}\)-adapted, i.e., the random variable \(S_t\) is \(\mathcal{F}_t\)-measurable. In other words, the price of the asset \(S\) at time \(t\) is supposed to be known at time \(t\) (and not in advance). We interpret \(S^0\) as the bank account or any other risk free asset. To fix the units of the asset prices once and for all, we define \(X_t^{i}=S_t^{i}/S_t^0\). Hence, \(S_t^0 = 1\) for all times, and we call \(X\) the \textit{discounted} price process.
\\\\
We may want to trade in these \(d+1\) assets by holding a certain number \(\vartheta_t\) of risky assets \(X_t^i,\;i=1,\dots,d\) and a number \(\eta\) of risk-free assets \(X_t^0\equiv 1\). We call the process \((\vartheta,\eta)\) a \textit{trading position} or a \textit{portfolio}. This trading position can be adjusted dynamically and generally evolves stochastically over time. We don't want to allow insider trading, i.e., knowing the price at time \(f\) of a risky asset before time \(t\), hence giving the possibility to adjust the prices in a way that makes a profit with probability 1. In order to enforce this requirement, we require strategies to \textit{predictable} that measurable with respect to the predictable \(\sigma\)-algebra \(\mathbb{P}\).
The discounted \textit{wealth process} is then defined as 
\begin{equation}
    V_t = \vartheta^{tr}X_t+\eta_t1\;.
\end{equation}
By trading the assets \(X\), the accumulated gains and losses up to time \(t\) amount to \(G_t(\vartheta) = (\vartheta \bullet X)_t\).
\\ In the case \(I=\mathbb{N},\;\overline{\mathbb{N}}\) we have: 
\begin{equation}\label{eq:discrete accumulated gains}
    (\vartheta \bullet X)_t=\sum_{k=1}^{n}\vartheta_k(X_{k}-X_{k-1})
\end{equation}
and for \(I=\mathbb{R},\;\overline{\mathbb{R}}\):
\begin{equation}\label{eq:continuous accumulated gains}
    (\vartheta \bullet X)_t=\int_0^t\vartheta_s\diff X_s\;,
\end{equation} 
where the right hand side denotes the stochastic integral of \(\vartheta\) with respect to \(X\). We deliberately use the generic notation \((\vartheta \bullet X)_t\), referring to both as stochastic integrals and will only use the explicit definition if needed. Nevertheless, the representation \cref{eq:discrete accumulated gains} explains why the adaptiveness of \(\vartheta\) suffices to prohibit insider trading and the theory of stochastic integration (see \parencite{Protter2005StochasticIntegration}).
\\
We denote \(\mathcal{L}(X)\) the set of \(X\)-integrable processes \(\vartheta\), that is the set,depending on \(X\), of \(\vartheta\) for which \(\vartheta\bullet X\) is well-defined.
\\\\
Up to this point, we allow for external cash flow or costs. The cumulative amount of cash injected or taken out is given by the cumulative cost process \(C_t = V_t - (\vartheta \bullet X)_t\). Requiring that there is no external cash flow for \(t\in I\) makes the portfolio \textit{self-financed} and the cumulative cost process constant, equal to the initial value of the  portfolio \(C_t=V_0\quad\forall t\in I\). Hence, for a self-financing strategy, the value process \(V\) takes the form
\begin{equation}\label{eq:self financing strategy}
    V_t(\vartheta) = V_0 + (\vartheta\bullet X)_t=V_0+G_t(\vartheta)\;.
\end{equation}
Notice that \cref{eq:self financing strategy} is fully determined by \(V_0\) and the process \(\vartheta\). Indeed, \(\vartheta\) and \(V_0\) determine \(V_t\) by  \cref{eq:self financing strategy}. Finally, \(\eta\) is then determined by \(V_t = \vartheta^TX_t+\eta_t\). Hence, we denote a self-financing strategy by \((V_0,\vartheta)\). For the clarity of exposition, we introduce the notation. In the spirit of the distinction made for Let \(\Theta\) be a subset of \(\mathbb{F}\)-predictable, real-valued stochastic processes:
\begin{equation}
    \Theta\coloneqq\left\{\vartheta=(\vartheta_{t\in I}):\vartheta \text{ is }\mathcal{P}-\text{measurable}\right\}\cap\mathcal{L}(X)\;
\end{equation}
where \(\mathcal{P}\) is the \(\sigma\)-algebra on \(\Omega\times I\) generated by the \(\mathbb{F}\)-adapted, left-continuous processes.
Then 
\begin{equation}
    \mathcal{A}(\Theta)\coloneqq \mathbb{R}+G_T(\Theta)=\left\{V_0+(\vartheta \bullet X)_T:(V_0,\vartheta)\in\mathbb{R}\times \Theta\right\}\;,
\end{equation}
where \(T\in I\) is the final time in \(I\) and \(G_T(\Theta)=\left\{(\vartheta \bullet X)_T:\vartheta\in \Theta\right\}\).\\
Hence, \(\mathcal{A}(\Theta)\) defines the space of value processes for self-financing strategies \((V_0,\vartheta)\in\mathbb{R}\times \Theta\) We fix \(X=(X_0,\dots,X_d)\) fix once and for all and don't explicitely the dependency.
\section{General idea of hedging}
For the sake of a simple exposition of the underlying concepts, we introduce a finite horizon \(T\in\mathbb{R}_+\); i.e., we consider all the processes up to here (S,X,V,C,\(\vartheta\),\(\eta\)) as processes indexed by \(t\in[0,T]\). We call \(T\) the final time or a fixed time of maturity. Then \(V_T = V_0 + \int_{0}^T\vartheta_s \diff X_s\) is the final time value of the self-financing strategy \((V_0,\vartheta)\). As explained in the \cref{sec:hedging} we want to hedge a liablity \(H\) which we model to be a real-valued random variable whose value is known at a time \(T\geq 0\), i.e. \(H\in\mathcal{L}^0(\Omega,\mathcal{F}_T;\mathbb{R})\). For an concrete example, one might think of \(H\) as a europian call option and \(T\) the strike time of that option.
\\
We call a contingent claim \textit{attainable} if there exists a self-financing strategy \((\vartheta,\eta)\) such that
\begin{equation}\label{eq:atteinable strategy}
    V_T = H\quad \mathbb{P}\text{-a.s.}\;.
\end{equation}
A market is \textit{complete} if all contingent claims are attainable. We will concentrate on hedging in incomplete markets. Hence, for a given contingent claim \(H\), there doesn't necessarily exist a self-financing strategy that replicates \(H\). Hence, we can approach this in three ways. Either we insist on \(V_T = H\quad\mathbb{P}\text{-a.s.}\) and let go of the self-financing assumption for the replicating strategy \((V_0,\vartheta)\). This would entail a non-constant cost process, and the criterion we want to minimize is exactly the cumulative costs of the trading strategy. Alternatively, we could insist on the self-financing condition and let go of the exact replication condition. Lastly, we could let go of both assumptions. We will concentrate on the second option; that is, we will insist on the self-replicating property but let go of the perfect replication property. 
\\\\
In particular, the corresponding hedging problem is to find \textit{the optimal} self-replicating strategy \((V_0,\vartheta)\) that best approximates the contingent claim at time \(T\). The nature of the problem now heavily relies on what we mean by optimal. For this, we will concentrate on two approaches: looking for the trading strategy that makes the terminal position of the seller \textit{acceptable} and, alternatively, the difference between the final value of the strategy and the liability. For the former, we will clarify what it means to be \textit{acceptable} , and for the latter, we specify in which sense we define the difference, namely by which notion of distance.
In the most general case, we define hedging for some \textit{hedging measure} \(d:\mathcal{L}^0(\mathbb{P,\mathcal{F};\mathbb{R}})\to\mathbb{R}\), where \(\mathcal{L}^0(\mathbb{P,\mathcal{F};\mathbb{R}})\) denotes the set of real-valued random variables. Then we define the optimal hedging strategy by the solution to the following minimization problem
\begin{equation}\label{eq:general hedging}
    \inf_{(V_0,\vartheta)\in\mathcal{A}(\Theta)}d(H-V_T(\vartheta))\;,
\end{equation}
for some liability \(H\in\mathcal{L}^0(\mathbb{P},\mathcal{F}_T;\mathbb{R})\) and terminal wealth of a self-financing strategy \(V_T(\vartheta)= V_0+(\vartheta\bullet X)_T\). At this point, no finiteness or attainability claims are set upon or deduced from \cref{eq:general hedging}. If, however, \cref{eq:general hedging} admits a solution \((V_0^*,\vartheta^*)\in\mathcal{A}(\Theta)\), in the sense that the infimum of \cref{eq:general hedging} is attained in \(\mathcal{A}(\Theta)\), we call it an \textit{optimal trading strategy}. The meaning of \textit{optimal} thus depends on the choice of the hedging measure.  
\\\\
In what follows, we will consider two different choices for \(d\): convex risk-measure and \(L^2\)-distance. We refer to \(H-V_T(\theta)\) as the \textit{hedging error} of the strategy \(\vartheta\) for the liability \(H\).
\subsection{convex risk measures}
As our first choice, we consider the hedging measure \(d\) to be a \textit{convex risk measure}. As the name suggests, it is supposed to measure the risk of holding an asset position. Hence, we demand that it follows certain rules
\begin{definition}[Convex Risk Measure]
    We call \(\rho:\mathcal{L}^0(\mathcal{F}_T,\mathbb{P};\mathbb{R})\to\mathbb{R}\)
\end{definition} a \textit{convex risk measure} if it is
\begin{itemize}
    \item Montone decreasing: If \(Z_1\geq Z_2\) then \(\rho(Z_1)\leq\rho(Z_2)\)
    \item Convexity: For \(\alpha\in[0,1]\), we have \(\rho(\alpha Z_1+(1-\alpha)Z_2)\leq \alpha\rho(Z_1)+(1-\alpha)\rho(Z_2)\) for any \(Z_1,Z_2\in\mathcal{L}^0(\mathcal{F}_T,\mathbb{P};\mathbb{R})\) 
    \item Cash-Invariant: \(\rho(Z+c)=\rho(Z)-c \quad \text{for all } c\in\mathbb{R}\)
\end{itemize}
Note that we are not enforcing continuity. Nevertheless, if \(\Omega=\{\omega_1,\dots,\omega_N\}\) is finite, then trivially \(\Omega\cong \mathbb{R}^N\) since any \(X\in\mathcal{L}^0(\Omega,\mathbb{P};\mathbb{R})\) is determined by the tuple \((X(\omega_1),\dots,X(\omega_N))\). Hence, in the finite case we can view \(\rho\) as a convex function over \(\mathbb{R}^N\), thus continuous.
\\\\
The above conditions can be translated in practical assmuptions. The monotonicity reflects the fact that a position with a larger payoff has a smaller risk assosciated to it. Convexity formalizes the idea that diversification lowers risk. Cash invariance then models the simple idea that injecting capital lowers the risk of the position by that same amount.
\\\\
Furthermore, we call a trading \(X\) \(\rho\)-acceptable if \(\rho(X)\leq 0\), and we call \(X\) to be a \(\rho\)-indifferent trading position if \(\rho(X) = \rho(0)\); i.e., holding the trading position \(X\) is equivalent, from the point of view of the convex risk measure \(\rho\), to not investing at all. We call \(\rho\) normalized if \(\rho(0)=0\).
\\\\
Notice that by cash invariance \(\rho(X+\rho(X))=0\). In other words, \(\rho(X)\) is the smallest amount of cash that needs to be added to the trading position \(X\) in order to make it \(\rho\)-acceptable. 
\bigbreak
We now consider the following minimization problem. Let \(\rho\) be a convex risk measure and \(H\in \mathcal{L}^0(\mathcal{F}_T,\mathbb{P};\mathbb{R})\) a liability. Then
\begin{equation}\label{eq:risk measure hedging}
    \pi(H) = \inf_{(V_0,\vartheta)\in\mathcal{A}(\Theta)}\rho\left(H-V_0-(\vartheta\bullet X)_T\right)\;.
\end{equation}
We interpret \(\pi(X)\) as the hedging error under optimal hedging, or equivalently, recalling the discussion of \(\rho(X+\rho(X))=0\), the minimal amount of cash, making the position \(H-(\vartheta\bullet X)_T\) \(\rho\)-acceptable. Notice that, if \(\Theta\) is a convex set, \(\pi\) itself is again a convex risk measure. Indeed, cash-invariance and monotonicity follow directly from the corresponding properties of \(\rho\). For convexity, notice that 
\begin{align*}
    &\inf_{\vartheta\in\Theta}\rho(\alpha H_1+(1-\alpha) H_2-V_0-(\vartheta\bullet X)_T)
    \\&=\inf_{\vartheta_1\in\Theta}\rho(\alpha\left( (H_1-V_0-(\vartheta_1\bullet X)_T\right)+(1-\alpha) \left(H_2-V_0-(\vartheta_2\bullet X)_T\right))
    \\&\leq \alpha\inf_{\vartheta_1\in\Theta}\rho( (H_1-V_0-(\vartheta_1\bullet X)_T)+(1-\alpha)\inf_{\vartheta_2\in\Theta}\rho( (H_2-V_0-(\vartheta_2\bullet X)_T)
\end{align*}
Taking \(\inf_{V_0\in\mathbb{R}}\) of the above inequality yields \(\pi(\alpha H_1+(1-\alpha) H_2)\leq \alpha\pi(H_1)+(1-\alpha)\pi(H_2)\), where we used the convexity of \(\Theta\) in the second step and the convexity of \(\rho\) in the third step.
\\
Notice an important subtlety about \cref{eq:risk measure hedging}. Due to the cash-invariance property of the convex risk measure \(\rho\) we have 
\begin{align*}
    \pi(H) = &\inf_{(V_0,\vartheta)\in\mathcal{A}(\Theta)}\rho\left(H-V_0-(\vartheta\bullet X)_T\right)\\
    &=\inf_{\vartheta\in \Theta}\inf_{V_0\in\mathbb{R}}\rho\left(H-V_0-(\vartheta\bullet X)_T\right)\\
    &=\inf_{\vartheta\in \Theta}\inf_{V_0\in\mathbb{R}}\rho\left(H-(\vartheta\bullet X)_T\right)+V_0=-\infty\;,
\end{align*}
if \(\rho\left(H-(\vartheta\bullet X)_T\right)\) is finite,
 we can take \(V_0\) arbitrarily negative.\footnote{ If \(\rho\left(H-(\vartheta\bullet X)_T\right)=+\infty\) the problem is again ill defined.} In this sense, the problem is ill defined. Hence, either we restrict our admissible strategies to those that are lower bounded by some real value \(-b\leq 0\), or we can put additional assumptions on the convex risk measure \(\rho\) (see \parencite{Xu2006RiskMeasurePricingHedging} for a sufficient treatment of this matter). 
Alternatively, we can take a different point of view and consider a fixed and finite \(V_0^*\in\mathbb{R}\) and analyze the minimization problem 
\begin{equation}\label{eq:risk measure heding with fixed initial capital}
    \inf_{\vartheta\in\Theta}\rho\left(H-V_0^*-(\vartheta\bullet X)_T\right)\;.
\end{equation}
This situation can be motivated and related to the realistic scenario where \(V_0^*\) is determined as a result of trading derivatives in the market at competitive prices. The infimum of \cref{eq:risk measure heding with fixed initial capital} might still not be finite.  In the case where \(\Omega\) is finite, it is enough to require \(\pi(0)>-\infty\), as shown in \parencite[Cor.~3.5.]{Buehler2018DeepHedging}, which we recall at this point. 
\begin{lemma}\label{lemma:Finiteness of inf for finite sample space}
    Assume \(\Omega = (\omega_1,\dots,\omega_N)\) for \(N\in\mathbb{N}\) and \(\pi(0)>-\infty\). Then \(\pi(Z)>-\infty\) for all \(Z\in\mathcal{L}^0(\mathcal{F}_T,\mathbb{P};\mathbb{R})\)\;.
\end{lemma}
\begin{proof}
    Indeed, as \(\Omega\) is finite, the se \(\sup_{\omega\in\Omega} X(\omega)\equiv\sup X<\infty\). Hence, using the monotonicity of the convex risk measures \(\pi\) and \(\pi(X)\geq\pi(\sup X)=\pi(0)-\sup X>-\infty\).
\end{proof}
In the general case one has to restrict to bounded cumulative gains 
\begin{equation}
G_T^b(\Theta)=\left\{(\vartheta \bullet X)_T:\vartheta\in \Theta, (\vartheta \bullet X)_T\geq -b\right\}\;,
\end{equation}for some \(b\geq 0\). Imposing further assumptions on the convex risk measure \(\rho\) one can show the existence of the solution 
\begin{equation}
    \pi_b(H) = \inf_{(V_0,\vartheta)\in\mathcal{A}_b(\Theta))}\rho(H-V_0-(\vartheta\bullet X)_T))=\rho(H-V_0^*-(\vartheta^*\bullet X)_T))\;,
\end{equation}
where \(\mathcal{A}_b(\Theta)=\mathbb{R}+G_T^b(\Theta)\). For the proof and the exact assumptions, see \parencite[Lemma~2.5]{Xu2006RiskMeasurePricingHedging} and  \parencite[Thm~2.6]{Xu2006RiskMeasurePricingHedging}.


\subsection{Mean-square hedging}\label{sec:high level discussion of mean-square hedging}
For the second choice of hedging measure \(d\) we take the \(L^2\)-norm that is \(d(X)=\mathbb{E}[|X|^2]\). Notice that this object to make sense (we mean finiteness in this context) we need to restrict the hedging measure to \(d:L^2(\mathbb{\mathbb{P},\mathcal{F};\mathbb{R}})\to\mathbb{R}\) and assume \(H\in L^2(\mathbb{P})\). The corresponding hedging problem then takes the form 
\begin{equation}\label{eq:mean squared hedging}
\inf_{(V_0,\vartheta)\in\mathcal{A}(\Theta)}\mathbb{E}[\left(H-V_0-(\vartheta\bullet X)_T)\right)^2]\;,
\end{equation}
for some fixed \(V_0\in\mathbb{R}\).
\\
First observe that \cref{eq:mean squared hedging} doesn't suffer from the same ill posedness issue as \cref{eq:risk measure hedging} as \(V_0\) enters quadratically.
Furthermore notice that if \(G_T(\Theta)\) is a closed subspace of \(L^2\) then \cref{eq:mean squared hedging} reduces to a \(L^2\) projection problem of the \(L^2\)-random variable \(H\) onto \(G_T(\Theta)\). 
\\
Indeed, then \cref{eq:mean squared hedging} asks for the element in the closed subspace \(\mathcal{A}(\Theta)\) that minimizes the \(L^2\)-distance. Recalling that \((L^2(\mathbb{R)},\lVert\cdot\rVert_{L^2})\) is a Hilbert space, we get existence and uniqueness of the minizer of \cref{eq:mean squared hedging} from \parencite[Thm.~12.3]{Rudin1991FunctionalAnalysis}. 
Furthermore, if \(\Theta\subset \Theta^2\), where \(\Theta^2\coloneqq\left\{\vartheta\in\mathcal{L}(X):(\vartheta\bullet X)_T\in L^2(\mathcal{F},\mathbb{P})\right\}\), then notice that due to continuity of the map, \footnote{It follows as the composition of continuous functions: \(L^2\)-norm and addition in \(\mathbb{R}\). This uses the general fact that for \(f:X\to\mathbb{R}\) continuous for some topological space \(X\), we have for any subset (\(Y\subset X\)) that \(\inf d(Y)=\inf d(\overline{Y})\). For a short proof, see \cref{Appendix:Infima over closure}}.
\begin{equation}
\begin{cases}
    (\mathcal{A}(\Theta),\lVert\cdot\rVert_{L^2})\to(\mathbb{R},|\cdot|_{\mathbb{R}})\\
    V_T\mapsto\lVert H-V_0-V_T\rVert_{L^2}\;,
\end{cases}
\end{equation}
the minimization problem \cref{eq:mean squared hedging} is equivalent to 
\begin{equation}\label{eq:mean squared hedging over closure}
\inf_{V_T\in \overline{G_T(\Theta)}}\mathbb{E}[\left(H-V_T\right)^2]\;,
\end{equation}
where \(\overline{G_T(\Theta)}\) denotes the \(L^2\)-closure of \(G_T(\Theta)\). In particular, due to the closedness of \(\overline{G_T(\Theta)}\) in the Hilbert space \(L^2(\mathbb{P})\), and hence by application of the Hilbert projection theorem, \cref{eq:mean squared hedging over closure} admits a unique solution \(V_T^*\in\overline{G_T(\Theta)}\). 
Hence, we don't get an explicit expression for the optimal value process at time \(T\). Thus, it is useful to study conditions of \(\Theta\) for which \(G_T(\Theta)\) is closed in \(L^2(\mathbb{P})\). Indeed, if this is the case, by the same Hilbert projection argument, we get the existence of a unique \(V_T^*\in G_T(\Theta)\), which, as an element of \(G_T(\Theta)\), takes the form \(V_T^*=\int_0^T\vartheta^*\diff X_s\).


\chapter{Deep Hedging in Discrete Time}\label{ch:Deep Hedging in Discrete Time}
After establishing a good understanding of the hedging problem in general terms we further discuss deep hedging. We will approach it generally at first and then go into more detail into the set up and treatement of \parencite{Buehler2018DeepHedging}. Recall that the hedging problem, as we treat it in this article, is given by \cref{eq:general hedging} as the approximation of a liability \(H\) by the value process of a self-financing strategy \(\vartheta\in\Theta\) for some set of admissible trading strategies \(\Theta\) and some notion of closeness. 

\section{The deep hedging framework of Bühler et al.}
The basic idea in deep hedging is to approximate the strategies \(\vartheta\in\Theta\) by a \textit{good} set of approximants. By good, we mean that the approximation should be arbitrarily exact and implementation friendly. The limiting exactness of our approximants is central to the theoretical development of \parencite{Buehler2018DeepHedging} as discussed in the following, and the practical importance of neural networks will motivate a specific choice for our approximation system. The specific setup we discuss in this chapter is based on a finite sample space \(\Omega\). 
\\\\
What makes \parencite{Buehler2018DeepHedging} particularly powerful is its practical generality, particularly allowing for trading costs. 
The main problem that is directly addressed theoretically in \parencite{Buehler2018DeepHedging} is convex risk-measure hedging as introduced in \cref{eq:risk measure hedging}. In light of our discussion of the continuous time setup, we will nevertheless consider a frictionless market, that is, without trading costs. In this case, we will show how the theoretical results from \parencite{Buehler2018DeepHedging} apply to mean-variance hedging in the finite setup. This in particular adds to \parencite[Remark.~3.4]{Buehler2018DeepHedging} and conceptually links their work to our discussion of deep mean-variance hedging.

What makes \parencite{Buehler2018DeepHedging} particularly powerful is its practical generality, especially in allowing for trading costs: \(C_t(\vartheta)\) and modeling of constrained strategies denoted by \(\mathcal{H}\subset\mathcal{H}û\), where the latter denotes the unconstrained strategies.
\\
As the purpose of this exposition is to understand deep hedging and, in particular, its relation to deep mean-variance hedging, we won't assume unnecessary generality and will specialize in this case. 
\begin{align}\label{eq:simplifying assumptions}
    &C_t(\vartheta)=0\quad\forall \vartheta\quad\forall t\geq 0\;,\\
    &\mathcal{H}=\mathcal{H}^u=\Theta\;,
\end{align}
where \(\Theta\) denotes the set of \(\mathbb{F}\)-predictable processes.
\subsection{Market model and trading strategies}
The current setup still follows, for the most part, the general setup described in \cref{sec:hedging}. Some simplifying assumptions are made that make the treatment technically easier compared to the general continuous time treatment, as discussed in \cref{sec:deep mean-square hedging in continuous time}. In particular, the main simplifications follow from the assumptions of a discrete financial market with a finite time horizon \(T\), finitely many trading times \(0=t_0<t_1<\dots<t_n=T\), and a finite sample space \(\Omega=\{\omega_1,\dots,\omega_N\}\). Hence, the market-filtration \(\mathbb{F}=(\mathcal{F})_{k=0}^n\), generated by the \(\mathbb{R}^r\)-valued market process \(Y=(Y_k)_{k=0}^n\), is also finitely generated in the sense that \(\mathcal{F}_k=\sigma(Y_0,\dots,Y_k)\). If we consider \(\mathcal{F}_k\) the available information up to time \(t_k\), we have, by the Doob-Dynkin lemma \parencite[Lemma.~1.14]{Kallenberg2002FoundationsModernProbability}, that any random quantity whose value can be deduced from the information available at time \(t_k\), that is any \(\mathcal{F}_k\)-measurable random variable, is a function of \(Y_0,\dots,Y_k\).  
\bigbreak
Recall that the optimization problem \cref{eq:risk measure hedging} is ill defined in the sense that \(\pi(H)=-\infty\) for all \(H\in\mathcal{L}^0(\mathbb{P},\mathcal{F}_T;\mathbb{R})\) if we don't put additional assumptions on \(\mathcal{A}(\Theta)\) and \(\rho\). An alternative, related problem, which doesn't suffer this degenerate behavior, is to consider 
\begin{equation}\label{eq:risk measure hedging from buehler}
    \pi(L)=\inf_{\vartheta\in\Theta}\rho(L-(\vartheta\bullet X)_T)\;,
\end{equation}
for \(L\in\mathcal{L}^0(\mathbb{P},\mathcal{F},\mathbb{R})\) and where one optimizes over \(\vartheta\in\Theta\). \cref{lemma:Finiteness of inf for finite sample space} still applies to \cref{eq:risk measure hedging from buehler} and ensures finiteness if \(\pi(0)>-\infty\). \cref{eq:risk measure hedging from buehler} is then the hedging problem that is explicitly addressed in \parencite{Buehler2018DeepHedging}.
\\
Commenting on the interpretation of \cref{eq:risk measure hedging from buehler} we observe that it corresponds to the optimal hedging error for \(V_0=0\) initial capital, which feels natural. Nevertheless,is due to the observations made earlier regarding the interpretation of the convex risk measure, \(\pi(H)\) of  \cref{eq:risk measure hedging from buehler} is the minimal additional cash that needs to be added to the traders position in order to make it \(\rho\)-acceptable. 
\section{Approximation of trading strategies}\label{sec:Approximation of trading strategies}
Adapting the point of view presented in \parencite{Buehler2018DeepHedging}, we are not so much interested in sufficient conditions for the infimum in \cref{eq:risk measure hedging from buehler} to be attained, but rather in the consistent approximation of \(\pi(L)\) for any \(L\in\mathcal{L}^0(\mathbb{P},\mathcal{F}_T;\mathbb{R})\) through approximating the trading strategies in \(\Theta\). The underlying idea is to leverage the \(\mathbb{F}\)-predictability of \(\vartheta\in\Theta\) by approximating its respective Doob-Dynkin functional representation with a suitable set of functions. Mathematically speaking, we want to find a suitable dense subset of the function space that encompasses the Doob-Dynkin representations of the \(\mathcal{F}_{k-1}\)-predictable \(\vartheta_k\). That is:
\begin{equation}
   \vartheta_k=f_k(Y_0,\dots,Y_{k-1})\;,
\end{equation}
where \(f_k:\mathbb{R}^{rk}\to\mathbb{R}^d\) is Borel-measurable.
\\
Notice that in our current setting, finite \(\Omega\) we have that 
\begin{equation}
    \mathbb{E}_{\mathbb{P}}\left[f_k(Y_0,\dots,Y_{k-1})\right]=\sum_{\omega\in\Omega}f_k(Y_0(\omega),\dots,Y_{k-1}(\omega))<\infty\;,
\end{equation}
is finite as a finite sum over finite terms.
In particular, due to
\begin{equation}
    \mathbb{E}_{\mathbb{P}}\left[f_k(Y_0,\dots,Y_{k-1})\right]=\mathbb{E}_{\mathcal{L}(Y_0,\dots,Y_{k-1})}\left[f_k\right]\;,
\end{equation}
where \(\mathcal{L}(Y_0,\dots,Y_{k-1})\) denotes the law of \((Y_0,\dots,Y_{k-1})\)\footnote{That is, the push-forward measure \(\mathbb{P}_*(Y_0,\dots,Y_{k-1})\coloneqq\mathbb{P}\circ(Y_0,\dots,Y_{k-1})^{-1}\).}, we have 
\begin{equation}\label{eq:f is in L1 for finite Omega}
    f_k\in L^1(\mathbb{R}^{rk},\mathcal{L}(Y_0,\dots,Y_{k-1}))\;.
\end{equation}
More generally, we have
\begin{equation}
    f_k\in L^p(\mathbb{R}^{rk},\mathcal{L}(Y_0,\dots,Y_{k-1}))\quad\forall p>0\;.
\end{equation}
In the probabilistic case, measure spaces are finite (probability spaces); hence \(L^p\subset L^q\) for \(p\geq q\). Hence, we need an approximation system \(\Psi\) that approximates functions from \(\subset\  L^1(\mathbb{R}^{r(k+1)},\mathcal{L}(Y_0,\dots,Y_{k-1}))\) arbitrarily well. Technically, we want the following density statement.
\begin{equation}
    \overline{\Psi}^{L^1}=L^1(\mathbb{R}^{r(k+1)},\mathcal{L}(Y_0,\dots,Y_{k-1}))\;.
\end{equation}
This is enough for the \(\Omega\)-finite case. Looking ahead, we want to treat the general \(\Omega\) case. We need to require \(\Psi\) to be in \(L^p\quad p>0\). We will do this by requiring elements in \(\Psi\) to be essentially bounded, that is almost surely bounded. Thus, what we want is \(\Psi\) such that:
\begin{align*}
  &\Psi\subset\mathcal{L}^{\infty}(\Omega,\mathbb{P})\\
  &\overline{\Psi}^{L^1}=L^1(\Omega,\mathbb{P})\;.
\end{align*}
We will only motivate and introduce a concrete incarnation of such a \(\Psi\) at the end of the current chapter and discuss it in \cref{ch:Neural Networks and Universal Approximation}. For now and the rest of this section we shall remain general.
\bigbreak
We now come back to the original task of approximating \(\pi(L)\) of \cref{eq:risk measure hedging from buehler} by approximating the strategies in \(\Theta\) using our approximation system \(\Psi\). Following \parencite{Buehler2018DeepHedging}  we introduce a stratification on \(\Psi\) as follows. Denote by \((\Psi_M)_{M\in\mathbb{N}}\) a sequence of subsets of \(\Psi\) such that 
\begin{itemize}
    \item \(\Psi_M\subset\Psi_{M+1}\quad\forall M\in\mathbb{N}\)\;,
    \item \(\bigcup_{M\in\mathbb{N}}\Psi_M=\Psi\)\;,
    \item for all \(M\in\mathbb{N}\) there exists a parameter space \(\Sigma_M\subset\mathbb{R}^t\) for some \(t\in\mathbb{N}\) depending on \(M\)\;.
\end{itemize}
The last property leads us to interpret \(\Psi_M\) as a subfamily of \(\Psi\) parameterized by \(\Sigma_M\). For instance, we might think of \(\Psi\) as polynomials of arbitrary degree and \(\Psi_M\) as polynomials of degree less than or equal to \(M\in\mathbb{N}\)\footnote{Note that this example only serves the purpose of providing concrete intuition about the introduction of this stratified family \(\bigcup_{M\in\mathbb{N}}\Psi_M=\Psi\). In particular, polynomials are not generally dense in \(L^p(\Omega,\mathbb{P})\) (They are if \(\Omega\subset \mathbb{R}^s,s\in\mathbb{N}\) is compact)}. \\
As we want to approximate strategies with \(\Psi\) we define the following subset of \(\Theta\):
\begin{equation}\label{eq:restricted strategies}
    \Theta_M\coloneqq\{\vartheta=(\vartheta_k)_{k=0}^{n-1}\in\Theta:\vartheta_k=F_k(Y_0,\dots,Y_{k-1},\vartheta_{k-1}),\quad F_k\in\Psi_M\}\;.
\end{equation}
Notice the analogy with the Doob-Dynkin representation of the strategy \(\vartheta\). Do to the same Doob-Dynking representation theorem \(\vartheta_{k-1}\) is itself already some measurable function of \((Y_0,\dots,Y_{k-1})\), thus we can omitt \(\vartheta_{k-1}\) dependency in the definition \cref{eq:restricted strategies}. Thus, \(\Theta_M\) is the set of strategies that have a Doob-Dynkin representation in \(\Psi_M\). 
We further denote
\begin{equation}
    \Theta_{\infty}\coloneqq \bigcup_{M\in\mathbb{N}}\Theta_M
\end{equation}

We can now state the restricted optimization problem. 
\begin{equation}\label{eq:restricted risk measure hedging from buehler}
    \pi^M(L)=\inf_{\vartheta\in\Theta_M}\rho(L-(\vartheta\bullet X)_T)\;.
\end{equation}
The natural theorem we would like to show is that the restricted hedging errors from \cref{eq:restricted risk measure hedging from buehler} converge to the unrestricted hedging error \cref{eq:risk measure hedging from buehler} for \(M\to\infty\). Indeed, in our current setup, this holds true and is the core theoretical result of \parencite{Buehler2018DeepHedging}. 
\begin{theorem}[Consistent Approximation of Hedging Errors]\label{thm:approximation by restricted hedging errors}
 Let \(\Theta_M\subset \Theta\) be as in \cref{eq:restricted strategies} and \(\pi^M\) as in \cref{eq:restricted risk measure hedging from buehler}. Then, for any \(L\in\mathcal{L}^0(\mathbb{P},\mathcal{F}_T;\mathbb{R})\),
 \begin{equation}
     \lim_{M\to\infty}\pi^M(L)=\pi(L)\;.
 \end{equation}
\end{theorem}
Before giving the proof, we conceptually argue how we expect the proof to proceed. Recall that the strategy \(\vartheta=(\vartheta_k)_{k=0}^{n-1}\) is a \(\mathbb{F}\)-adapted process, where \(\mathcal{F}_k=\sigma(Y_0,\dots,Y_{k})\). By Doob-Dynkin, for every \(k=0,\dots n-1\), there exists a \(f_k:\mathbb{R}^{r(k+1)}\to\mathbb{R}^d\) measurable function such that. 
\begin{equation}
    \vartheta_k=f_k(Y_0,\dots,Y_{k-1})\;.
\end{equation}
As we have argued previously in \cref{eq:f is in L1 for finite Omega} we have \(f\in L^1(\mathbb{R}^{r(k+1)},\mathcal{L}(Y_0,\dots,Y_{k-1}))\). We can thus approximate it by elements in \(\Psi\). Because \(\bigcup_{M\in\mathbb{N}}\Psi_M=\Psi\), for a given approximation error bound \(\epsilon>0\) we can do so by elements of \(\Psi_M\) for \(M\in\mathbb{N}\) large enough. The hope is then that we can also approximate the hedging error arbitrarily well in the sense that we can get \(\epsilon\)-close with large enough \(M\). To show that this works requires careful argumentation and will be done in the following proof
\begin{proof}
 First, notice that due to the monotonicity \(\Psi_M\subset\Psi_{M+1}\subset\Psi\) we have \(\Theta_M\subset\Theta_{M+1}\subset\Theta\) hence \(\pi^M(L)\geq\pi^{M+1}(L)\geq\pi(L)\). Hence, the convergence of \(\pi(L)\) from above is equivalent to convergence from above. Thus, it is enough to show that for any \(\epsilon>0\), there exists \(M\in\mathbb{N}\) such that 
 \begin{equation}
     \pi^M(L)\leq \pi(L)+\epsilon\;.
 \end{equation}
 Furthermore, by the definition of the infimum over a real set, for every \(\epsilon/2>0\), there exists \(\vartheta\in\Theta\) such that
 \begin{equation}
     \rho(L-(\vartheta\bullet X)_T)\leq\pi(L)+\frac{\epsilon}{2}\;.
 \end{equation}
 Consider this \(\vartheta=(\vartheta_k)_{k=0}^{n-1}\). By definition, \(\vartheta_k\) is \(\mathcal{F}_k\)-measurable. So, by the Doob-Dynkin lemma, there exists \(f_k:\mathbb{R}^{rk}\to\mathbb{R}^d\) measurable such that 
 \begin{equation}
     \vartheta_k=f_k(Y_0,\dots,Y_{k-1})\;,
 \end{equation}
 for all \(k=0,\dots,n-1\). As argued in \cref{eq:f is in L1 for finite Omega} \(f_k\in L^1(\mathbb{R}^{rk},\mathcal{L}(Y_0,\dots,Y_{k-1}))\) because \(\Omega\) is finite. Let's denote \(\mu\coloneqq\mathcal{L}(Y_0,\dots,Y_{k-1})\). By the density of \(\Psi\) in \(L^1\), there exists a sequence of functions \(F_k^n\) converging to \(f_k\) in \(L^1(\mathbb{R}^{rk},\mu;\mathbb{R}^d)\) or, equivalently, \(\vartheta_k^n\coloneqq F_k^n(Y_0,\dots,Y_{k-1})\) converges to \(f_k(Y_0,\dots,Y_{k-1})\) in \(L^1(\Omega,\mathbb{P};\mathbb{R}^d)\) as \(n\to\infty\)\footnote{in the sense that \(\int_{\mathbb{R}^{r(k+1)}}\lVert f_k-F_k^n\rVert_{d}\diff \mu\to 0\) as \(n\to\infty\), where \(\lVert\cdot\lVert_d\) denotes the Euclidean norm on \(\mathbb{R}^d\)}. Passing to a subsequence and by a diagonal argument, we get \(\mathbb{P}\)-almost sure convergence for all \(k\) simultaneously. Morally speaking, we want to pass the limit \(n\to\infty\) inside \(\rho\). Even if \(\rho\) has sufficient regularity, \(\mathbb{P}\)-a.s. convergence is not enough to get the claim. Lukely we are in the finite setting, hence \(\mathbb{P}(\{\omega\})>0\) for all \(\omega\in\Omega\), hence \(\mathbb{P}\)-almost sure convergence is equivalent to convergence everywhere. Thus
 \begin{equation}
     \lim_{n\to\infty}\vartheta_k^n(\omega)=\vartheta_k(\omega)\quad\forall\omega\in\Omega\;,k=0,\dots,n-1\;.
 \end{equation}
 \\
 In particular, 
 \begin{align}
(\vartheta^n\bullet X)_T(\omega)=\sum_{k=1}^{n}\vartheta^n_k(\omega)\cdot(X_{k}(\omega)-X_{k-1}(\omega))\overset{n\to\infty}{\rightarrow}(\vartheta\bullet X)_T(\omega)\;.
\end{align}
 Furthermore, notice that \(\mathcal{L}^0(\Omega,\mathcal{F}_T;\mathbb{R})\simeq\mathbb{R}^{N}\), hence we can identify \(\rho\) with a convex map \(\tilde{\rho}:\mathbb{R}^{N}\to\mathbb{R}\). In the following we will not differentiate \(\rho\) adn \(\tilde{\rho}\). Notice that \(\rho\) is convex function on a finite dimensional vector space and as such continuous (see \parencite[Thm.~1.3.14.]{Cheridito2013ConvexAnalysis}). Thus we have
 \begin{equation}
     \lim_{n\to\infty}\rho(L-(\vartheta^n\bullet X)_T)=\rho(L-(\vartheta\bullet X)_T)\;.
\end{equation}
In particular, there exists \(n\in\mathbb{N}\) large enough such that 
\begin{equation}
    \rho(L-(\vartheta^n\bullet X)_T)\leq\rho(L-(\vartheta\bullet X)_T)+\frac{\epsilon}{2}\leq\pi(L)+\epsilon\;.
\end{equation}
Since for every \(\vartheta^n\) there exists \(M\in\mathbb{N}\) large enough for which \(\vartheta^n\in\Theta_M\). Thus, we have 
\begin{equation}
    \pi^M(L)\leq\pi(L)+\epsilon\;,
\end{equation}
which ends the proof.
\end{proof}
\section{From convex risk measures to quadratic loss}
As previously pointed out, the minimization problem in \cref{eq:risk measure hedging from buehler}, used in particular in the convergence result \cref{thm:approximation by restricted hedging errors}, is not quite the same minimization problem \cref{eq:risk measure hedging}. Even worse, we have to study mean-squared hedging separately from the convex risk measure hedging, as the \(L^2\)-norm is not a convex risk measure (it lacks cash-invariance and the monotonicity is reversed). Nevertheless, the methods used in the proof of \cref{thm:approximation by restricted hedging errors} apply analogously to the mean variance hedging setup, as remarked in \parencite[Rem~.3.4.]{Buehler2018DeepHedging}. In the setting of \parencite{Buehler2018DeepHedging} and more precisely \parencite[Rem~.3.4.]{Buehler2018DeepHedging}, the initial capital \(V_0\) is fixed exogenously, for instance by the market. In the following, we show an analogous statement to \cref{thm:approximation by restricted hedging errors} for the mean variance setup, still with finite \(\Omega\). We argue that this is a stronger statement than the mean-variance minimization problem \cref{prop:finite deep mean variance hedging}. Lastly, we show how to generalize to the joint minimization over self-financing strategies \((V_0,\vartheta)\).
\subsection{Approximate optimal hedging in the context of deep hedging}
As before, compared to the general setup in \parencite{Buehler2018DeepHedging} we specialize to the case where do not consider the differentiation between constrained and unconstrained strategies. We also neglect trading costs, that is, \(C_t=0\quad \forall t\geq 0\). The minimization problem we study now takes the form 
\begin{equation}\label{eq:hedging error}
    \Delta(H) = \inf_{\vartheta\in\Theta}\mathbb{E}[\left(H-V_0-\vartheta\bullet X\right)_T)^2]
\end{equation}
and in analogy to \parencite[Thm.~4.9.]{Buehler2018DeepHedging} we introduce
\begin{equation}\label{eq:restricted mean variance hedging error}
    \Delta^M(H) = \inf_{\vartheta\in\Theta_M}\mathbb{E}[\left(H-V_0-\vartheta\bullet X\right)_T)^2]
\end{equation}
for \(\Theta_M\) as in \cref{eq:restricted strategies} and \(V_0\in\mathbb{R}\) fixed as in the setup of \cref{eq:risk measure heding with fixed initial capital}. We want to show the following discrete time deep mean hedging theorem.
\begin{proposition}[Finite Deep Mean variance Hedging]\label{prop:finite deep mean variance hedging}
    \begin{equation}
        \inf_{\vartheta\in\Theta}\mathbb{E}[\left(H-V_0-(\vartheta\bullet X)_T\right)^2]=\inf_{\vartheta\in\Theta_{\infty}}\mathbb{E}[\left(H-V_0-(\vartheta\bullet X)_T\right)^2].
    \end{equation}
\end{proposition}
More precisely, we will show the stronger statement
\begin{theorem}\label{thm:strong discrete deep mean variance hedging}
Let \(\Theta_M\) be as in \cref{eq:restricted strategies}, \(\Delta^M(H)\) as in \cref{eq:restricted mean variance hedging error}, and \(\Delta(H)\) as in \cref{eq:hedging error} respectively. Then, for any \(H\in\mathcal{L}^0(\mathcal{F}_T,\mathbb{P})\),
    \begin{equation}
        \lim_{M\to\infty}\Delta^M(H)\to\Delta(H)\;.
    \end{equation}
\end{theorem}
Indeed \cref{thm:strong discrete deep mean variance hedging} implies \cref{prop:finite deep mean variance hedging}. To see this more generally, let \(J:\mathcal{H}\to\mathbb{R}\) and assume \(\lim_{M\to\infty}\inf J(\mathcal{H}^M)=\inf J(\mathcal{H}^u)\). Then there exists \((x_n)_{n\in\mathbb{N}}\subset\bigcup_{n\geq 1}\mathcal{H}^M\) such that \(\lim_{n\to\infty}J(x_n)=\inf J(\mathcal{H}^u)\). This then implies \(\inf J(\mathcal{\bigcup_{M\geq 1}\mathcal{H}}^M)=\inf J(\mathcal{H}^u)\). Indeed, because \(\bigcup_{M\geq 1}\mathcal{H}^M\subset \mathcal{H}^u\) hence \(\inf J(\mathcal{H}^u)\) is a lower bound for \(\inf J(\mathcal{\bigcup_{M\geq 1}\mathcal{H}}^M)\). Those it remainse to show that \(\forall \delta>0\) there exists \(x_n\in \bigcup_{M\geq 1}\mathcal{H}^M)\) such that \(\inf J(\mathcal{\bigcup_{M\geq 1}\mathcal{H}}^M)<\inf J(\mathcal{H}^u)+\delta\). This is equivalent to finding a sequence \((x_n)_{n\in\mathbb{N}}\subset \bigcup_{M\geq 1}\mathcal{H}^M\) such that \(\lim_{n\to\infty}J(x_n)=\inf J(\mathcal{H)}\).\footnote{This equivalence is more or less by definition but rigourusly it goes like this. For once if the the property holds for every \(\delta>0\) we can consider the nullsequence \(\delta_n = 1/n\) and get for every \(n\in\mathbb{N}\) an \(x_n\) such that \(J(x_n)\) is \(\delta_n\) close to \(\inf J(\mathcal{H})\). Conversely, if there exists such a sequence then \(\forall \delta>0\quad \exists N\in\mathbb{N}\) for which \(\left|J(x_n)-\inf J(\mathcal{H})\right|<\delta\). But \(J(x_n)\geq \inf J(\mathcal{H}^M)\geq \inf J(\mathcal{H})\) for some \(M\in\mathbb{N}\), so \(J(x_n)<\inf J(\mathcal{H})+\delta\).} For \(J(\vartheta)\coloneqq\mathbb{E}[\left(H-p_0-(\vartheta\bullet X)_T\right)^2]\) we get the claim.
\\\\
For the proof of \cref{thm:strong discrete deep mean variance hedging} we proceed very similar to the strategy used in the proof of \parencite[Thm.~4.9.]{Buehler2018DeepHedging} with the simplification as previously meantioned:
\begin{proof}
Notice first that by the definition, we have \(\Psi=\bigcup_{M\geq 1}\Psi_M\subset\Theta\) and \(\Psi_M\subset\Psi_{M+1}\subset\bigcup_{M\geq 1}\Psi_M\subset \Theta\) for all \(M\in\mathbb{N}\). Hence \(\Delta^M(H)\leq\Delta^{M+1}(H)\leq\Delta(H)\). So in particular it suffised to show that \(\forall \epsilon>0\quad\exists M\in\mathbb{N}\) such that \(\Delta^M(H)\leq\Delta(H)+\epsilon\).
The proof now is exactly the same as in \parencite[Thm.~4.9.]{Buehler2018DeepHedging} just using the fact that the expectations
\begin{equation}\label{eq:Continuity of L^2 norm}
    \mathbb{E}[\left(H-V_0-(\vartheta\bullet X)_T\right)^2]=\sum_{\omega\in\Omega}\left(H(\omega)-V_0-(\vartheta\bullet X)_T(\omega)\right)^2\;,
\end{equation}
are finite sums because the sample space \(\Omega = \left\{\omega_1,\dots,\omega_N\right\}\) is finite and the square function is continuous everywhere. As we are dealing with expectations, we don't even need convergence in every \(\omega\in\Omega\) but almost everywhere convergence of the approximation functions \(\vartheta_k^n\coloneqq F_{k,n}(Y_0,\dots,Y_{k-1})\) to \(\vartheta_k\) (where the latter is \(\vartheta\in\Theta\) such that \(\mathbb{E}[\left(H-p_0-(\vartheta\bullet X)_T\right)^2]\leq\Delta(H)+\epsilon\)).
\end{proof}
This shows that, as alluded to in \parencite[Remark~3.4]{Buehler2018DeepHedging}, the approximate optimal hedging problem, and in particular the mean hedging problem, is subsumed in the framework of \parencite{Buehler2018DeepHedging}. Indeed the proof above for \(l(\cdot)=(\cdot)^2\) goes through for any continuous function \(l:\mathbb{R}\to\mathbb{R}_+\).
\\
We may want to extend the setting to the joint optimization problem of \cref{eq:mean squared hedging}. For this consider the following modified set up.
\begin{align*}
 &(\tilde{X}_k^i)_{k=0}^n \text{ for } i=0,\dots,d \text{, where } X_k^0 = k/n, \tilde{X}_k^i = X_k^i\quad i = 1,\dots,d \text{ and } k=0,\dots n.\\
 &
    (\tilde{\vartheta}_k^i)_{k=0}^n \text{ for } i=0,\dots,d \text{, where } \tilde{\vartheta}_k^0 = V_0, \delta_k^i = \tilde{\vartheta}_k^i\quad i = 1,\dots,d \text{ and } k=0,\dots n.
\end{align*}
Then observe that \(\tilde{X}\) and \(\tilde{\vartheta}\) are still \(\mathbb{F}\)-adapted and 
\begin{equation}
    (\tilde{\vartheta}\bullet\tilde{X})_T=V_0+(\vartheta\bullet X)_T\;.
\end{equation}
Due to adaptability, we still have the Doob-Dynkin representation of the \(\mathcal{F}_k\)-measurable function \(\tilde{\vartheta}_k=f_k(Y_0,\dots,Y_{k-1})\) for some measurable function \(f_k:\mathbb{R}^{r(k+1)}\to\mathbb{R}^{d+1}\). Because \(\Omega\) is finite, \(\tilde{\vartheta}\) is bounded, hence \(f_k\in L^1(\mathbb{R}^{r(k+1),\mu})\), where \(\mu=\mathcal{L}(Y_0,\dots,Y_{k-1})\). Hence we can approximate \(f_k\) by elements of our dense approximants \(\Psi\). Whence we get the same results as in \cref{prop:finite deep mean variance hedging} and \cref{thm:strong discrete deep mean variance hedging} for the corresponding joint minimization over \((V_0,\vartheta)\in\mathcal{A}(\Theta)
=
\left\{
\tilde{\vartheta}
=
(\tilde{\vartheta}_k)_{k=0}^{n-1}
:\;
\tilde{\vartheta}_k
=
(V_0,\vartheta_k),
\quad
V_0\in\mathbb R,
\quad
\vartheta\in\Theta
\right\}\). This now gives a full account for the treatment of mean-variance hedging in the setup of \parencite{Buehler2018DeepHedging}.
\bigbreak
We remark that in the current setting, \cref{prop:finite deep mean variance hedging} can also be proved directly by observing that \(\overline{\Theta_{\infty}}^{\lVert\cdot\rVert}=\Theta\) and that the map \(\vartheta\lVert H-V_0-(\vartheta\bullet X)_T\rVert_{L^2}\) is continuous with respect to the induced by the norm \(\lVert\cdot\lVert\). For details we refer to the appendix \cref{Appendix:Direct proof of Finite Deep mean variance Hedging}
\bigbreak
Lastly, notice the same conditions on the cost process \(C\) (upper semi-continuity), and the constraining map \(H\)(continuity) in \parencite{Buehler2018DeepHedging}, all the results generalize to the general case with a non-zero cost process and the constrained hedging problem \parencite[Lemma~4.4.]{Buehler2018DeepHedging}, when going through the exact argument given in the proof of \parencite[Prop.~4.9.]{Arandjelovic2025}, using that, by the observation \cref{eq:Continuity of L^2 norm}.
\chapter{Neural Networks and Universal Approximation}\label{ch:Neural Networks and Universal Approximation}
\section{Feedforward neural networks}
\subsection{Architectures and notation}

\section{Universal approximation}
\subsection{Cybenko's theorem:\((C(K),||\cdot||_{\infty})\)and \((L^p(\mu),||\cdot||_{p})\)-density}
\subsection{Discriminatory activation functions}
Based on \parencite[Thm.~1,2]{Hornik1991}
\subsection{Generalization of discriminatory property to locally convex spaces}
Assumption on activation function and \parencite[Assumptions~5.13.]{Hornik1991}. Then {Hornik1991}\parencite[Prop.~5.15.]{Hornik1991}\\
Show that the generalized neural network, with an activation function obeying the same assumptions as in \parencite{Cybenko1989}, is still discriminatory; hence, the density results still follow.
\subsection{Neural approximation on locally convex spaces}
\parencite[Thm~5.16.]{Hornik1991}
Emphasize the role of Hornik's and Cybenko's arguments, and in particular, remark that we will only use the case \(X = \mathbb{R}^{k}\) (e.g. \(k = n\cdot d\)) in what follows.

\chapter{Continuous-time Mean-Square Hedging}
In this section, we discuss in more detail hedging, more precisely mean-variance hedging, in the continuous time setup and the related notions required for this generalization. In particular, we shall assume a filtered probability space \((\Omega,\mathcal{F},\mathbb{F},\mathbb{P})\) with filtration \(\mathbb{F}=(\mathcal{F_t})_{t\in\mathbb{R}_+}\) satisfying the usual conditions. In particular, when we speak of martingales, we shall take its càdlàg version.
\section{Self financing strategies in continuous times}\label{sec:Self financing strategies in continuous times}
\textcolor{blue}{
Explain briefly its definition and provide a short discussion about stochastic integration. For the square integrable case, briefly discuss the space \(\mathcal{M}^2,L^2(X)\) and the construction of the stochastic integral via Itô-Isometry. \\
For the general semimartingale case, give only most important result (similar to Rheinländer and Sexton) and give the ucp continuity thoerem. Maybe give some motivation why it is crucial to consider predictable processes as integrands.}
\bigbreak
We briefly introduce the self financing strategies discussed in \cref{eq:self financing strategy} with more details in the context of continuous-time mathematical finance. We concentrate on the objects introduced, namely the stochastic integrals, and only briefly allude to the conceptually important restriction leading to the subset of integrable processes called admissible strategies. For the latter, we follow \parencite{RheinlaenderSexton2011HedgingDerivatives} in that we define a baseline notion of an admissible trading strategy and specify modified notions of admissibility for particular settings, although always explicitly. 
In order to prevent weird scenarios such as the "suicide strategy" (see \parencite[Example 3.3]{RheinlaenderSexton2011HedgingDerivatives}), where one is almost guaranteed positive gains but potentially requires infinite resources to sustain the strategy, we require that, for price processes that are local martingales the value process for all admissible strategies shall be super-martingales. The property of being a (super-)martingale depends on the underlying probability measure. Hence, to prevent ill behaved scenarios like the suicide strategy, we shall require that the value process of an admissible strategy is a super-martingale under all equivalent measures for which the price process is itself a local martingale. We denote the set of equivalent probability measures for which the price process is a local martingale by \(\mathcal{M_e}\), and call it \textit{equivalent martingale measures}. We can now state our "baseline" definition for an admissible strategy
\begin{definition}\label{def:admissible trading strategy}
    Assume \(\mathcal{M}_e\neq\emptyset\). A strategy \(\vartheta\) is an \textit{admissible strategy} if the associated values process \(V\) is a \(\mathbb{Q}\)-super-martingale for every measure \(\mathbb{Q}\in\mathcal{M}_e\)\;. 
\end{definition}
In the continuous case, the gains process of a self-financing strategy \cref{eq:self financing strategy} is given by the stochastic integral
\begin{equation}
    V_t = V_0 + \int_0^t \vartheta_s\diff X_s\;,
\end{equation}
Hence, we additionally require in this case that \(\vartheta\) is \(X\)-integrable. We shall provide a sufficient condition for integrability in different cases. Without going into further detail on arbitrage theory, recall that the condition \(\mathcal{M}_e\neq\emptyset\) implies that there are no arbitrage opportunities with admissible strategies. We will now go into more detail about the construction of the stochastic integral \(\int\vartheta_s\diff X_s\) for different integrands \(\vartheta\) and integrators \(X\). 
\section{Square integrable martingale case}
\textcolor{blue}{
In this section, we will solve the mean-variance problem (...) in the Itô-isometry setting, which assumes the price process \(X\) to be a locally square integrable martingale and that the integrands are in \(L^2(X)\). One might oppose to this approach in the sense that it is very restrictive. Nevertheless, from a theoretical point of view the solution will be short and insightfull and even from a practical view one can make sense of the integrability condition of the price process as remarkes in \parencite[Section~3.1.]{RheinlaenderSexton2011HedgingDerivatives}. Indeed this amounts to solving the mean-vairance problem under a risk-free measure \(\mathbb{Q}\) instead of the real-world statistical measure \(\mathbb{P}\).\footnote{We assume an incomplete market without arbitrage hence there exist multiple equivalent martingale measures}}
In this section, we sketch the construction of the stochastic integral for square integrable price processes and apropriate integrands. For this we define the space \(\mathcal{H}^p\) and call elements in \(\mathcal{H}^p\) p-integrable processes.
\begin{definition}\label{def:H^p}
    Let \(p\in[1,\infty)\), and write
    \begin{equation}
        \lVert M\rVert_{\mathcal{H}^p}\coloneqq \lVert M_{\infty}^*\rVert_{p}=\mathbb{E}[\sup_t|M_t|^p]^{1/p}\;,
    \end{equation}
    where we use the notation \(M^*_{\infty}\coloneqq\sup_{t\geq 0}|M_t|\) and \(\lVert\cdot\rVert_p\) denotes the \(L^p\)-norm. The \(\mathcal{H}^p\) is the space of martingales such that 
    \begin{equation}
         \lVert M\rVert_{\mathcal{H}^p}<\infty\;.
    \end{equation}
\end{definition}
We shall denote \(\mathcal{H}^p_0\) processes \(X\) in \(\mathcal{H}^p\) with \(X_0=0\). The following lemma will be crucial in characterizing the space \(\mathcal{H}^p\).
\begin{lemma}\label{lemma:order of H^p spaces}
 Let \(\mathcal{H}^p\) be as in \cref{def:H^p}. Then \(\mathcal{H}^p\) has the following property
    \begin{enumerate}
        \item For \(q\leq p\), then \(\lVert M\rVert_{\mathcal{H}^q}\leq\lVert X\rVert_{\mathcal{H}^p}\) thus \(\mathcal{H}^p\subset\mathcal{H}^q\). In particular, \(\mathcal{H}^p\subset\mathcal{H}^1\) for all \(p\in[1,\infty)\)\;.
        \item If \(p\in(1,\infty)\) and \(M\) are uniformly integrable, then 
        \begin{equation}
            \lVert M_{\infty}\rVert_p\leq\lVert M_{\infty}^*\rVert_p\leq q\lVert M_{\infty}\rVert_p
        \end{equation}
        for \(p,q\) Hölder-conjugates, that is \(\frac{1}{p}+\frac{1}{q}=1\)\;.
    \end{enumerate}
\end{lemma}
\begin{proof}
    (1) is the usual inclusion of \(L^p(\mu)\) spaces for finite measure \(\mu\) and follows by Jensen's inequality (...). For (2), first notice that by the Martingale convergence theorem (...) \(M_{\infty}\coloneqq\lim_{n\to\infty}M_n\) is well defined as an almost sure limit. The first inequality is obvious, and the second follows by Doob's \(L^p\) inequality. 
\end{proof}
For the following lemma, we need an easy criterion for which a local martingale is actually a true martingale. Furthermore, it gives sufficient conditions for a local martingale to be a uniformly integrable martingale
\begin{lemma}
    Let \(X\) be a local martingale such that \(\mathbb{E}[X_t^*]<\infty\) for all \(t\in\mathbb{R}_+\). Then \(X\) is a martingale. If \(E[X_\infty^*]<\infty\), then \(X\) is a uniformly integrable martingale. 
\end{lemma}
For a proof, see \parencite[Thm.~5.1, II]{Protter2005StochasticIntegration}.
\begin{lemma}
    For all \(p\in[1,\infty)\), we have that \(X\in\mathcal{H}^p\) is uniformly integrable and \(|M_t|^p\to|M_{\infty}|^p\)in \(L^1\).
\end{lemma}
\begin{proof}
    For the first part, it suffices to notice that \(X\in\mathcal{H}^p\) for \(p\in[1,\infty)\) implies  \(E[X_\infty^*]<\infty\), by point (1) of \cref{lemma:order of H^p spaces}. Thus \(M\) is uniformly integrable and \(M_{\infty}\) exists as an almost sure limit \(\lim_{n\to \infty}M_t=M_{\infty}\). Now, because \(|M_t|^p\leq (M_{\infty}^*)^p\) and \((M_{\infty}^*)^p\) are integrable, the family \((|M_t|^p)_{t\in\mathbb{R}_+}\) is uniformly integrable by \parencite[Example.~2.5.3.]{CohenElliott2015StochasticCalculusApplications} for \(\phi(\cdot)=|\cdot|^p\). Hence, the almost sure convergence of \(|M_t|^p\) to \(|M_{\infty}|^p\) implies the respective \(L^1\) convergence of the latter.
\end{proof}
Now, without proof, we state the following Theorem. A proof is given in \parencite[Lemma~10.1.5.]{CohenElliott2015StochasticCalculusApplications}.
\begin{theorem}\label{thm:H^p is a Banach space}
    Identifying indistinguishable martingales, \(\mathcal{H}^p\) is a Banach space under its norm for all \(p\in[1,\infty)\)\;.
\end{theorem}
We can use our general insight from \cref{lemma:order of H^p spaces} and \cref{thm:H^p is a Banach space} to identify \(\mathcal{H}^p\) with the Banach space \(L^p(\Omega,\mathcal{F},\mathbb{P})\). Indeed, the identification is explicitly provided by the following bijective mapping
\begin{align}\label{eq:H^p identification}
    &\mathcal{H}^p\to L^p(\Omega,\mathcal{F}_{\infty},\mathbb{P})\\
    &M\mapsto M_{\infty}
\end{align}
First of all, because the processes in \(\mathcal{H}^p\) are uniformly integrable, the map is well defined. Injectivity follows from the fact that \(M\in\mathcal{H}^p\) is uniformly integrable, and hence by (....) the limit \(M_{\infty}=\lim_{n\to\infty} M_t\) exists and \(M_t = \mathbb{E}[M_{\infty}|\mathcal{F}_t]\). Surjectivity follows by setting \(M_t = \mathbb{E}[X|\mathcal{F}_t]\in\mathcal{H}^p\) for \(X\in L^p(\Omega,\mathcal{F}_{\infty},\mathbb{P})\). Moreover, \cref{lemma:order of H^p spaces} (2) states that the norm on \(\mathcal{H}^p\) is equivalent to the norm on \(L^p\) under the stated identification. In particular, this implies that the dual of \(\mathcal{H}^p\) can be \(\mathcal{H}^q\) where \(p,q\) are Hölder conjugates. More explicitly, any continuous linear form \(\phi\in(\mathcal{H}^p)^*\) can be identified with some \(N\in\mathcal{H}^q\) and written as \(\phi(M)=\mathbb{E}[M_{\infty}N_{\infty}]\). 
\\\\
As usual for \(L^p\) spaces, the case \(p=2\) is particularly interesting as it yields a Hilbert space structure and a rich geometry. According to the identification \cref{eq:H^p identification}, \(\mathcal{H}^2\) can be identified with the Hilbert space \(L^2(\Omega,\mathcal{F}_{\infty},\mathbb{P})\). Explicitly, we can define an inner product on \(\mathcal{H}^2\) that is given by 
\begin{equation}\label{eq:scalar product on H^2}
    (M,N)_{\mathcal{H}^2}=\mathbb{E}[M_{\infty}N_{\infty}]\;.
\end{equation}
This inner-product structure on \(\mathcal{H}^2\) induces a natural notion of orthogonality. We call \(X,Y\in\mathcal{H}^2\) \texit{weakly orthogonal} if \((X,Y)_{\mathcal{H}^2}=0\). Why we refer to this notion of orthogonality as "weak" is due to the fact that we will later encounter a stronger notion of orhtogonality on the subspace \(\mathcal{H}_0^2=\left\{X\in\mathcal{H}^2:X_0=0\right\}\) which will imply the latter. We finish this subsection with the following standard convetion: For simplicity, we will call elements in \(\in\mathcal{H}^2\) square integrable martingales.

\subsection{Stochastic Integration for square integrable Martingales}\label{sec:Stochastic Integration for square integrable Martingales}
The strategy will be as follows: it will mimic classic Lebesgue integration theory but with a few important caveats. In particular, we will define the stochastic integral for a set of simple functions on the product space \(\Omega\times [0,\infty)\). This is possible for general semimartingale integrators \(X\). Due to the representation of the semimartingales \(X\), we observe that the main caveat is the definition of the stochastic integral for square integrable martingales. Hence, we study the latter in more detail. The crucial idea to notice is that for simple integrands and square integrable integrators, we find that the integral map defines an isometric embedding. The latter is the key property enabling us to extend the stochastic integral to a larger class of integrands in a way that preserves the isometry property.
\\\\
Consider a general semimartingale \(X\). By definition a semimartingale can decomposed as 
\begin{equation}\label{eq:semimartingale}
    X = X_0+M+A\;,
\end{equation}
where \(M\in\mathcal{M}_{0,loc}\): local martingale starting at \(0\) and \(A\in FV_{0}\):finite variation process starting at \(0\). We first define the stochastic integral for "simple" integrands. First, notice that due to \parencite[Thm.~85,VI]{DellacherieMeyer1982ProbabilitiesPotentialB}\footnote{\parencite[Thm.~85,VI]{DellacherieMeyer1982ProbabilitiesPotentialB} asserts that every local martingale can be decomposed as a sum of a local martingale with bounded jumps and a local martingale of finite variation. The later can be absorbed in the finite variation part of \cref{eq:semimartingale}.}, we can choose the local martingale \(M\) to have bounded jumps, hence locally square integrable. Assume that integration with respect to locally square integrable processes \(M\) is understood for a respective class of integrable processes. Up to the issue of the uniqueness of decomposition \cref{eq:semimartingale} we naturally define the integral of some integrable process \(H\) with respect to \(X\) as 
\begin{equation}\label{eq:heuristic integral with respect to semimartingale}
    \int H\diff X\coloneqq\int H\diff M+\int H\diff A\;.
\end{equation}
In particular integration with respect to finite variation measures is well understood via the Lebesgue-Stieljes integral. Hence, the main conceptual gap in \cref{eq:heuristic integral with respect to semimartingale} is, up to localization, the construction of \(\int H\diff M\) for a square integrable martingale \(M\). In what is to come we will proceed with the definition of the stochastic integral of "simple" integrands and then specialized to square integrable martingales as integrators. 
\begin{definition}\label{def:elementary predictable process}
    We call \(H\) an \textit{elementary predictable process} if there exists a finite, increasing sequence of deterministic times, \(0=t_0<\dots<t_{n+1}<\infty\) and a sequence of bounded random variables \((H^i)_{i=0}^n\subset\mathcal{L}^{\infty}(\mathcal{F}_{t_i},\mathbb{P})\) such that 
    \begin{equation}
        H = H^01_{\left\{0\right\}}+\sum_{i=1}^nH^i1_{(t_i,t_{i+1}]}\;.
    \end{equation}
In particular, we have that \(H^i\) is \(\mathcal{F}_{t_i}\)-measurable. We denote the set of elementary predictable processes by \(\mathcal{E}\) and we say \(H\in\mathcal{E}\) is associated with a sequence \((t_k)_{k=0}^{n+1}\).
\end{definition}
Notice that elementary processes are, as defined, left-continuous processes. In particular, observe that clearly \(\mathcal{E}\) is closed under summation, but even more, it is also closed under multiplication, as it is closed under multiplication. To see this, take two elements of \(H,G\in\mathcal{E}\) and their assosciated sequences \((t^H_k)_{n=0}^{n+1}\), \((t^G_k)_{k=0}^{l+1}\) for some \(n,l\in\mathbb{N}\). They both form a partition of \([0,\infty)\) if we add one point \(t_{n+2}=t_{l+2}=\infty\) to both sequence. Take their common refinement and notice that the pointwise product of two indicator function of disjoint sets vanishes.\\\\
For later use, we denote by \(\mathcal{E}^{\infty}\) the set of processes as in \cref{def:elementary predictable process} but for infinite sets \((t_k)_{k\in\mathbb{N}}\) where \(0=t_0<\dots,t_n<\infty\) and \(\lim_{n\to\infty}t_n=\infty\). Hence processes of the form 
\begin{equation}
      H = H^01_{\left\{0\right\}}+\sum_{i\geq 1}^nH^i1_{(t_i,t_{i+1}]}\;.
\end{equation}
We identify a process in \(\mathcal{E}\) associated with the sequence \((t_k)_{k=0}^n\), with the process in \(\mathcal{E}^{\infty}\), the sequence \(\tilde{t}_t=t_k\quad k=0,\dots,n\), and \(t_k = t_n\) for \(k>n\), following the convention \(1_{(t,t]}=1_{\emptyset}=0\), thus identifying \(\mathcal{E}\) as a subset of \(\mathcal{E}^{\infty}\).
\bigbreak
The corresponding integral of \(H\) with respect to \(X\) is then naturally defined as  
\begin{definition}\label{def:integral of simple processes}
Let \(H\in\mathcal{S}\). Then we set 
\begin{equation}\label{eq:integral of simple processes}
    I_X(H)=H^0X_0+\sum_{i}^n H^i(X_{\tau_{i+1}}-X_{\tau_{i}})\;.
\end{equation}
Furthermore, for every \(t\in\mathbb{R}_+\), we define
    \begin{equation}
        I_X(H)_t=H^0X_0+\sum_{i}^n H^i(X_{\tau_{i+1}\wedge t}-X_{\tau_{i}\wedge t})\;.
    \end{equation}
    We call \(I_X(H)\) the stochastic integral process of \(H\) with respect to \(X\) and use the following notations interchangeably
    \begin{equation}
        I_X(H)=H\bullet X=\int H\diff X\;.
    \end{equation}
\end{definition}
Notice that the expressions \cref{eq:integral of simple processes} are pathwise and hence independent of the representation of \(H\in\mathcal{S}\).
\bigbreak
We now turn to the integration for \(X\) a square integrable martingale. The key observation is the so called \textit{Itô Isometry} , which we state below. 
Recall that for \(X\in\mathcal{H}^2\) one can define its quadratic variation process \(\langle X\rangle\) (see \parencite[Def.11.1.1.]{CohenElliott2015StochasticCalculusApplications}. The crucial observation is the following.
\begin{proposition}\label{prop:Ito isometry for simple processes}
    For \(H\in\mathcal{E}\) and \(X\in\mathcal{H}_0^2\), we have that \(\int H\diff X\in\mathcal{H}^2_0\) and the property
    \begin{equation}\label{eq:Ito isometry for simple processes}
        \mathbb{E}[\left(\int_0^{\infty} H_s\diff X_s\right)^2]=\mathbb{E}[\int_0^{\infty}H_s^2\diff\langle M\rangle_s]\;,
    \end{equation}
    where \(\int_0^{\infty}H_s^2\diff\langle M\rangle_s\) denotes the Lebesgue-Stieltjes integral for the finite variation process \(\langle M\rangle\).
\end{proposition}
The statement involves two things. First, it postulates the equality \cref{eq:Ito isometry for simple processes}, and secondly, it assumes that the terms in \cref{eq:Ito isometry for simple processes} are finite. 
The proof of the martingale property, as well as the final equality \cref{eq:Ito isometry for simple processes}, is proven by direct calculation (see \parencite[Lemma.~12.1.4]{CohenElliott2015StochasticCalculusApplications}).
The main idea is that one wants to elevate \cref{prop:Ito isometry for simple processes} to an isometry statement. Recall that, knowing in the setting of \cref{prop:Ito isometry for simple processes}, we have \(\int_0^{\infty} H_s\diff X_s\in\mathcal{H}_0^2\). From our identification, \cref{eq:scalar product on H^2} we have
\begin{equation}
    \mathbb{E}[\left(\int_0^{\infty} H_s\diff X_s\right)^2]=\lVert\int H_s\diff X_s\rVert_{\mathcal{H}^2}\;.
\end{equation}
Thus we want to understand \cref{prop:Ito isometry for simple processes} as the statement that the map 
\begin{align}\label{eq:integral map as isometry}
     I:&H\mapsto \int H\diff X\\
     &\mathcal{L}\to\mathcal{H}_0^2\;,
\end{align}
is an linear isometry between \(\mathcal{L}\) and \(\mathcal{H}_0^2\) for some appropriate space normed vector \(\mathcal{L}\). From \cref{eq:Ito isometry for simple processes} the natural candidate is the following space:
\begin{definition}\label{def:L^2(M)}
    Let \(X\in\mathcal{H}^2\). We denote \(L^2(X)\) as the space of predictable processes \(H\) such that 
    \begin{equation}
        \lVert H\rVert_{L^2(M)}=\mathbb{E}[\int_0^{\infty}H_s^2\diff\langle M\rangle_s]<\infty\;.
    \end{equation}
\end{definition}
Hence \cref{prop:Ito isometry for simple processes} reads exactly as the statement that the map \(I:L^2(X)\cap\mathcal{E}\to\mathcal{H}^2_0\) given by \(H\mapsto \int H\diff X\) as defined in \cref{eq:integral of simple processes} is an isometry of normed vector spaces. This important structural insight will enable us to extend the stochastic integral to all of \(L^2(X)\). Indeed we will show that \(\mathcal{E}\) is a dense subset of \(L^2(X)\), hence we define the stochastic integral of processs in \(L^2(X)\) for \(X\in\mathcal{H}_0^2\) as the unique continuous extension of the linear, bounded map \cref{eq:integral map as isometry} for \(\mathcal{L}=L^2(X)\).
\bigbreak
We first make a crucial observation, justifying the choice of notation \(L^2(X)\). For this we define the following measure on the measure space \(\left(\Omega\times[0,\infty),\mathcal{F}\otimes\mathcal{B}([0,\infty)\right)\):
\begin{equation}
    \mathbb{P}_M(A)=\mathbb{E}[\int_0^{\infty}1_A\diff\langle M\rangle_s]\quad\forall A\in\mathcal{F}\otimes\mathcal{B}([0,\infty))\;.
\end{equation}
 Then we have, by \cref{def:L^2(M)}:
\begin{equation}\label{eq:L^2(X) as Hilbert space}
    L^2(M) = L^2(\Omega\times [0,\infty),\mathcal{P},\mathbb{P}_M)\;,
\end{equation}
where \(\mathcal{P}\) is the sub-\(\sigma\)-algebra of \(\mathcal{F}\otimes\mathcal{B}[0,\infty)\) generatey by adapted, left-continuous processes. As usual \cref{eq:L^2(X) as Hilbert space} is to be understood as the quotient space \(\mathcal{L}^2(\mathcal{P},\mathbb{P}_M)/\sim\) with \(X\sim Y\iff \lVert X-Y\rVert_{L^2(M)}=0\), making \(\lVert \cdot\rVert_{L^2(M)}\) a well real norm on \(L^2(M)\). In particular this implies \(L^2(M)\) is a Hilbert space with inner product \(\langle X,Y\rangle_{L^2(M)}=\mathbb{E}[\int_0^{\infty}X Y\diff\langle M\rangle_s]\) and in particular \(L^2(M)\) is complete. 
Consider the following usefull lemma, assuring that for \(M\in\mathcal{H}_0^2\) the measure \(\mathbb{P}_M\) is finite.
\begin{lemma}\label{eq:P_M is finite measure}
    Suppose \(M\) is a local martingale starting at \(0\). Then \(M\in\mathcal{H}^2_0\) if and only if \(\mathbb{E}[\langle M\rangle_{\infty}]<\infty\)
\end{lemma}
\begin{proof}
    See \parencite[]{}
\end{proof}
One can show that for the process \(\langle M\rangle\) exists if and only if \(M\in\mathcal{H}_{0,loc}^2\) where \(\mathcal{H}_{0,loc}^2\) is the space of processes that are locally in \(\mathcal{H}_{0}^2\). Observe that the space of locally continuous martingales, denoted \(\mathcal{M}_{0,loc}^c\) is a subset of \(\mathcal{H}_{0,loc}^2\). 
\bigbreak
Now we can proof the processes in \(L^2(M)\) are approximated by elementary predictable process. The proof will be based on an application of the functional version of the monotone class theorem (\parencite[Thm.~7.4.1.]{CohenElliott2015StochasticCalculusApplications}) which exploits the fact that \(\sigma(\mathcal{E})=\mathcal{P}\). We will actually show that \(\sigma(b\mathcal{E})=\mathcal{P}\), where \(b\mathcal{E}\) denotes the processes in \(\mathcal{E}\) that are bounded. The next lemma shows that this is equivalent to the full statement \(\sigma(\mathcal{E})=\mathcal{P}\).
\begin{lemma}\label{lemma:sigma algebra of bounded functions}
    Let \(\mathcal{M}\) be a class of real-valued functions on \(\Omega\). We denote by \(b\mathcal{M}\) the class of bounded functions in \(\mathcal{M}\). Then
    \begin{equation}
        \sigma(b\mathcal{M})=\sigma(\mathcal{M})\;.
    \end{equation}
\end{lemma}
\begin{proof}
    The inclusion \(\sigma(b\mathcal{M})\subset\sigma(\mathcal{M})\) is clear. For the converse let consdier \(X\in\mathcal{M}\) and define
    \begin{equation}
        X^n = -n\vee X\wedge n\in b \mathcal{M}\;.
    \end{equation}
    Then it is clear that \(X^n\to X\) pointwise \footnote{It should be clear that pointwise means: \(X^n(\omega,t)\to X(\omega,t)\quad\forall (\omega,t)\in\Omega\times[0,\infty)\)} as \(n\to \infty\), hence the claim follows. 
\end{proof}
\begin{lemma}\label{lemma:elementary processes generated predictable sigma algebra}
    For \(\mathcal{E}\) as in \cref{def:elementary predictable process} and \(\mathcal{P}\) the predictable \(\sigma\)-alebgra, we have that 
    \begin{equation}
        \sigma(\mathcal{E})=\mathcal{P}\;.
    \end{equation}
    \begin{proof}
        By definition, \(\mathcal{P}\) is generated by the family of adapted, left-continuous processes. Every element of \(\mathcal{E}\) is clearly adapted and left-continuous hence \(\sigma(\mathcal{E})\subset\mathcal{P}\) is clear. Hence it remains to show that \(\mathcal{P}\subset\sigma(\mathcal{E})\). For this we  first show that any adapted left-continuous process is the limit of elements of \(\mathcal{E}^{\infty}\). Let \(X\) be an adapted left-continuous process. By \cref{lemma:sigma algebra of bounded functions} we might assume \(X\) to be bounded. We define
        \begin{equation}\label{eq:approximating adapted left-continuous process}
            X^n= X_01_{\left\{0\right\}}+\sum_{k\geq 1}X_{k/n}1_{(k/n,(k+1)/n]}
        \end{equation}
        As \(X\) is adapted, \(X_{k/n}\) is bounded and \(\mathcal{F}_{k/n}\)-measurable; hence, indeed \(X^n\in\mathcal{E}^{\infty}\). Then for every fixed \(t\in[0,\infty)\) there exists a unique \(k(n)\in\mathbb{N}\) such that \(\frac{k(n)}{n}<t\leq \frac{k(n)+1}{n}\). Hence \(X_t^n=X_{k(n)/n}\), thus \(X_t^n\) is \(\mathcal{F}_{k/n}\subset\mathcal{F}_t\)-measurable, hence adapted. Furthermore, notice that \(\lim_{n\to\infty}\frac{k(n)}{n}\leq\lim_{n\to\infty}|\frac{k(n)}{n}-t|+t\leq\lim_{n\to\infty}\frac{1}{n}+t=t\) and \(\frac{k(n)}{n}<t\); hence \(k(n)\nearrow t\). Thus, for every \((\omega,t)\in\Omega\times[0,\infty),\;X^n_t(\omega)\to X_t(\omega)\) as \(n\to\infty\) by the left-continuity of \(X\). This implies 
        \begin{equation}
            \mathcal{P}\subset\sigma(\mathcal{E}^{\infty})\;.
        \end{equation}
        Indeed, every adapted, left-continuous process \(X\) is \(\sigma(\mathcal{E}^{\infty})\) as the point-wise limit of elements in \(\mathcal{E}^{\infty}\). But \(\mathcal{P}\) is the smalles such \(\sigma\)-algebra.\\\\
        Furthermore every \(H\in\mathcal{E}^{\infty}\) is trivially the pointwise limit of processes
        \begin{equation}\label{eq:approximating approximating-sequence of adapted left-continuous processes}
            H^{m}=H^01_{\left\{0\right\}}+\sum_{k=1}^mH^{i}1_{(t_k,t_{\frac{k+1}{n}}]}\in\mathcal{E}\;,
        \end{equation}
        for \(m\to\infty\).
        By the same argument as before we have
        \begin{equation}
        \mathcal{P}\subset\sigma(\mathcal{E}^{\infty})\subset\sigma(\mathcal{E})\;.
        \end{equation}
        \end{proof}
\end{lemma}
Taking a closer look we see that the previous proof shows the following corollary
\begin{corollary}
    Let \(\mathcal{A}=\left\{A\times(s,t], A\in\mathcal{F}_s, s<t\in\mathbb{R}_+\right\}\cup \left\{\left\{0\right\}\times A, A\in\mathcal{F}_0\right\}\). Then \(\sigma(\mathcal{A})=\mathcal{P}\).
\end{corollary}
\begin{proof}
    Every indicator function \(1_A\) with \(A\in\mathcal{A}\) is adapted and left-continuous hence predictable. The converse inclusion follows from the very same argument as in \cref{lemma:elementary processes generated predictable sigma algebra}. Indeed every adapted left-continuous process \(X\) is the limit of processes \(X^n\) of the form \cref{eq:approximating adapted left-continuous process}. Every \(X^n\) is the limit of processs \(X^n_m\) of the form \cref{eq:approximating approximating-sequence of adapted left-continuous processes}. The latter is clearly seen to be \(\sigma(\mathcal{A})\)-measurable as \(X_{k/n}\) is \(\mathcal{F}_{k/n}\). Indeed for \(B\in\mathcal{B}(\mathbb{R})\)
    \begin{align}
    (X_{t/k}1_{(t/k,k+1/n]})^{-1}(B)=
    \begin{cases}
        X_{t/k}^{-1}(B)\times (k/n,k+1/n] \text{ if }0\notin B\\
        X_{t/k}^{-1}(B)\times (k/n,k+1/n]\cup \Omega\times \mathbb{R}/(k/n,k+1/n]\text{ if }0\in B
    \end{cases}
    \end{align}
    but \( \Omega\times \mathbb{R}/(k/n,k+1/n]=(\Omega\times(k/n,k+1/n])^c\in\sigma(\mathcal{A})\). In every case \\
    \((X_{t/k}1_{(t/k,k+1/n]})^{-1}(B)\in\sigma(\mathcal{A})\).
    Thus \(X^n\) is \(\sigma(\mathcal{A})\)-measurable as a pointwise limit of measurable processes and \(X\) is \(\sigma(\mathcal{A})\)-measurable by the same reasoning. 
\end{proof}
We denote \(b\mathcal{P}\) bounded predictable processes. We can finally state and proof the density of elementary processes in \(L^2(M)\).
\begin{proposition}\label{prop:density of E in L2(M) via monotone class}
    Let \(M\in\mathcal{H}^2_0\) and \(\mathcal{E}\) as defined in \cref{def:elementary predictable process}. Then \(\mathcal{E}\) is dense in \(L^2(M)\)\;.
\end{proposition}
\begin{proof}
    The proof uses essentially the functional version of the monotone class theorem and dominated convergence. By \cref{lemma:sigma algebra of bounded functions} and \cref{lemma:elementary processes generated predictable sigma algebra} we have 
    \begin{equation}
        \sigma(b\mathcal{E})=\mathcal{P}\;.
    \end{equation}
    Let 
    \begin{equation}
        \mathcal{H}=\left\{ H\in b\mathcal{P}:H\in L^2(M)\;,\forall \epsilon>0\;, \exists H_{\epsilon}\in b\mathcal{E}\text{ with }\lVert H_{\epsilon}-H\rVert_{L^2(M)}<\epsilon\right\}\;.
    \end{equation}
    Suppose we where able to show that \(\mathcal{H}\) was a monotone class containing the algebra of functions \(b\mathcal{E}\). We would conclude by the monotone class theorem (\parencite[Thm.~7.4.1.]{CohenElliott2015StochasticCalculusApplications}) that \(\mathcal{H}\) containes all bounded \(\sigma(b\mathcal{E})=\mathcal{P}\)-measurable maps, which would show that \(b\mathcal{E}\) is dense in \(L^2(b\mathcal{P},\mathbb{P}_M)\). To get density in \(L^2(\mathcal{P},\mathbb{P}_M)\) we observe that \(L^2(b\mathcal{P},\mathbb{P}_M)\)  is dense in \(L^2(\mathcal{P},\mathbb{P}_M)\). Indeed for every \(H\in L^2(b\mathcal{P},\mathbb{P}_M)\) we have that by dominated convergence that \(H^n=-n\vee H\wedge n\in L^2(b\mathcal{P},\mathbb{P}_M)\) converges to \(H\) in \(L^2(M)\). 
    We proceed by showing that we can rightfully apply the montone class theorem. \\\\
    First, by \cref{eq:P_M is finite measure}, \(\mathbb{P}_M\) is a finite measure, hence \(\mathcal{H}\) contains \(b\mathcal{E}\) and hence also \(1_{\Omega\times[0,T]}\). and forms a vector space uf functions. Indeed if \(H_1,H_2\in\mathcal{H}\), for every \(\epsilon>0\) there exist \(H^1_{\frac{\epsilon}{2}},H^2_{\frac{\epsilon}{2|\alpha|}}\) such that, for all \(\alpha\in\mathbb{R}\):
    \begin{equation}
        \lVert H^1_{\frac{\epsilon}{2}}+\alpha H^2_{\frac{\epsilon}{2|\alpha|}} \rVert_{L^2(M)}\leq \epsilon.
    \end{equation}
    It remainse to show that \(\mathcal{H}\) is closed under bounded monotone limits. Hence let \(0<H_0<H_1<\dots\) with \(H_j\in\mathcal{H}\) and  uniformly bounded in \(j\) with \(\lim_{n\to\infty}H_n = H\) pointwise. Fix \(\epsilon>0\). Then due to dominated convergence 
    \begin{equation}
        \lim_{n\to\infty}\lVert H-H_n\rVert_{L^2(M)}=0\;.
    \end{equation}
    Hence there exists \(n\in\mathbb{N}\) large enough such that \(\lVert H-H_n\rVert_{L^2(M)}<\epsilon/2\). Now \(H_n\in\mathcal{H}\), hence there exists \(H^n_{\epsilon/2}\in b\mathcal{E}\) such that \(\lVert H_n-H^n_{\epsilon/2}\rVert_{L^2(M)}<\epsilon/2\). By the triangle inequality we have
    \begin{equation}
        \lVert H-H^n_{\epsilon/2}\rVert_{L^2(M)}\leq\lVert H-H_n\rVert_{L^2(M)}+\lVert H_n-H^n_{\epsilon/2}\rVert_{L^2(M)}<\epsilon
    \end{equation}
    Hence we conclude by the monotone class theorem. 
\end{proof}
Notice that in this proof goes through analogously. Indeed for \(p\in[1,\infty)\) and \(\mathcal{H}^p\) we define the according seminorm
\begin{equation}
    \lVert H\rVert_{L^p(M)}=\mathbb{E}[\left(\int_0^\infty |H_s|^p\diff \langle M\rangle_s\right)^{p/2}]
\end{equation}. The truncation argument works by two succesive applications of dominated convergence showing that \(V^n=-n\vee V\wedge n\to V\) in \(L^p(M)\) we , hence we get the following proposition 
\begin{proposition}\label{prop:density of E in Lp(M) via monotone class}
    Let \(M\in\mathcal{H}^p_0\) for \(p\in[1,\infty)\) and \(\mathcal{E}\) as defined in \cref{def:elementary predictable process}. Then \(\mathcal{E}\) is dense in \(L^p(M)\)\;.
\end{proposition}
Finally, having proven that the elementary processes are dense in \(L^2(M)\) we can use the isometry property for the map \cref{eq:integral map as isometry} to extend the integral to a bounded linear operator, by the standard extension theorem (see for instance ......).
We get the following theorem.
\begin{theorem}\label{thm:square integrable stochastic integral}
    The linear map \(H\mapsto \int H\diff M\) fo \(\mathcal{E}\) into \(\mathcal{H}_0^2\) extends uniquely to an isometric linear map \(L^2(M)\mapsto\mathcal{H}_0^2\). We call the this extension, the integral of \(H\) with respect to \(M\).
\end{theorem}
Notice that as we define the stochastic integral as the extension of a map \(I:\mathcal{E}\to\mathcal{H}_0^2\) to \(L^2(M)\) where \(\mathcal{E}\subset L^2(M)\) is a dense susbet, the resulting integral is in \(\mathcal{H}_0^2\) by construction. 
\\\\
Furthermore, the isometry property is generally know as Itô's isometry. We state it explictely due to its importance in this chapter and the following chapters
\begin{theorem}[Itô Isometry]\label{thm:Ito Isometry}
    Let \(M\in\mathcal{H}^2_0\), then for \(H\in L^2(M)\)
    \begin{equation}
        \lVert H\rVert_{L^2(M)}=\mathbb{E}[\int_0^{\infty}H_s^2\diff \langle M\rangle_s]=\mathbb{E}[\left(\int_0^{\infty}H_s^2\diff M_s\right)^2]=\lVert \int H_s\diff M_s\rVert_{\mathcal{H}^2}
    \end{equation}
\end{theorem}
We will refer to the setting in \cref{thm:square integrable stochastic integral} as the \textit{square integrable setting}. From here one can start and broaden the space of integrands and integrators. \textcolor{blue}{We will only generlize to \(\mathcal{H}_{0,loc}^2\) integrators as this is all we need. The step from intgrands \(L^2(M)\) to \(L_{loc}^2(M)\) will be scipped as it follows a simple localization argument as well. For the step from \(X\in\mathcal{H}^2\toX\in\mathcal{H}_{loc}^2\) we will need to understand how the stopped integral process looks like. For this we give a characterisation which is independently useful
\begin{proposition}\label{prop:characterisation of stochastic integral}
    Let \(M\in\mathcal{H}_0^2\) and \(H\in L^2(M)\). The stochastic integral \(H\bullet M\) is the unique process in \(\mathcal{H}^2_0\) such that, for every \(N\in\mathcal{H}^2\),
    \begin{equation}
        \mathbb{E}[L_{\infty}N_{\infty}]=\mathbb{E}[\int_0^\infty H_s\diff\langle M,N\rangle_s]\;.
    \end{equation}
\end{proposition}
For the proof we refer to \parencite[Thm.~2.5.]{Schweizer2025BrownianMotionStochasticCalculus}. Conceptually not that the right-hand side is a linear form on in \(N\) that is on \(\mathcal{H}^2\), which has been identified with a Hilbertspace. By the Kunita-Watanabe inequality \cref(...) it is also continuous. Then \(L\) is the unique Riesz-Fischer vector in \(\mathcal{H}^2\) such that the continuous linear form can be written as an inner product. The latter characterises the continous linear form on a Hilbertspace uniquely. 
From \cref{prop:characterisation of stochastic integral} it is now easy to proof the following
\begin{corollary}\label{cor:stopped integral}
    For \(M\in \mathcal{H}^2_0\), \(H\in L^2(M)\) and \(\tau\) a stopping time, we have 
    \begin{equation}
        (H\bullet M)^{\tau}=H\bullet(M)^{\tau}=(H1_{(0,\tau]})\bullet M\;.
    \end{equation}
\end{corollary}
\begin{proof}
    ...
\end{proof}
It is now straight forward to define the stochatic integral for \(M\in\mathcal{H}^2_{0,loc}\) and \(H\in L^2(M)\) (notice that in the definition of \(L^2(M)\) we only need the existence of the process \(\langle M\rangle\). The latter exists if and only if \(M\in\mathcal{H}^2_{0,loc}\). Given a localizing sequence \((\tau_n)_{n\in\mathbb{N}}\) of \(M\) we define 
\begin{align}
    \int H\diff M\coloneqq\lim_{n\to\infty}(\int H\diff M^{\tau_n})\;.
\end{align}}
One can further extend the the class of integrands to the localised space \(L_{loc}^2(M)\) (for more details see \parencite[Prop.~2.20]{Schweizer2025BrownianMotionStochasticCalculus}) when properly defining \(L_{loc}^2(M)\). We shall hence define the class of integrable processes for \(M\in\mathcal{H}_{0,loc}^2\) by 
\begin{equation}
    \mathcal{L}(M)=L_{loc}^2(M)\;.
\end{equation}
We give the definition of \(L_{loc}^2(M)\).
\begin{definition}\label{def:L^2(M)_loc}
    For \(M\in\mathcal{H}^2_{0,loc}\) the space \(L_{loc}^2(M)\) consists of predictable processes \(H\) for which there is  a sequence of stopping times \(\tau_n\nearrow\infty\) almost surely, such that 
    \begin{equation}
        \mathbb{E}[\int_0^{\tau_n}H_s^2\diff \langle M\rangle_s]<\infty\quad\forall n\in\mathbb{N}\;.
    \end{equation}
\end{definition}
Notice that we are saying that for \(M\in\mathcal{H}_{0,loc}^2\), \(H\in L^2(M)\) is equivalent to \(H1_{(0,\tau_n]}\in L^2(M)\) for each \(n\in\mathbb{N}\) and also to \(H\in L^2(M^{\tau_n})\) for each \(n\in\mathbb{N}\), for a localizing sequence \((\tau_n)_{n\in\mathbb{N}}\) in the sense of \cref{def:L^2(M)_loc}. A useful characterisation of \(L^2(M)\) for \(M\in\mathcal{H}_{0,loc}^{2}\) is given as follows
\begin{lemma}
    Let \(M\in\mathcal{H}_{0,loc}^{2}\). Then \(H\in L^2_{loc}(M)\) iff 
    \begin{equation}
        \int_0^t H_s^2\diff\langle M\rangle_s<\infty \quad\mathbb{P}\text{-a.s.}\;,\forall t\geq 0
    \end{equation}
\end{lemma}
We meantioned that the stochastic integral of processes in \(L^2_{0,loc}(M)\) with respect to processes in \(\mathcal{H}_{0,loc}^2\) is defined through localisation of the integral defined in \cref{thm:square integrable stochastic integral}. From this it is clear that by construction the stochsatic integral in the local case is a local martingale. Notice that in general it is not more than that.
\subsection{Kunita-Watanabe Decomposition}
Before turning towards stochastic integration with respect to general semimartingales we profit from the current familiarity of the reader with the square integrable setup to introduce and and state the celebrated Kunita Watanabe decomposition. For this we will need a stronger notion of orthonogality on \(\mathcal{H}^2_0\) then the one induced by the inner product structure. The fact that \(\mathcal{H}^2_0\) is a space of uniformly integrable martingales leeds to a richer geometry than the Hilbertspace structure discussed in \cref{eq:H^p identification}. 
\begin{definition}
    Let \(X,Y\in\mathcal{H}_0^2\). We say that \(X\) is strongly orthogonal to \(Y\) and vise versa if their product \(XY\) is a martingale. In this case, we denote \(X\perp Y\).
\end{definition}
In particular, for \(X\perp Y\),
\begin{equation}
    (X,Y)_{\mathcal{H}_0^2}=\mathbb{E}[X_{\infty}Y_{\infty}]=\mathbb{E}[X_0Y_0]=0\;.
\end{equation}
Hence, strong orthogonality implies the weak orthogonality of \(X,Y\in\mathcal{H}_0^2\). 
\bigbreak
As shortly hinted in \cref{sec:high level discussion of mean-square hedging}, in the context of mean-squared hedging we are generally interested, for a fixed prices process \(X\), what are the strategies for which the value process is again square integrable. In the square integrable setting, we gave a partial answer in affirming that for \(X\) square integrable the space \(L^2(M)\) is in \(G_T(\vartheta)\). In this case, we already argued that the mean-squared hedging problem reduced to a projection problem in \(L^2\). The Kunita Watanabe inquality will give a complete answer for the square integrable setting and even the locally square integrable case. The crucial observation is given by the following \cref{prop:weak and strong stable subspaces}, which necessits the introduction of the following concept
\begin{definition}\label{def:stable subspace}
    A \textit{stable} subspace of \(\mathcal{H}_0^2\) is a closed linear subsapce \(\mathcal{A}\) of \(\mathcal{H}_0^2\) which is closed under stopping: For very \(M\in\mathcal{A}\) and every stopping time \(\tau\), the stopped process \(M^{\tau}\) is in \(\mathcal{A}\).
\end{definition}
The following example of a stable subspace instantly underlines the importance of stable subspaces in our case
\begin{lemma}\label{lemma:stochastic integrals as stable subspace}
    For every \(M\in\mathcal{H}_{0,loc}^2\) the space 
    \begin{equation}
        \mathcal{I}(M)=\left\{\int H\diff M:H\in L^2(M)\right\}
    \end{equation}
    is a stable subspace of \(\mathcal{H}_0^2\). 
\end{lemma}
\begin{proof}
    ....
\end{proof}
\begin{proposition}\label{prop:weak and strong stable subspaces}
    Suppose \(\mathcal{A}\) is a stable subspace of \(\mathcal{H}_0^2\) and denote by 
    \begin{equation}
        \mathcal{A}^{\Delta}=\left\{N\in\mathcal{H}_0^2:(N,M)_{\mathcal{H}^2}=0\quad\forall M\in\mathcal{A}\right\}
    \end{equation}
    the weakly orhtogonal compelement of \(\mathcal{A}\) in \(\mathcal{H}_0^2\). Then \(\mathcal{A}^{\Delta}\) is also strongly orthogonal to \(\mathcal{A}\) and a stable stubapce of \(\mathcal{H}_0^2\). 
\end{proposition}
\begin{theorem}[Kunita-Watanabe Decomposition]\label{thm:KW}
    Fix \(M\in\mathcal{H}_{0,loc}^2\). Then every \(N\in\mathcal{H}_{0}^2\) can be uniquely written as 
    \begin{equation}
        N = \int H\diff M +L
    \end{equation}
    for some \(H\in L^2(M)\) and some \(L\in\mathcal{H}_{0}^2\) with \(L\perp M\). Furthermore \(H=\frac{\diff \langle M,N\rangle}{\diff \langle M\rangle}\) is a predictable density of \(\langle M,N\rangle\) with respect to \(\langle M\rangle\). 
\end{theorem}
\begin{proof}
    ...
\end{proof}
\begin{corollary}
    Fix \(M\in\mathcal{H}_{0,loc}^2\). Then ervery random variable \(F\in L^2(\mathcal{F}_{\infty},\mathbb{P})\) can be uniquely written as 
    \begin{equation}
        F = \mathbb{E}[F|\mathcal{F}_0]+\int_0^{\infty}H_s\diff M_s+L_{\infty}
    \end{equation}
    for some \(H\in L^2(M)\) and some \(L\in \mathcal{H}_{0}^2\) such that \(L\perp M\). 
\end{corollary}
\begin{proof}
    We apply \cref{thm:KW} to the process \(V_t\coloneqq \mathbb
    E[F|\mathcal{F}_t]- E[F|\mathcal{F}_0]\in\mathcal{H}_0^2\) for \(t\in[0,\infty]\). We get 
    \begin{equation}\label{eq:KW of V^F}
        V = V_0 + \int H\diff M+L\;.
    \end{equation}
    Noticing that \(F=\mathbb{E}[F|\mathcal{F}_{\infty}]\) and \(V_0=\mathbb{E}[F|\mathcal{F}_{0}]\) we obtain the claim by considering \cref{eq:KW of V^F} at \(t=\infty\). 
\end{proof}
Notice that if \(\mathcal{F}_0\) is trivial and completed, that is \(\mathcal{F}_0=\left\{\emptyset,\Omega\right\}\vee\mathcal{N}_{\mathbb{P}}\), then \(\mathbb{E}[F|\mathcal{F}_0]=\mathbb{E}[F]\), \(\mathbb{P}\)-almost surely (with exact equality if \(\mathcal{F}_0=\left\{\emptyset,\Omega\right\}\). 
\subsection{Stochastic Integration with respect to Semimartingales}
The construction of the stochastic integral with respect to square integrable martingales constitutes the biggest conceptual development towards stochstic integration with respect to general martingales. The case for bounded predictable integrands can be studied in \parencite[Thm.3 VIII]{DellacherieMeyer1982ProbabilitiesPotentialB}. In view of our developement of \cref{sec:From simple processes to stochastic integral} we take a different route presented by \parencite[Chapter.~4]{Protter2005StochasticIntegration}. For the interested reader we note that the notion of semimartingales used in \parencite[Chapter.~4]{Protter2005StochasticIntegration} is different from the classical one, used in \cref{eq:semimartingale} at first sight. In \parencite[Chapter.~4]{Protter2005StochasticIntegration} are essentially defined to be processes for which integration of simple predictable processes is continuous under uniform convergence. The latter essentially garanties a notion of bounded convergence for the respective integral, which is given a central role in the construction of a sensible integral. It is the celebrated Bitcheler Dellacherie Theorem (see \parencite[Thm.~47]{DellacherieMeyer1982ProbabilitiesPotentialB}) that makes the junction between these two notions. 
\bigbreak
For our purposes we need to enlarge the definition of elementary predictable processes from \cref{def:elementary predictable process} to the notion of \textit{simple predictable processes}. 
\begin{definition}\label{def:simple predictable process}
    We call \(H\) an \textit{simple predictable process} if there exists a finite, increasing sequence of stopping times, \(0=\tau_0<\dots<\tau_{n+1}<\infty\) and a sequence of bounded random variables \((H^i)_{i=0}^n\subset\mathcal{L}^{\infty}(\mathcal{F}_{\tau_i},\mathbb{P})\) such that 
    \begin{equation}
        H = H^01_{\left\{0\right\}}+\sum_{i=1}^nH^i1_{(\tau_i,\tau_{i+1}]}\;.
    \end{equation}
In particular, we have that \(H^i\) is \(\mathcal{F}_{\tau_i}\)-measurable. We denote the set of elementary predictable processes by \(\mathcal{S}\) and we say \(H\in\mathcal{S}\) is associated with a sequence \((\tau_k)_{k=0}^{n+1}\).
\end{definition}
Clearly \(\mathcal{E}\subset\mathcal{S}\). We further denote \(\mathcal{S}_{ucp}\) the set \(\mathcal{S}\) endowed with the topology of \(ucp\)-convergence. Furthermore we dentoe by \(\mathbb{L}\) the class of adapted càglàd processes, \(b\mathbb{L}\) the respective bounded subset and by \(\mathbb{D}\) the class of adapted càdlàg processes. The construction of the stochastic integral \(\int H\diff X\) for \(H\in b\mathbb{L}\) and \(X\) a semimartingale. The construction laid out in \parencite[Chapter.~4]{Protter2005StochasticIntegration} works as follows:\\
The integral for the simple processes is defined as in \cref{def:integral of simple processes}. One shows that the integral map 
\begin{equation}\label{eq:ucp continuity of integral map}
    I_X:\mathcal{S}_{ucp}\to\mathbb{D}_{ucp}
\end{equation}
is continuous with respect to the \(ucp\)-topology. Furthermore one shows that \(\mathcal{S}_{ucp}\) is dense (for the \(ucp\)-topology) in \(\mathbb{L}\). Under the additional insight that \(\mathbb{D}_{ucp}\) is a complete metric space the existence of a unique continuous extension 
\begin{equation}
    \tilde{I}_X:\mathbb{L}_{ucp}\to\mathbb{D}_{ucp}
\end{equation}
follows as usual from (....). For the proof of continuity of the map \cref{eq:ucp continuity of integral map} we refer to \parencite[Thm.~11]{Protter2005StochasticIntegration}. For the density of \(\mathcal{S}_{ucp}\) in \(\mathbb{L}\) we give some details as it forms an interesting contrast to the monotone class argument used in \cref{prop:density of E in L2(M) via monotone class}. 
\begin{proposition}\label{prop:ucp density of simple process in cadlag processes}
    The space \(\mathcal{S}\) is dense in \(\mathbb{L}\) for the \(ucp\)-topology
\end{proposition}
As often, before going into the details, we give some conceptual idea of the proof strategy. Given some \(Y\) We want to approximate it by elements of \(\mathcal{S}\) in the \(ucp\)-sense. After arguing that we can take \(Y\in b\mathbb{L}\) w.l.o.g. the idea is to approximate every path of \(Y\) by step functions that are uniformly \(\epsilon\)-close. As we want to do this for every path and we don't know in know how every path looks like as they are randome, we have to contruct this step functions by cutting up \(Y\) by stopping times in a way that the so obtained step functions are \(\epsilon\) close in every \((\omega,t)\). Notabene this yields an approximating sequence in \(\mathcal{S}\) and not in \(\mathcal{E}\) in general. It then remainse to check that the apprimating sequence in \(\mathcal{S}\) indeed approximates \(Y\) in \(ucp\). The main technical difficulties come from the fact that the random times we will define need not be stopping times for arbitary processes so we will have to choose apropriate version with the right path properties. 
\begin{proof}
    First notice that for \(Y\in\mathbb{L}\), we can stop \(Y\) with the stopping times \(F_n=\inf\)
\end{proof}
\textcolor{blue}{Notice that an monotone class argument as in \cref{prop:density of E in L2(M) via monotone class} would not be well adapted as it deeply relies on a class closed under bounded, monotone, pointwise convergence, which is ill conditioned to make statements about convergenc in \(ucp\). More precisely \(ucp\) convergence is to strong in the sense that it generally not implied by uniformly bounded monotone convergence, hence a class \(\mathcal{H}=\left\{ U \in \mathcal{A}:\forall \epsilon\exists H\in\mathcal{S}:\sup_{0\leq s<\infty}|H_t-U_t|<\epsilon \mathbb{P}-a.s. \forall t\in\mathbb{R}+\right\}\) is not closed under bounded montone limit, as this is a pointwise property and doens't converserve a uniform bound (not in probability or almost surely).}
\begin{theorem}\label{thm:ucp continuity of general stochastic integral}
    Let \(X\) be a semimartingale. Then the mapping 
    \begin{equation}
        I_X:\mathcal{S}_{ucp}\to\mathbb{D}_{ucp} 
    \end{equation}
    is continuous. 
\end{theorem}
This allows to define
\begin{definition}\label{def:stochastic integral w.r.t. semimartingales}
    Let \(X\) be a semimartingale. We define the stochastic integral
    \begin{align*}
        I_X:\;&\mathbb{L}\to\mathbb{D}\\
        &H\mapsto\int H\diff X
    \end{align*}
    as the unique continuous extension of the map in \cref{thm:ucp continuity of general stochastic integral} from \(\mathcal{S}\) to \(\overline{S}^{ucp}=\mathbb{L}_{ucp}\).  
\end{definition}
\subsection{Mean Variance Hedging in square integrable setting}\label{sec:Mean Variance Hedging in square integrable setting}
Finally we return to the main our main goal, the study of the mean-variance hedging problem outlined in \cref{eq:mean squared hedging}. We consider it in the square-integrable setting an show that in this setup it solution neetly reduces to an observation and interpretation of of the Kunita-Watanabe decomposition for the situation at hand. We restate the mean variance hedging problem for the square integrable setting
\begin{proposition}\label{prop:square integrable mean variance via KW}
    Let \(H\in L^2(\mathbb{P},\mathcal{F}_{\infty})\), \(X\in\mathcal{H}_{0,loc}^2\) and \(L^2(X)\) defined as in \cref{def:L^2(M)}. Furthermore assume \(\mathcal{F}_0\) to be trivial.
    Consider the following minimization problem
    \begin{equation}\label{eq:square integrable mean square hedging}
        \inf_{(\vartheta,V_0)\in L^2(X)\times\mathbb{R}}\lVert H-V_0-\int_{0}^{\infty}\vartheta_s\diff X_s\rVert_{L^2}
    \end{equation}
    Then the infimum is attained in \(L^2(X)\times\mathbb{R}\) and the optimal stratgey \((V_0^*,\vartheta^*)\) is given by \((\mathbb{E}[H|\mathcal{F}_0],\vartheta^H)\), where the latter are given by the Kunita Watanabe decomposition of the process \(V^H_t\coloneqq \mathbb{E}[H|\mathcal{F}_t]\) for \(t\in[0,\infty)\):
    \begin{equation}\label{eq:KW of V}
        V^H = \mathbb{E}[H]+\int\vartheta^H\diff X+L\;.
    \end{equation}
\end{proposition}
\begin{proof}
    To show existence we translate \cref{eq:square integrable mean square hedging} into a Hilbert projection problem. For this notice that the above can be written as a minimization problem for stochastic processes. To see this recall that we identified \(\mathcal{H}^2\) with \(L^2(\mathcal{F}_{\infty},\mathbb{P})\) via the map \(\phi:M\mapsto M_{\infty},\mathcal{H}^2\to L^2(\mathcal{F}_{\infty})\) with 
    \begin{equation}\label{eq:norm perserving property of L^2 identification}
        \lVert M\rVert_{\mathcal{H}^2}=\lVert M_{\infty}\rVert_{L^2}\;.
    \end{equation}
    Furthermore, setting
    \begin{equation}\label{eq:identifying r.v. H with process V}
        V^H_t\coloneqq \mathbb{E}[H|\mathcal{F}_t]\quad 0\leq t<\infty\;,
    \end{equation}
    which defines a unfiromly integrable martingale. 
    In particular \(H=V^H_{\infty}\)
    Under this identification we have due to \cref{eq:norm perserving property of L^2 identification}
    \begin{equation}\label{eq:equivalent formulation of mean varianc hedging for processes}
        \inf_{Z\in\mathbb{R}+ \mathcal{I}(M)}\lVert V-Z\rVert_{\mathcal{H}^2} =\inf_{Z\in\mathbb{R}+ \mathcal{I}(M)}\lVert H-Z_T\rVert_{L^2} 
    \end{equation}
    Now by \cref{lemma:stochastic integrals as stable subspace} \(\mathcal{I}(M)\) is a stable subspace of \(\mathcal{H}^2\), in particular closed. Does the same is true for \(\mathbb{R}+\mathcal{I}(M)\) as a subspace of \(\mathcal{H}^2\). Hence the right-hand side of \cref{eq:equivalent formulation of mean varianc hedging for processes} is a Hilbert projection problem of \(V\in\mathcal{H}^2\) onto the closed subsapce \(\mathbb{R}+\mathcal{I}(M)\). By (...) this shows existence. 
\\\\
It remains to show that the solution is given by the Kunita Watanabe decomposition. For this consider the Kunita-Watanabe decomposition of \(V\)
\begin{equation}\label{eq:KW of H}
    V = \mathbb{E}[F]+\int H_s\diff M_s + L\;,
\end{equation}
with \(H\in L^2(X)\) and \(L\in\mathcal{H}_0^2\) with \(L\perp X\). In particular, recall from the proof of \cref{thm:KW} that \(L\) is constructed to be in the orthogonal complement of \(\mathcal{I}(X)\) and in particular an element of \(\mathcal{H}_0^2\). We show that \(\mathbb{E}[H|\mathcal{F}_0]+\int \vartheta^H\diff X\) is the (weak) orthogonal projection of \(H\) onto the closed subspace \(\mathbb{R}+\mathcal{I}(M)\subset\mathcal{H}^2\)\footnote{\(\mathcal{I}(M)\subset \mathcal{H}_0^2\) is closed by \cref{...} and hence \(\mathbb{R}+\mathcal{I}(X)\) is also close in \(\mathcal{H}^2\)}. Indeed as \(L=V-\mathbb{E}[H|\mathcal{F}_0]+\int \vartheta^H\diff X\) we have
\begin{align*}
    &(L,Y_0+\int \vartheta^Z\diff X)_{\mathcal{H}^2}= (L,Y_0)_{\mathcal{H}^2}+ (L,\int \vartheta^Z\diff X)_{\mathcal{H}^2}=Y_0\lVert L_{\infty}\rVert_{L^2}=0\\&\forall Z_0+\int\vartheta^Z\diff X\in\mathbb{R}+\mathcal{I}(X)\;,
\end{align*}
because \(L\in\mathcal{H}_0^2\), in particular \(L_{\infty}\) exists and \(\mathbb{E}[L_{\infty}]=\mathbb{E}[L_0]=0\). Hence \cref{eq:KW of V} is the orthogonal decomposition of \(V^H\) onto \(\mathbb{R}+\mathcal{I}(X)\subset \mathcal{H}^2\) and the claim follows by pythagoras
\begin{equation}
    \lVert H-Z_0+\int\vartheta^Z\diff X\rVert_{\mathcal{H}^2}\geq\lVert H-\mathbb{E}[H]+\int \vartheta^H\diff X\rVert_{\mathcal{H}^2}
\end{equation}
for all \(Z_0+\int\vartheta^Z\diff X\in\mathbb{R}+\mathcal{I}(X)\). This shows that \(\mathbb{E}[H]+\int \vartheta^H\diff X\) is the minizer in \(\mathbb{R}+\mathcal{I}(X)\). Retranslating into the random variable form we get that \(\mathbb{E}[H|\mathcal{F}_0]+\int_0^{\infty} \vartheta^H_s\diff X_s\) minizies \cref{eq:square integrable mean square hedging} and hence the claim. 
\end{proof}
\section{General case}
These sections will be based primarily on \parencite{Schweizer2001GuidedTourQuadraticHedging} and the helpful, short overview in \parencite{RheinlaenderSexton2011HedgingDerivatives}. 
\\
Moreover, state the main results and underlying assumptions for the projection results:
\parencite[Thm.~46.]{Schweizer2001GuidedTourQuadraticHedging}. 

\chapter{Approximation by Algorithmic Strategies in Continuous Time}
Based on \parencite{Arandjelovic2025}.
\section{Outline of main ideas used in \parencite{Arandjelovic2025}}\label{sec:main ideas Alexandar}
\parencite[Chapter~5]{Arandjelovic2025} should be understood as an extension of the setup in \parencite{Buehler2018DeepHedging} in the sense that it leverages similar ideas developed in the latter, such as the Doob-Dynkin representation of measurable functions and the density of neural networks, but in the context of continuous time stochastic calculus. One has to deal with measurability and regularity issues not present in the finite setup; nevertheless, the core ideas remain similar. In this section, we guide the reader through this extension with the goal of arriving at a neural approximation theory for wealth processes of self-financing strategies and thus a continuous time version of \cref{prop:finite deep mean variance hedging}. Along the way, we highlight the conceptual proximity needed to the development of \parencite{Buehler2018DeepHedging} as described in \cref{ch:Deep Hedging in Discrete Time}.
\bigbreak
Upon closer inspection of the proof of \cref{prop:finite deep mean variance hedging}, we can infer that one of the crucial properties used is the fact that approximating the self-financed strategy \(\vartheta\) point-wise, for every time separately, also approximated the respective wealth processes. In the finite case, this is trivial, as the wealth processes were finite sums over the strategy at any point in time, i.e., \(\vartheta_k\) for \(k=0,\dots,n-1\). More generally, this is a statement about the continuity of the integral map
\begin{equation}\label{eq:stochastic integral map}
   I: \vartheta\mapsto \vartheta\bullet X\;,
\end{equation}
As we have seen in \cref{sec:Self financing strategies in continuous times} there are different settings and different such that we have indeed continuity but in different topologies. A very broad framework would be to approximate the strategies \(\vartheta\) in the sense that we find a subset \(D\) such that \(\overline{D}=\Theta\), where the closure is taken in the topology where \cref{eq:stochastic integral map} is continuous. If \(\Theta\) is first countable \footnote{In our setting these spaces will be one of the three: real-valued, adapted càglàd processes: \(\mathcal{L}_{ucp},\), real-valued, dapted càdlàg: \(\mathcal{D}_{ucp}\), \(L^p(X)\) for a square integrable martingale \(X\). All these space are first-countable. In the first two cases this follows from the fact that \(ucp\)-topology is metrizable. In the last case \(L^2(X)\) is even a Hilbert space.} this implies that for any such strategy \(\vartheta\in\Theta\) we can find an approximating sequence \((\vartheta_n)_{n\in\mathbb{N}}\) such that \(\lim_{n\to\infty}\vartheta_n=\vartheta\). By continuity of the integral \cref{eq:stochastic integral map} we have that also the integral \(\vartheta\bullet X\) is aproximated by \(\vartheta_n\bullet X\). Now this is not tell us how to construct the dense subset \(D\). In the context of deep hedging the general ideas should be to construct \(D\) by approximating parts of the strategies by neural networks, in an apropriate sense. In \cref{sec:Approximation of trading strategies} we did this by approximating the Doob-Dynkin representatiation for every time \(k=0,\dots,n\). We will quickly give motivation why a direct generalization of this argument to the continuous time setting is not the right way to appraoch this problem. 
\bigbreak
\textcolor{blue}{
To gain a deeper understanding of the approach taken in \parencite{Arandjelovic2025}, we will make some straightforward attempts on how one could think to generalize the discrete setup of \parencite{Buehler2018DeepHedging} to a time-continuous one. Understanding which ideas generalize well to the continuous case will also shed light on the approach taken in \parencite{Buehler2018DeepHedging}. First, let us take a look at the proof of \parencite[Prop.~4.9]{Buehler2018DeepHedging} as presented in \cref{thm:approximation by restricted hedging errors}. We will not be concentrating on the fact that the theorem is formulated for convex risk-measures, as we have deep mean-variance hedging in mind. Taking a discrete, predictable strategy \((\vartheta_k)_{k=0}^{n-1}\). The key ingredients to the proof are as follows:
\begin{itemize}
    \item Using the fact that \(\vartheta_k\) is \(\mathcal{F}_{k-1}=\sigma(Y_0,\dots,Y_{k-1})\)-measurable, we have by Doob-Dynkin that for \(k=0,\dots,n-1\) there exists a measurable map \(f_k\) such that \(\vartheta_k=f_k(Y_0,\dots,Y_{k-1})\).
    \item We have that \(f_k\in L^1(\mathbb{R}^{rk})\). Then for \(k=0,\dots,n-1\) we can find a sequence \((f_k^n)_{n\in\mathbb{N}}\subset\Psi\) that approximates \(f_k\) in \(L^1(\mathbb{R}^{rk})\). Then \(f_k(Y_0,\dots,Y_{k-1})\) approximates \(f(Y_0,\dots,Y_{k-1})\) in \(L^1(\mathbb{P})\). We define \(\vartheta^n=(f^n_k(Y_0,\dots,Y_{k-1}))_{k=0}^{n-1}\)
    \item For every \( k=0,\dots,n-1\) we can extract a subsequence such that the convergence holds \(\mathbb{P}\)-a.s.
    \item By a diagonalization argument, we can find a subsubsequence such that convergence holds \(\mathbb{P}\)-a.s \(\forall k\) 
    simultaneously.
    \item For the strategy \(\vartheta\) and the price process \(X\) we have that \((\vartheta^n\bullet X)_T\) converges to \((\vartheta\bullet X)_T\) \(\mathbb{P}\)-a.s.
\end{itemize}
As we are more interested in mean-variance hedging, we are content with \(L^2\)-convergence of \(((\vartheta^n-\vartheta)\bullet X)_T\). 
\\\\
\textbf{Example 1.}
The most natural and also naive approach would be to consider a filtration \(\mathbb{F}=(\mathcal{F}_t)_{t\in\mathbb{R}_+}\) where \(\mathcal{F}_t=\sigma(Y_s:0\leq s\leq T)\), on a complete probability space \((\Omega,\mathcal{F},\mathbb{P})\), in which we take the market information process \(Y\) RCLL. Even more, assume that it takes values in a locally convex path space \(D([0,T],\mathbb{R}),\), which we equip with its cylinder \(\sigma\)-algebra. Consider a \(\mathbb{F}\)-predictable strategy \(\vartheta\). In particular, \(\vartheta\) is \(\mathbb{F}\)-adapted, thus, for every \(t\in\mathbb{R}_+\) we have by Doob-Dynkin that there exists some (Borel-) measurable \(f^t:\mathbb{R}\to E\):
\begin{equation}
    \vartheta_t=f_t(Y^t)\;,
\end{equation}
where we define the stopped process \(Y^t=(Y_{s\wedge t})_{s\in\mathbb{R}_+}\). Here we implicitly used the fact that \(\sigma(Y_s:0\leq s\leq t)=\sigma(Y^t)\). Notice, crucially, the dependence in \(t\) on the map \(f_t\). Now, assuming \(f^t\in L^1(D([0,T],\mathbb{R}),\mathbb{P}\circ(Y^t)^{-1})\), we can approximate \(f_t\) by a sequence of neural networks \((f_t^n)_{n\in\mathbb{N}}\subset\mathcal{NN}(D^*,\psi)\) in \(L^1\) (...). Now considering \(\vartheta^n_t\coloneqq f^n_t(Y^t)\).
\\\\
Now, there is no reason why \(\vartheta^n_t\) should be predictable, which a claim about the measurability of the map 
\begin{equation}
    (t,\omega)\mapsto f_t(Y^{t}(\omega))\;,
\end{equation}
where \(Y^{t}(\omega)\) denotes the stopped path of \(Y(\omega)\) on \([0,T]\).
\\\\
\textbf{Example 2.}
Nevertheless, we see from the above approach that the problem arises due to the pointwise application of the Doob-Dynkin lemma, where we lose predictability of a joint measurability condition. This motivates us to apply Doob-Dynkin on a process level. That is, we consider the equivalent description of a stochastic process.
\begin{align}
    Y&:\Omega\times[0,T]\to\mathbb{R}\\
    &(\omega,t)\mapsto Y_t(\omega)\;.
\end{align}
Then \(\sigma(Y)=\left\{(\omega,t):Y_t(\omega)\in A,\; A\in\mathcal{B}(\mathbb{R})\right\}\) is a \(\sigma\)-field on \(\Omega\times[0,T]\). Now consider a strategy predictable \(\vartheta\) and assume that \(\vartheta\) is \(\sigma(Y)\)-measurable. We can now apply Doob-Dynkin for the whole process \(\vartheta\). We get that there exists a measurable function \(f:\mathbb{R}\to\mathbb{R}\) such that
\begin{equation}
    \vartheta=f(Y)\;.
\end{equation}
 in order to obtain a mean-variance result, we want to show that we can approximate the \(L^2\)-norm of \(\int_0^T \vartheta\diff X\) by \(\int_0^T \vartheta^n\diff X\) for some sequence \(\vartheta^n\). To make our point, we take the setup where the asset processes \(X\in \mathcal{M}_{0,loc}^2\) and \(\vartheta\in L^2(X)=L^2(\Omega,\mathcal{P},\mathbb{P}_X)\), with \(\mathbb{P}_X=\mathbb{P}\otimes\diff [X]\). Denoting 
 \begin{equation}
     \mathbb{P}_Y=\mathbb{P}\circ Y^{-1}
 \end{equation}
 the push-forward measure on \((\mathbb{R},\mathcal{B}(\mathbb{R}))\) directly implies that the condition \(\vartheta\in L^2(X)\) implies that \(f\in L^2(\mathbb{R},\mathbb{P}_Y)\) (iff). We can find a sequence of neural networks  \(f^n\) that approximate \(f\) in a \(L^2(\mathbb{R},\mathbb{P}_Y)\)-sense. As usual, we define
 \begin{equation}
     \vartheta^n= f^n(Y)
 \end{equation}
 Again, we find the same problem: \(\vartheta\) is not predictable in general. So again, we cannot even start to consider the convergence of stochastic integrals, as they are not defined. We need a completely different approach. 
 }
 In light of the issues found above and the importance of continuity of the map \cref{eq:stochastic integral map} we will start with simple bounded predictable integrands. For this recall form \cref{def:simple bounded predictable processes} that a simple bounded predictable process \(\vartheta\) has the representation
 \begin{equation}
     \vartheta_t = \vartheta^01_{\left\{ 0\right\}}+\sum_{k=1}^n \vartheta^k1_{]\tau_k,\tau_{k+1}]}
 \end{equation}
 for a non-decreasing family of stopping times \(0<\tau_1,\dots<\tau_n<\infty\) and \(\vartheta^i\in\mathcal{L}^{\infty}(\mathcal{F}_{t_k},\mathbb{P})\). The crucial idea is now to apply Doob-Dynkin to the random variables \(\vartheta^i\) rather than directly to \(\vartheta_t\). The approach taken in \parencite{Arandjelovic2025} will be presented here. Approximating simple processes will preserve the continuity of \cref{eq:stochastic integral map}, in the appropriate topology, by the very construction of the respective stochastic integrals. Furthermore, predictability will also be preserved. Indeed, for some element of some approximating sequence \((\vartheta^k_n)_{n\in\mathbb{N}}\) in \(\mathcal{L}^{\infty}(\mathcal{F}_{t_k},\mathbb{P})\), we have that \(\vartheta_n^01_{\left\{ 0\right\}}+\sum_{k=1}^n \vartheta_n^k1_{]\tau_k,\tau_{k+1}]}\) is still a simple bound predictable process. We have seen in \cref{sec:Self financing strategies in continuous times} that both cases, \(\tau_i\) general \(\mathbb{F}\)-stopping times and \(\tau_i=t_i\in\mathbb{R}_+\) deterministic times are interesting for stochastic integration. We will does discuss the developpment of the right approximation system for \(\mathcal{L}^{\infty}(\mathcal{F}_{t_k},\mathbb{P})\) in both cases. 
 \bigbreak
 \subsection{Algorithmically Generated Random Variables}\label{sec:algorithmically generated random variables}
 In all that follows, assume \((\Omega,\mathcal{F},\mathbb{P})\) to be a complete probability space. This will be a standing assumption.
 \begin{proposition}\label{prop:reduction to countable measurability}
     Let \((\Omega,\mathcal{F})\) and \(S,\mathcal{S}\) be measurable spaces. Let \(I\) be a non-empty set and \((F_i)_{i\in I}\) be an indexed family of sub-\(\sigma\)-algebras of \(\mathcal{F}\). For any subset \(J\subset I\), we define \(\mathcal{F}_J=\sigma(\bigcup_{j\in J}\mathcal{F}_j)\). Let \(X:\Omega\to S\) be and \(\mathcal{F}_I/\mathcal
     S\)-measurable map such that the restricted \(\sigma\)-algebra on \(X(\Omega)\); \(\mathcal{S}_X\coloneqq\left\{B\cap g(\Omega):B\in \mathcal{S}\right\}\), is \textbf{countably} generated. Then there exists a \textbf{countable} subset \(J\subset I\) such that \(g\) is \(F_{J}/\mathcal{S}\)-measurable. 
 \end{proposition}
 \begin{remark}\label{Remark:countably generated sigma algebra}
      Notice that the claim is not only that we can reduce measurability to a countably generated \(\sigma\)-algebra, but that for \(F_I =\sigma(\bigcup_{i\in I}\mathcal{F}_i)\), it has the form \(\mathcal{F}_J=\sigma(\bigcup_{j\in J}\mathcal{F}_j)\) for \(J\subset I\) countable. Indeed, the mere fact that if \(\mathcal{S}_X\) is countably generated, one can find a sub-\(\sigma\)-algebra \(\mathcal{F}_X\subset\mathcal{F}_I\) that is countably generated and is generated by a countable set of sets. Indeed, as \(\mathcal{S}_X\) is countably generated, there exists a sequence \((S_n)_{n\in\mathbb{N}}\) in \(\mathcal{S}\) such that \(\mathcal{S}_X=\sigma(S_n\cap X(\Omega):n\in\mathbb{N})\). Setting 
 \(\mathcal{F}_X=\sigma(X^{-1}(S_n):n\in\mathbb{N})\), we have that \(\mathcal{F}_X\subset\mathcal{F}_I\) and that \(X\) is \(\mathcal{F}_X\)-measurable. Indeed, this is part of the proof of \cref{prop:reduction to countable measurability}. 
 \end{remark}
 How to get the further property of is due to the following general property, which we formulate as a lemma
 \begin{lemma}
     Let \(I\) be a non-empty set and \(\mathcal{F}_J=\sigma(\bigcup_{j\in J}\mathcal{F}_j)\) for any \(J\subset I\). Then 
     \begin{equation}
         \mathcal{F}_I=\bigcup_{\substack{J\subset I\\ J\text{ countable}}}\mathcal{F}_J
     \end{equation}
 \end{lemma}
\begin{proof}
 Clearly, part of the claim is that the right-hand side is a \(\sigma\)-algebra. For notational convenience in this proof we denote it by \(\mathcal{G}\). We will check this first before proving equality.
 Indeed, \(\Omega\in \mathcal{F}_J\) for every countable \(J\subset I\); in particular, \(\Omega\in \mathcal{G}\). If \(A\in\mathcal{G}\), then there exists a \(J\subset I\) countable such that \(A\in\mathcal{F}_J\). Hence, \(A^c\in\mathcal{F}_J\); thus, \(A^c\in\mathcal{G}\). Consider a sequence \((A_n)_{n\in\mathbb{N}}\) in \(\mathcal{G}\). 
 Then \(\forall n\in\mathbb{N}\); there is a \(J_n\subset I\) countable such that \(A_n\in \mathcal{F}_{J_n}\). Observe that by the definition of \(\mathcal{F}_J\), for a \(J_1\subset J_2\subset I\) countable, we have \(\mathcal{F}_{J_1}\subset\mathcal{F}_{J_2}\). Thus, setting \(J=\bigcup_{n\in\mathbb{N}}J_n\), which is countable as a countable union of countable sets, we have \(\bigcup_{n\in\mathbb{N}}A_n\in\mathcal{F}_J\). Hence, \(\bigcup_{n\in\mathbb{N}}A_n\in\mathcal{F}_{J}\). Hence, \(\mathcal{G}\) is a \(\sigma\)-algebra. 
 \\\\
 To show equality, it remains to show that 
 \begin{equation}
     \mathcal{F}_I\subset \bigcup_{\substack{J\subset I\\ J\text{ countable}}}\mathcal{F}_J\;.
 \end{equation}
 As commented 
 Indeed, by construction, \(\bigcup_{i\in I}\mathcal{F}_i\) generates \(\mathcal{F}_I\); hence, it is enough to show that this generating set is in \(\mathcal{G}\). For \(B\in\bigcup_{i\in I}\mathcal{F}_i\), there exists an \(i\in I\) such that \(B\in \mathcal{F}_i\subset\mathcal{G}\), as the singleton \(\left\{i\right\}\) is a countable subset of \(I\). This concludes the proof. 
\end{proof}
 
 Before giving the proof, recall that the Borel-\(\sigma\)-algebra of a separable metric space is countably generated; hence, the above can be applied to \((S,\mathcal{B}(S))\), for a separable metric space \(S\). Indeed, recall that the open balls induce a topology on \(S\), in the sense that they form a basis. If there exists a countable subset \(D\) dense in \(S\), then the family \(\left\{B(x,r):r\in D,x\in \mathbb{Q}_+\right\}\) also forms a basis for the same topology; thus, the respective Borel-\(\sigma\)-algebra is countably generated. Furthermore, we will be especially interested in the case where \(I=\mathbb{R}_+\), \(\mathcal{F}_t=\sigma(Y_s:0\leq s\leq t)\) for some stochastic process \(Y\). Then consistently, for \(J=[0,t]\subset\mathbb{R}_+\) we have \(\mathcal{F}_t=\sigma(\bigcup_{s\in[0,t]}\mathcal{F}_s)=\sigma(Y_s:0\leq s\leq t)\). Taking \((S,\mathcal{S})\) to be \((\mathbb{R},\mathcal{B}(\mathbb{R}))\) and denoting \(\mathcal{F}_{\infty}
 \coloneqq \sigma(\bigcup_{t\in\mathbb{R}_+}\mathcal{F}_t)\),  \cref{prop:reduction to countable measurability} implies, in particular, that for the real-valued random variable X, which is \(\mathcal{F}_{\infty}/\mathcal{B}(\mathbb{R})\)-measurable, there exists a countable subset \(\left\{t_1,t_2,\dots\right\}\subset\mathbb{R}_+\) such that \(X\) is \(\sigma(Y_{t_1},Y_{t_2},\dots)/\mathcal{B}(\mathbb{R})\)-measurable.
 \bigbreak
 Eventually, we give the proof of \cref{prop:reduction to countable measurability}
 \begin{proof}
 The case where \(I\) is countable is trivial, so we assume \(I\) to be uncountable
 The first part follows the \cref{Remark:countably generated sigma algebra}. We call it here for completeness. By assumption, \(\mathcal{S}_X\) is countably generated; hence, there exists a sequence \((S_n)_{n\in\mathbb{N}}\) in \(\mathcal{S}\) such that \(\mathcal{S}_X=\sigma(S_n\cap X(\Omega):n\in\mathbb{N})\). 
 Clearly, \(X\) is \(\mathcal{F}_X\)-measurable. Now we conclude by \cref{lemma:Finiteness of inf for finite sample space}. We have \(\mathcal{F}_X\subset\mathcal{F}_I\), or equivalently
 \begin{equation}
     \sigma(X^{-1}(S_n):S_n\in\mathcal{S})\subset\bigcup_{\substack{J\subset I\\ J\text{ countable}}}\mathcal{F}_J\;.
 \end{equation}
 In particular, for every \(n\in\mathbb{N}\), there exists a countable set \(J_n\subset I\) such that
 \(X^{-1}(S_n)\in \mathcal{F}_{J_n}\). Hence, setting \(J=\bigcup_{n\in\mathbb{N}}J_n\) we have 
 \begin{equation}
     X^{-1}(S_n)\in\mathcal{F}_J\quad\forall n\in\mathbb{N}\;,
 \end{equation}
 where \(J\), a countable union of countable sets, is countable.
 But \(\mathcal{F}_X=\sigma(X^{-1}(S_n):S_n\in\mathcal{S})\) is the smallest \(\sigma\)-algebra for which \(X^{-1}(S_n)\), for all \(n\in\mathbb{N}\), are measurable, thus \(\mathcal{F}_X\subset \mathcal{F}_J\). We conclude by recalling again that \(X\) is \(\mathcal{F}_X\) and thus also \(\mathcal{F}_J\)-measurable.
   \end{proof}
As mentioned, we will be especially interested in the case where \(I=\mathbb{R}_+\) and where the collection of \(\sigma\)-algebras \((\mathcal{F}_i)_{i\in I }\) forms an increasing sequence \(\mathcal{F}_t=\{Y_s:0\leq s\leq t\}\) for some stochastic process \(Y\). Then \(\mathbb{F}=\{\mathcal{F}_t\}_{t\in\mathbb{R}_+}\) defines a filtration on \((\Omega,\mathcal{F},\mathbb{P})\). It is usual and advantageous to require that the filtration satisfies the usual assumption; that is, that it is right-continuous and complete. The completion can be done canonically by defining \(\mathcal{\overline{F}}_t=\sigma(\mathcal{F}_t\cup\mathcal{N}_{\mathbb{P}})\), where \(\mathcal{N}_{\mathbb{P}}\) is as in \cref{eq:nullsets}. We denote \(\overline{\mathbb{F}}=(\overline{\mathcal{F_t}})_{t\in\mathbb{R}_+}\)\textcolor{blue}{As for right continuity....}. \\
As motivated in the introduction of the current chapter, our goal is to develop an approximation theory for \(\mathcal{F_{\tau}}\)-measurable, bounded random variables, where \(\tau\) is an \(\overline{\mathbb{F}}\)-stopping time (truly random or deterministic) and \(\overline{\mathbb{F}}\) is the completed filtration generated by a Lévy process \(Y\), hence satisfying the usual assumptions. 
Thanks to \cref{prop:reduction to countable measurability}, we now have the machinery to do so. The following corollary is essentially a special case of \cref{prop:reduction to countable measurability}. 
\begin{proposition}\label{prop:reduction to countable measurability with filtration}
    In the setup of \cref{prop:reduction to countable measurability} with \(Y\) a stochastic process and filtrations \(\mathbb{F}=(\mathcal{F}_t)_{t\in\mathbb{R}_+}\) and \(\overline{\mathbb{F}}=(\overline{\mathcal{F}_t})_{t\in\mathbb{R}_+}\) where \(\mathcal{F}_t=\sigma(Y_s:0\leq s\leq t)\) and \(\overline{\mathcal{F}}_t=\sigma(Y_s:0\leq s\leq t)\vee\mathcal{N}_{\mathbb{P}}\) respectively. Then
    \begin{enumerate}[label=(\alph*)]
        \item For fixed \(t\in\mathbb{R}_+\), let \(X:\Omega\to\mathbb{R}\) be \(\mathcal{F}_t\)-measurable. Then, there exists a \textit{countable} set \(J\subset [0,t]\) such that \(X\) is \(\mathcal{F}_J\)-measurable, where \(\mathcal{F}_J = \sigma(Y_t:t\in J)\). 
        \item Assume Y is a Lévy process. Let \(\tau\) be a \(\overline{\mathbb{F}}\)-stopping time. Furthermore, let \(X:\Omega\to\mathbb{R}\) be \(\mathcal{F}_{\tau}\)-measurable. Then, there exists a countable set \(J\subset\mathbb{R}_+\) such that \(X\) is \(\mathcal{F}_J\)-measurable, where \(\mathcal{F}_J=\sigma(\tau)\vee\sigma(Y_t^{\tau}:t\in J)\vee \mathcal{N}_{\mathbb{P}}\).
    \end{enumerate}
\end{proposition}
\begin{proof}
    The proof is an application of \cref{prop:reduction to countable measurability} with different choices of \(I\) and \(\mathcal{F}_i\). For (a) take \(I=\mathbb{R}_+\) and \(\mathcal{F}_i=\sigma(Y_i)\) for \(i\in I\). For (b) by \parencite[Thm.~13.45]{HeWangYan1992} (recalling that \(\overline{\mathcal{F}}_t=\mathcal{F}_t\vee \mathcal{N}_{\mathbb{P}}\)) we have \(\mathcal{F}_{\tau}=\sigma(\tau)\vee\sigma(Y_t^{\tau}:t\in J)\vee \mathcal{N}_{\mathbb{P}}\) take \(I=\mathbb{R}_+\cup\left\{a,b\right\}\) and \(\mathcal{F}_i=\sigma(Y_i^{\tau})\) for \(i\in \mathbb{R}_+\) and \(\mathcal{F}_a=\sigma(\tau)\) and \(\mathcal{F}_b=\sigma(\mathcal{N}_{\mathbb{P}})\). 
\end{proof}
\cref{prop:reduction to countable measurability with filtration} now motivates the following definition, which will be the building block for our approximation theory of simple bounded predictable processes based on neural networks. For the following, let \(\mathcal{F}_{\infty}\coloneqq\bigvee_{t\in\mathbb{R}_+}\mathcal{F}_t\) and \(\overline{\mathcal{F}_{\infty}}=\mathcal{F}_{\infty}\vee \mathcal{N}_{\mathbb{P}}\).
\begin{definition}\label{def:algorithmic strategies}
    \begin{itemize}
        \item We call an \(\overline{\mathcal{F}_{\infty}}\)-measurable, real-valued random variable \(\vartheta\) \textit{neural} if there exists a subset \(I\) of \([0,\infty)\)\footnote{\(I\) must not necessarily be a proper subset. We include the case \(I=\mathbb{R}_+\)} and a neural network \(f\in\mathcal{NN}(\mathbb{R}^I,\psi)\), such that 
        \begin{equation}
            \vartheta = f(Y_i:i\in I)\;.
        \end{equation}
        \item For \(t\in[0,\infty]\), we call an \(\overline{\mathcal{F}_t}\)-measurable, real-valued random variable \(\vartheta\) \textit{algorithmically generated} if there exist deterministic times 
        \(t_1,\dots,t_n\) in \([0,t]\) for \(n\in\mathbb{N}\) and a neural network \(f\in\mathcal{NN}((\mathbb{R}^{dn})^*,\psi)\), such that
        \begin{equation}
            \vartheta = f(Y_{t_1},\dots,Y_{t_n})\;.
        \end{equation}
        We denote the class of these processes by \(\Theta_{\mathcal{NN}}(t,\psi)\). 
        \item
        For \(Y\), a Lévy process, and \(\tau\) an \(\overline{\mathbb{F}}\) stopping time, we call a \(\mathcal{F}_{\tau}\)-measurable, real-valued random variable \(\vartheta\), \textit{algorithmically generated}, if there exist deterministic times 
        \(t_1,\dots,t_n\) in \([0,t]\) for \(n\in\mathbb{N}\) and a neural network \(f\in\mathcal{NN}((\mathbb{R}^{(1+d)n})^*,\psi)\), such that
        \begin{equation}
            \vartheta = f(\tau,Y^{\tau}_{t_1},\dots,Y^{\tau}_{t_n})\;.
        \end{equation}
        We denote the class of these processes by \(\Theta_{\mathcal{NN}}(\tau,\psi)\). 
    \end{itemize}
\end{definition}
We now turn to the main approximation theoretic result, upon which the rest will be based. We will show that the algorithmically generated random variables are dense in \(\mathcal{L}^p(\Omega,\overline{\mathcal{F}},\mathbb{P})\) for \(p\in[0,\infty)\), where \(\overline{\mathcal{F}}=\sigma(Y_s:0\leq s\leq t)\). As expected from the finite setting \parencite{Buehler2018DeepHedging} the proof will consist of applying the Doob-Dynkin representation theorem. The application of the latter will be standard, as we have previously argued how to reduce measurability with respect to countably many random variables \cref{prop:reduction to countable measurability with filtration}. This is not fundamental from a theoretical point of view, as we shall argue; nevertheless, it makes for a clearer connection with \parencite{Buehler2018DeepHedging} and allows for a more general setup.
\\\\
For the proof of the following proposition, we will need the well known martingale convergence theorem, which we restate here for the convenience of the reader. The proof for a more general version on Orlicz spaces can be found in \parencite[Thm.~5.12]{Arandjelovic2025}. We state the special case for the Orlicz space \(L^p(\mu)\) where \(\mu\) is a finite measure on a measure space \((\Omega,\mathcal{F})\) and \(p\in[1,\infty)\).
\begin{lemma}[Martingale Convergence Theorem]\label{lemma:Martingale Convergence Theorem}
    Let \((\mathcal{F}_n)_{n\in\mathbb{N}}\) be a filtration of \(\mathcal{F}\) and \(\mathcal{F}_{\infty}\coloneqq\bigvee_{n\in\mathbb{N}}\mathcal{F}_n\). Then for every \(H\in L^p(\mu)\) for some finite measure \(\mu\) and \(p\in[0,\infty)\), we have
    \begin{equation}
        \mathbb{E}_{\mu}[H|\mathcal{F}_n]\to\mathbb{E}_{\mu}[H|\mathcal{F}_{\infty}],\quad n\to\infty
    \end{equation}
 where convergence holds \(\mu\)-almost surely and in \(L^p(\mu)\).
\end{lemma}
The claim is generally wrong if \(\mu\) is not a finite measure.
\begin{proposition}\label{prop:approximation of measurable maps by algorithmic strategies}
\begin{enumerate}
    \item Let \(I\) be a non-empty set, and \((Z_i)_{i\in I}\) an indexed family of \(\mathcal{F}\)-measurable, \(\mathbb{R}^d\)-valued random variables. Let 
    \begin{equation}
        \overline{\mathcal{F}_I}=\mathcal{N}^\mathbb{P}\vee\sigma(Z_i:i\in I)
    \end{equation}
    and \(X\in\mathcal{L}^p(\mathbb{P})\) for \(p\in[1,\infty)\), \(\overline{\mathcal{F}_I}\)-measurable. Then for every \(\epsilon>0\), there exists \(n\in\mathbb{N}\) and a \textbf{finite} set \(J=\left\{j_1,\dots,j_n\right\}\subset I\), and a neural network \(f\in\mathcal{NN}(\mathbb{R}^{dn})^*,\psi\), such that 
    \begin{equation}
        \lVert X-f(Z_{j_1},\dots, Z_{j_n})\rVert_{p}<\epsilon\;.
    \end{equation}
    \item Let \(t\in[0,\infty]\) and \(X\in\mathcal{L}^p(\mathbb{P})\) for \(p\in[1,\infty)\), \(\overline{\mathcal{F}_I}_t\)-measurable. Then, for every \(\epsilon>0\), there exists \(\vartheta\in\Theta_{\mathcal{NN}}(t,\psi)\) such that 
    \begin{equation}
        \lVert \vartheta-X\rVert_{p}<\epsilon\;.
    \end{equation}
    \item Assume \(Y\) is a Lévy process. Let \(\tau\) be an \(\overline{\mathbb{F}}\)-stopping time, and \(X\in\mathcal{L}^p(\mathbb{P})\) for \(p\in[1,\infty)\), \(\overline{\mathcal{F}_I}_{\tau}\)-measurable. Then, for every \(\epsilon>0\), there exists \(\vartheta\in\Theta_{\mathcal{NN}}(\tau,\psi)\) such that 
    \begin{equation}
        \lVert \vartheta-X\rVert_{p}<\epsilon\;.
    \end{equation}
\end{enumerate}
\end{proposition}
\begin{proof}   
The main work is to prove part (1). Part (2) and (3) will then follow from the same arguments applied to special cases. \\
The main idea for part (1) is to rewrite the \(\overline{\mathcal{F}_I}\)-measurable function \(X\) as a conditional expectation and find a filtration such that \(X=\mathbb{E}[X|\overline{\mathcal{F}_I}]\) is the \(L^p\) limit of \(X_n=\mathbb{E}[X|\overline{\mathcal{F}}_n]\). We will use \cref{prop:reduction to countable measurability} to construct the filtration. Then we shall apply Doob-Dynkin to \(X_n\), and approximate the latter by neural networks. Having laid out the proof stratgey we give the details:\\
First notice that for \(\overline{\mathcal{F}_I}=\mathcal{F}_I\vee\mathcal{N}_{\mathbb{P}}\), we have
\begin{align}
    X=\mathbb{E}[X|\overline{\mathcal{F}_I}]\overset{\mathbb{P}-a.s.}{=}\mathbb{E}[X|\mathcal{F}_I]. 
\end{align}
Indeed, by definition, the conditional expectation is the \(\mathbb{P}\)-almost surely unique, integrable random variable \(\mathbb{E}[X|\overline{\mathcal{F}_I}]=\) such that \(\mathbb{E}[X1_B]=\mathbb{E}[\mathbb{E}[X|\overline{\mathcal{F}_I}]1_B]\) for all \(B\in\overline{\mathcal{F}_I}\). By \cref{lemma:characterisation of completed sigma-algebra} we have \(B=F\cup E\) for \(F\in\mathcal{F}_I\) and \(E\in\mathcal{N}_{\mathbb{P}}\), and we can choose w.l.o.g. \(F\cap E=\emptyset\); hence \(1_{B}=1_E+1_F\). However, \(\mathbb{P}(E)=0\)\footnote{ or,More precisely, the completed measure \(\overline{\mathbb{P}}\) defined by \(\overline{\mathbb{P}}(E\cup F)=\mathbb{P}(E)\) for all \(E\cup F\in\overline{\mathcal{F}_I}\).} Hence, \(\mathbb{E}[X1_E]=\mathbb{E}[\mathbb{E}[X|\overline{\mathcal{F}_I}]1_E]\) for all \(E\in\mathcal{F}_I\) and \(\mathbb{E}[X|\overline{\mathcal{F}_I}]\) is integrable by definition; hence \(\mathbb{E}[X|\overline{\mathcal{F}_I}]=\mathbb{E}[X|\mathcal{F}_I]\). For the scope of this proof we denote \(X_{\mathcal{F}_I}=\mathbb{E}[X|\mathcal{F}_I]\) and \(X_{\overline{\mathcal{F}_I}}=\mathbb{E}[X|\overline{\mathcal{F}_I}]\).
\\\\
Recall that \(\mathcal{F}_I=\sigma(Z_i:i\in I)\) and \(\mathbb{E}[X|\mathcal{F}_I]\) is \(\mathcal{F}_I\)-measurable by definition. 
Now by \cref{prop:reduction to countable measurability} we can find \(J= \left\{j_n:n\in\mathbb{N}\right\}\subset I\) such that \(\mathbb{E}[X|\mathcal{F}_I]\) is \(\mathcal{F}_J=\sigma(Z_j:j\in J)\)-measurable. We can now define a filtration on \(\mathcal{F}_J\) by setting \(\mathcal{F}_n=\sigma(Z_{j_1},\dots,Z_{j_n})\) and \(\mathbb{F}=(\mathcal{F}_{n})_{n\in\mathbb{N}}\). We even have \(\mathcal{F}_J=\bigvee_{n\in\mathbb{N}}\mathcal{F}_n\). We set for \(n\in\mathbb{N}\):
\begin{equation}
    M_n=\mathbb{E}[X_{\mathcal{F}_I}|\mathcal{F}_n]\;.
\end{equation}
and denote the resulting \(\mathbb{F}\)-adapted process by \(M=(M_n)_{n\in\mathbb{N}}\). By \cref{lemma:Martingale Convergence Theorem} we have that 
\begin{equation}
    M_n\overset{L^p}{\to}\mathbb{E}[X_{\mathcal{F}_I}|\mathcal{F}_J]=X_{\mathcal{F}_I}\quad n\to\infty
\end{equation}
In particular, because \(X_{\mathcal{F}_I}\overset{\mathbb{P}-a.s.}{=}X_{\overline{\mathcal{F}_I}}\), we have \(\lim_{n\to\infty}M_n=X_{\overline{\mathcal{F}_I}}\) in \(L^p\). Now, we are in a position to apply the Doob-Dynkin lemma to \(M_n\), which yields the existence of a \(\mathcal{B}(\mathbb{R}^n)/\mathcal{B}(\mathbb{R})\)-measurable map \(f_n\) such that 
\begin{equation}
    M_n = f_n(Z_{j_1},\dots,Z_{j_n})\;.
\end{equation}
Because of the convergence of \(M_n\) to \(X_{\mathcal{F}_I}\) in \(L^p\), we have that \(f_n\in L^p(\mu_n)\), where \(\mu_n=\mathcal{L}(Y_{j_1},\dots,Y_{j_n})=\mathbb{P}\circ(Y_{j_1},\dots,Y_{j_n})^{-1}\). By \cref{.....} we can approximate \(f_n\) by a sequence \((f_n^m)_{m\in\mathbb{N}}\subset \mathcal{NN}({\mathbb{R}^n,\psi})\) in \(L^p(\mu_n)\). \\\\
Fix some \(\epsilon>0\). By convergence of \(M_n\) to \(X_{\mathcal{F}_I}\) in \(L^p(\mathbb{P})\) there exists \(n\in\mathbb{N}\) large enough such that \(\lVert M_n-X_{\mathcal{F}_I}\rVert_{L^p(\mathbb{P})}<\epsilon/2\). By convergence of \(f_n^m(Z_{j_1},\dots,Z_{j_n})\) to \(M_n\) in \(L^p(\mu_n)\), for every \(n\in\mathbb{N}\) there exists \(m\in\mathbb{N}\) large enough such that \(\lVert f_n^m-M_n\rVert_{L^p(\mu_n)}<\epsilon/2\). Hence, we have by the triangle inequality
\begin{equation}
    \lVert M_n-X_{\overline{\mathcal{F}_I}}\rVert_{L^p(\mathbb{P})}=\lVert M_n-X_{\mathcal{F}_I}\rVert_{L^p(\mathbb{P})}\leq \lVert M_n-X_{\mathcal{F}_I}\rVert_{L^p(\mathbb{P})}+\lVert f_n^m-M_n\rVert_{L^p(\mu_n)}<\epsilon\,,
\end{equation}
which finishes the proof of the first part. \\\\
For part (2), the claim follows by the same argument as in part (1) with \(I=[0,t]\) for \(t\in\mathbb{R}_+\) as in \cref{prop:reduction to countable measurability with filtration} part (a). For part (3), again by the same argument, with \(I=[0,t]\cup\{\tau\}\) for \(t\in\mathbb{R}_+\) and \cref{prop:reduction to countable measurability with filtration} part (b). 
\end{proof}
Let us make an interesting observation, phrased as an proposition, which should be read as a comment to \cref{prop:approximation of measurable maps by algorithmic strategies} part (2). 
\begin{proposition}
    Let the market information process \(Y:(\Omega,\mathcal{F})\to(E,\mathcal{E}),\; \omega\mapsto \left(t\mapsto Y_t(\omega)\right)\) take values in a locally-convex path space \(E\) and take \(\mathcal{E}\) to be the cylinder \(\sigma\)-algebra on \(E\). Let \(I\subset\mathbb{R}_+\) and \(X\) be \(\overline{\mathcal{F}_I}\)-measurable and in \(\mathcal{L}^p(\mathbb{P})\). Then for every \(\epsilon>0\) there exists \(F\in\mathcal{NN}(E^*,\psi)\) such that 
    \begin{equation}
        \lVert X-F(Z)\rVert_{L^p(\mathbb{P})}<\epsilon\;.
    \end{equation}
\end{proposition}
\begin{proof}
    Again set \(X=\mathbb{E}[X|\overline{\mathcal{F}_I}]\overset{\mathbb{P}-a.s.}{=}\mathbb{E}[X|\mathcal{F}_I]\) and set \(X_{\mathcal{F}_I}=\mathbb{E}[X|\mathcal{F}_I]\). In our setting \(\mathcal{F}_I=\mathcal{F}_{\infty}=\sigma(Y_t:0\leq t<\infty)=\sigma(Y)\). Now \(X\) is, as usual, a real-valued random variable which is \(\sigma(Y)\) measurable, where \(Y\), as a stochastic process, is a path-valued random variable. Hence by Doob-Dynkin there exists a \(\mathcal{E}/\mathcal{B}(\mathbb{R})\)-measurable \(F\) such that 
    \begin{equation}
        X = F(Y)\;.
    \end{equation}
    Because \(X\in\mathcal{L}^p(\mathbb{P})\) we have \(F\in\mathcal{L}^p(\mu)\), where \(\mu=\mathcal{L}(Y)=\mathbb{P}\circ(Y)^{-1}\) the law of \(Y\) on the path space \((E,\mathcal{E})\). Now by the generalized density result \cref{....} we have the existence of an \(L^p(\mu)\)-approximating sequence \((F^n)_{n\in\mathbb{N}}\) of \(F\). Hence for every \(\epsilon>0\) there exists \(n\in\mathbb{N}\) large enough such that 
    \begin{equation}
        \lVert f(Z)-F_n(Z)\rVert_{L^p(\mathbb{P})}=\lVert f-F_n\rVert_{L^p(\mu)}<\epsilon\;.
    \end{equation}
    Setting \(F=F_n\) concludes the proof.
\end{proof}
One might take \(E=C([0,\infty))\) with topology induced by the countable family of semi-norms \(\rVert f\lVert_{n}=\sup_{x\in[0,n]}|f(x)|\), or the in the case where \(I\subset\mathbb{R}_+\) is compact, the Banach space \(C([0,\infty)),\lVert\cdot\rVert_{\infty}\). \\\\

Hence, under the additional assumption that the path-space of the market information process \(Y\) is locally convex, we don't need the reduction to measurability with respect to a countably measurable sub-\(\sigma\)-algebra approximate \(X\) with a \textit{neural} random variable. In light of \cref{Remark:locally convex not fundamental} this condition is not fundamental either. It is in this sense that we say the use of \cref{prop:reduction to countable measurability} is not of fundamental importance in approximating random variables in \(L^p\). 
\bigbreak

\subsection{From simple processes to stochastic integral}\label{sec:From simple processes to stochastic integral}
As discussed in \cref{sec:main ideas Alexandar}, using the approximation theory for \(L^p\) random variables which are measurable with respect to the market information process, we will approximate elementary predictable processes. This will enable us to build up stochastic integration theory. The appropriate density properties of the approximated elementary predictable processes will then enable us to leverage continuity properties of the stochastic integral, as seen in \cref{.....}, to finally approximate stochastic integrals. For the convenience of the reader we recall the definition of elementary predictable processes and simple predictable processes.
\begin{definition}\label{def:elementary predictable process}
    We call \(H\) an \textit{simple predictable process} if there exists a finite, increasing sequence of \(\mathbb{F}\)-stopping times, \(0=\tau_0<\dots<\tau_{n+1}<\infty\) and a sequence of bounded random variables \((H^i)_{i=0}^n\subset\mathcal{L}^{\infty}(\mathcal{F}_{t_i},\mathbb{P})\) such that 
    \begin{equation}\label{eq:simple representation}
        H = H^01_{\left\{0\right\}}+\sum_{i=1}^nH^i1_{(\tau_i,\tau_{i+1}]}\;.
    \end{equation}
In particular, we have that \(H^i\) is \(\mathcal{F}_{\tau_i}\)-measurable. We denote the set of simple predictable processes by \(\mathcal{S}\) and we say \(H\in\mathcal{S}\) is associated with a sequence \((\tau_k)_{k=0}^{n+1}\). \\\\
If the sequence of stopping times is deterministic, we denote them by \(\tau_i=t_i\in\mathbb{R}\quad\forall i=0,\dots,n\) and call \(H\) a \textit{elementary predictable process} and denote the corresponding set by \(\mathcal{E}\). Analogously we say \(H\in\mathcal{E}\) is associated with a sequence \((t_k)_{k=0}^{n+1}\).
\end{definition}
For statements that are true for both \(\mathcal{E}\) and \(\mathcal{S}\) we shall formulate the result for \(\mathcal{S}\) only, as the case for \(\mathcal{E}\) results as a special case of the latter. Analogously we use \(\tau\) to denote both cases, where \(\tau\) is truly random and the case where it is deterministic. If we write \(t\) we only refer to a deterministic time. \\\\
In particular, as we are in a finite measure space setting, notice that \(\mathcal{L}^{\infty}(\mathcal{F}_{\tau_i})\subset\mathcal{L}^{p}(\mathcal{F}_{\tau_i})\) for any \(p\in[1,\infty]\). By \cref{prop:approximation of measurable maps by algorithmic strategies} we can find approximating sequences of \((H^i_n)_{n\in\mathbb{N}}\subset\Theta_{\mathcal{NN}}(\tau_i,\psi)\) that converge to \(H^i\in\mathcal{L}^{\infty}(\mathcal{F}_{\tau_i})\) in \(L^p\). Indeed, take any null-sequence of positive real numbers \((\epsilon_n)_{n\in\mathbb{N}}\) and apply \cref{prop:approximation of measurable maps by algorithmic strategies} (2) repeatedly. This motivates the definition of the following space
\begin{definition}
    Let \(H\) be a stochastic process with representation \cref{eq:simple representation}, where \(H^i\in\Theta_{\mathcal{NN}}(\tau_i,\psi)\) for \(i=0,\dots,n\). We denote this subspace of \(\mathcal{S}\) by \(\mathcal{S}(\psi)\) and call its elements: \textit{simple algorithmic strategies}.
\end{definition}
\begin{definition}
    Let \(H\) be a stochastic process with representation \cref{eq:simple representation}, where \(H^i\in\Theta_{\mathcal{NN}}(t_i,\psi)\) for \(i=0,\dots,n\). We denote this subspace of \(\mathcal{E}\) by \(\mathcal{E}(\psi)\) and call its elements: \textit{elementary algorithmic strategies}.
\end{definition}
The main motivation for the definition of these sets, and in a sense for all of \cref{sec:algorithmically generated random variables}, is that there exists a suitable topology in which \(\mathcal{E}(\psi)\) and \(\mathcal{E}(\psi)\) are dense in \(\mathcal{E}\) and \(\mathcal{S}\) respectively. By a suitable topology, we mean firstly that density actually holds and that it is compatible with the respective notion of stochastic integration; that is, that stochastic integration is continuous with respect to this topology. The first two candidates that come to mind are \(ucp\)-convergence and, in the square integrable setting, \(\mathcal{H}^2\)-convergence. Indeed, by (...) in the general setting and the Itô-isometry in the square integrable setting, the respective stochastic integral maps \(H\mapsto H\bullet X\) are continuous in the respective topologies. We treat both settings separatly as they involve convergence with respect to very different norms. We start with the \(ucp\) setting. 
\begin{proposition}\label{prop:p-ucp density of S}
    Vor \(H\in\mathcal{S}\) and \(\epsilon>0\), there exists \(H^{\epsilon}\in\mathcal{S}(\psi)\), such that \(\lVert(V-U_{\infty})^*\rVert_p<\epsilon\) for \(p\in[1,\infty)\). In other words \(S(\psi)\) is dense in \(\mathcal{S}\) for the \(p-ucp\) topology for \(p\in[1,\infty)\).
\end{proposition}
\begin{proof}
For \(V\in\mathcal{S}(\psi)\), consider \(H^i\) for some fixed \(i=0,\dots,n\). We have \(H^i\in\mathcal{L}^{\infty}(\mathcal{F}_{\tau_i})\subset \mathcal{L}^p(\mathcal{F}_{\tau_i})\). Using \cref{prop:approximation of measurable maps by algorithmic strategies}(2) there exists an approximating sequence \((H^i_m)_{m\in\mathbb{N}}\subset\Theta_{\mathcal{NN}}(\tau_i,\psi)\) in \(L^p\). Thus we consider the obvious candidate sequence for our approximation, given by \((H_m)_{m\in\mathbb{N}}\)
\begin{equation}
    H_m = H^0_m1_{\left\{0\right\}}+\sum_{i=1}^nH^i_m1_{(\tau_i,\tau_{i+1}]}\;.
\end{equation}
We have the obvious bound, uniformly in time
\begin{equation}
   (H-H^m)_{\infty}^*(\omega)\leq \sum_{i=1}^n|H^i(\omega)-H^i_m(\omega)|\quad\forall\omega\in\Omega\;.
\end{equation}
Hence, by the monotonicity of the \(L^p\) norm and the triangle inequality
\begin{equation}
    \lVert (H-H^m)_{\infty}^*\rVert_{p}\leq\sum_{i=1}^n\lVert H^i-H^i_m\rVert_p\;.
\end{equation}
But by construction, the right hand side converges to \(0\) as a finite sum of terms, each converging to \(0\). This shows density, as for every \(\epsilon>0\) we can find \(m\) big enough such that \(\lVert (H-H^m)_{\infty}^*\rVert_{p}<\epsilon\). 
\end{proof}
\begin{corollary}\label{cor:p-ucp density of S}
    \(S(\psi)\) is dense in \(\mathcal{S}\) for the \(\mathbb{P}-ucp\) topology.
\end{corollary}
\begin{proof}
    The standard way in probability theory to get from convergence in \(L^p\) to convergence in probability is via Chebyshev's inequality. We will show that \(L^p\) convergence of \( (H-H^m)_{\infty}^*\) implies convergence in probability of the latter, which proves a slightly stronger version of the claim. The claim thus follows from Chebychev's inequality. Indeed,
    \begin{equation}
        \mathbb{P}((|H-H^m)_{\infty}^*|>\epsilon)\leq\frac{1}{\epsilon^p}\int_{\Omega}|(H-H^m)_{\infty}^*|^p\mathbb{P}(\diff \omega)=(\frac{1}{\epsilon}\lVert (H-H^m)_{\infty}^*\rVert_p)^p
    \end{equation}
By the continuity of the function \(x\mapsto x^{1/p}\), for any \(p\in[1,\infty)\), the right hand side converges to \(0\) as \(m\to\infty\) by \cref{prop:p-ucp density of S}. The claim thus follows. 
\end{proof}
\
Notice that \cref{cor:ucp density of S} we can substitute the stopping times \(\tau_i\) with deterministic times \(t_i\) everywhere, and the proof goes through unchanged. Thus we get the immediate result for \(\mathcal{E}\) and \(\mathcal{E}(\psi)\) respectively.
\begin{corollary}\label{cor:ucp and p-ucp density of E}
    \(\mathcal{E}(\psi)\) is dense in \(\mathcal{E}\) for the \(\mathbb{P}-ucp\)-topology as well as the \(p-ucp\)-topology for \(p\in[1,\infty)\).
\end{corollary}
Following the results of \cref{thm:ucp continuity of general stochastic integral}, the simple processes are the basic building blocks of stochastic integration càdlàg processes with respect to semimartingales in the sense that the integral extends from simple integrands to càglàd integrands by continuity. By \cref{cor:ucp density of S}, the integral of \(S(\psi)\)-integrands also extends to càglàd integrands by continuity. Indeed this is the content of the following theorem
\begin{theorem}[Approximation of Stochastic Integrals]\label{thm:Approx of stoch. Integrals}
    Let \(X\) be a semimartingale, and \(\vartheta\in \mathbb{L}\). Then, there exists a sequence \((\vartheta^n)_{n\in\mathbb{N}}\) in \(\mathcal{S}(\psi)\) that converges to \(\vartheta\) in ucp, such that
    \begin{equation}
        \int \vartheta_s^n\diff X_s \to \int \vartheta_s\diff X_s,\quad n\to\infty
    \end{equation}
    where convergence holds in ucp.
\end{theorem}
\begin{proof}
    By \cref{thm:ucp continuity of general stochastic integral} the integral map
    \begin{align}
     I_X:\mathcal{S}_{ucp}\to\mathbb{D}_{ucp}   
    \end{align}
    is continuous. By \cref{prop:p-ucp density of S} \(\mathcal{S}(\psi)\) is dense in \(\mathcal{S}\) in the \(ucp\)-topology. Furthermore, by \cref{prop:ucp density of simple process in cadlag processes}. Hence, there exists a unique continuous extension to a map \(\tilde{I}_X:\overline{S}^{ucp}\to\mathbb{D}_{ucp}\) with \(\overline{S}^{ucp}=\mathbb{L}_{ucp}\). We have that \(\tilde{I}_X\) agrees with the map from \cref{def:stochastic integral w.r.t. semimartingales}. Indeed assume not. Denote by \(\tilde{\tilde{I}}\) the unique continuous extension of \(I:\mathcal{S}(\psi)\to\mathbb{D}\) and by \(\tilde{I}\) the unique continuous extension of \(I:\mathcal{S}\to\mathbb{D}\). Then there must exist \(x\in\mathbb{D}\) for which \(\tilde{\tilde{I}}(x)\neq\tilde{I}(x)\). By density \(\exists(x_n)_{n\in\mathbb{N}}\subset\mathcal{S}(\psi)\) such that \(\lim_{n\to\infty}x_n=x\) in \(ucp\) and in particular also as a sequence in \(\mathcal{S}\). But then by continuity
    \begin{equation}
        \tilde{\tilde{I}}(x)=\lim_{n\to\infty}\tilde{\tilde{I}}(x_n)=\lim_{n\to\infty}I(x_n)=\lim_{n\to\infty}\tilde{I}(x_n)=\tilde{I}(x).
    \end{equation}
    Hence the contradiction. 
\end{proof}
We goal is to formulate a deep mean-variance hedging result in our continuous time set up. We will formulate in in the square integrable setup where the theory is not too technical and the relevant objects have been discussed in \cref{sec:Stochastic Integration for square integrable Martingales}. In this context the relevant contuinity of the integral map is given by the Itô-isometry which asserts that for some \(X\in\mathcal{H}_0^2\) the linear operator
\begin{align*}
    I_X:& H\mapsto H\bullet X\\
    &\mathcal{H}^2\to L^2(X)
\end{align*}
is bounded; hence, it is (uniformly) continuous. We still need a similar density result of elementary or simple algorithmic strategies in \(L^2(X)\). Note however that this is a funadmentally different kind of density as the \(ucp\) density discussed in \cref{cor:ucp density of S} and \cref{cor:ucp denstiy of E} as convergence in \(L^2(X)\) is with respect to the integral norm 
\begin{equation}
    \lVert H\rVert_{L^2(M)}^2=\mathbb{E}[\int_0^\infty H_s^2\diff \langle M\rangle_s]
\end{equation}
Hence we need a different argument for density, which is reflected in the proof strategy of the following proposition
\begin{proposition}\label{prop:Lp(M) density of E(psi)}
    Assume there exists \(\epsilon>0\) such that \(M\in\mathcal{H}_0^{p+\epsilon}\). The the set \(\mathcal{E}(\psi)\) is dense in \(L^p(M)\) with respect to the seminorm on \(\lVert\cdot\rVert_{L^p(M)}\). Moreover, for every \(U\in L^p(M)\), there exists a sequence \((U_n)_{n\in\mathbb{N}}\subset\mathcal{E}(\psi)\), such that
    \begin{equation}
        \lim_{n\to\infty}\int U^n\diff M=\int U\diff M\;,
    \end{equation}
    where convergence holds i \(\mathcal{H}^p\). 
\end{proposition}
\begin{proof}
    First notice that we can assume without loss of generality that \(V\in L^p(M)\) is bounded. Indeed if not, use \(V^n=-n\vee V\wedge n\) instead. If we are able to prove the result for \(V^n\) the result follows for \(V\) by dominated convergence. \\
    Furthermore, recall from \cref{prop:density of E in Lp(M) via monotone class} that \(\mathcal{E}\) is dense in \(L^p(M)\) by monotone class, using  dominated convergence and the fact that \(\sigma(\mathcal{E})=\mathcal{P}\). Thus, for the frist claim suffices to show that \(\mathcal{E}(\psi)\) is dense in \(\mathcal{E}\) for the \(L^p(M)\) seminorm (again, we draw attention to the fact that we have shown the density of \(\mathcal{E}(\psi)\) in \(\mathcal{E}\) for the \(ucp\) topology, which is not the task at hand.) \\\\
    For this we will make use of the fact that \(M\in\mathcal{H}^p\), thus we have Doob's \(L^p\)-inequality and the Burkohlder-Davis-Gundy (BDG) inequality at hand.\\
    Fix \(\epsilon>0\).
    Recall from \cref{cor:ucp and ucp density of E} that for for every \(\epsilon>0\) and any \(q\in[1,\infty)\)there is a \(U\in\mathcal{E}(\psi)\) such that \(\lVert (V-U)^*_{\infty}\rVert_{q}<\epsilon\). For the convenience of the proof, we choose \(\lVert (V-U)^*_{\infty}\rVert_{q}<\epsilon^{1/p}\) where we take \(q=\frac{p(p+\epsilon)}{\epsilon}\). We use the trivial bound \((V_s(\omega)-U_s(\omega))\leq(V(\omega)-U(\omega))_{\infty}^*\quad\forall\omega\in\Omega\), to get 
    \begin{equation}
        \rVert V-U\lVert_{L^p(M)}^p=\mathbb{E}[\left(\int_{0}^{\infty}(V_s-U_s)^2\diff\langle M\rangle\right)^{p/2}]\leq\mathbb{E}[(V-U_{\infty}^*)^p\langle M\rangle_{\infty}^p/2]
    \end{equation}
    By Hölder for the Hölder conjugate pair \(\tilde{p}=\frac{p+\epsilon}{\epsilon}\) and \(\tilde{q}=\frac{p+\epsilon}{p}\), we further get the bound
    \begin{equation}
        \mathbb{E}[(V-U_{\infty}^*)^p\langle M\rangle_{\infty}^p/2]\leq\lVert(V-U_{\infty}^*)^p\rVert_{\tilde{p}}\lVert\langle M\rangle_{\infty}^p/2 \rVert_{\tilde{q}}\;.
    \end{equation}
    We now use first BDG and then Doob's \(L^p\)-inequality to finally get
    \begin{equation}
        \lVert(V-U_{\infty}^*)^p\rVert_{\tilde{p}}\lVert\langle M\rangle_{\infty}^{p/2} \rVert_{\tilde{q}}\lesssim \lVert(V-U)_{\infty}^*)\rVert_{q}^p\lVert M^*_{\infty}\rVert_{p+\epsilon}^p\leq\epsilon\lVert M\rVert_{\mathcal{H}^{p+\epsilon}}\;,
    \end{equation}
    where we don't care about the exact BDG constants as they only depend on \(p\) and not on the process. 
    As \(\mathcal{H}^{p+\epsilon}\) is finite by assumption, the latter can be made arbitrarily small, hence showing the claim. \\\\
    For the last claim, we use the fact that for every \(H\in L^p(M)\) we have \(H\bullet M\in\mathcal{H}^p\) by \parencite[Corollary.~12.3.6.]{CohenElliott2015StochasticCalculusApplications}. Furthermore, by the first part, for every \(V\in L^p(M)\) we can find a sequence \((U_n)_{n\in\mathbb{N}}\subset\mathcal{E}(\psi)\) that converges to \(U\) in \(L^p(M)\). Furthermore, we have \([\vartheta\bullet M]=H^2\bullet[M]\) (again by \parencite[Corollary.~12.3.6.]{CohenElliott2015StochasticCalculusApplications}). Hence by the BDG inequality we get 
    \begin{align*}
        \lVert \int U\diff M-\int U^n\diff M\rVert_{\mathcal{H}^p}=&\mathbb{E}[\left(\left(U-U^n)\bullet M\right)^*_{\infty}\right)^p]^{1/p}\lesssim \mathbb{E}[[\left(U-U^n)\bullet M\right)_{\infty}]^{p/2}]^{1/p}\\&=\mathbb{E}[\left(U-U^n)^2\bullet [M]\right)_{\infty}]^{p/2}]^{1/p}
        =\lVert U-U^n\rVert_{L^p(M)}\;,
    \end{align*}
    which concludes the proof. 
\end{proof}
As one of our main goals is to discuss the deep mean-variance hedging result in continuous time, we only really need \(p=2\). Nevertheless, we gave the statement for general \(p\in[1,\infty)\) as we want to show the correct choices for the Hölder exponents. 
\chapter{Deep Mean-Variance Hedging in Continuous Time}\label{sec:deep mean-square hedging in continuous time}

\section{Deep mean-square hedging in the martingale case}
We finally state one of the main results we were building up to: deep mean-variance hedging. It is formulated in the square integrable setup and thus should be seen complementary to the the mean-variance discussion of \cref{sec:Mean Variance Hedging in square integrable setting}. In particular the claim is that instead of looking among all of strategies \(V_0,\vartheta)\in\mathbb{R}\times L^2(M)\) we can restrict our search to \(\mathbb{R}\times\mathcal{E}(\psi)\). The error will make by only looking in that subset is aribtrarily small. The exact mathematical statement is as follows. 
\begin{theorem}[Continuous deep mean-square hedging in the martingale case]\label{thm:martingale setup, continuous deep variance hedging}
Assume there exists \(\epsilon>0\) such that \(M\in\mathcal{H}_0^{2+\epsilon}\), and let \(H\in L^2(\mathbb{P})\) be a real-valued \(\mathcal{F}_{\infty}\)-measurable random variable. Furthermore assume \(\mathcal{F}_0\) to be trivial.  Then :
\begin{equation}\label{eq:martingale setup, continuous deep variance hedging}
    \min_{c\in\mathbb{R},\,\vartheta\in L^2(M)}\mathbb{E}\Bigl[\bigl(H-c-\int_{0}^T\vartheta_t\diff M_t\bigr)^2\Bigr]
    \,=\,\inf_{c\in\mathbb{R},\,\vartheta\in \mathcal{E}(\psi)}\mathbb{E}\Bigl[\bigl(H-c-\int_{0}^T\vartheta_t\diff M_t\bigr)^2\Bigr]\;,
\end{equation}
where the minimum on the left-hand side is unique. 
In particular, there are two statements. First, the existence and uniqueness of the minimum for the left hand side of the equation, and second, the equality both expression. 
\end{theorem}
For the second claim, observe the following
\begin{remark}\label{rem:inf and sequences}
Assume the first claim has been shown. Then, as the minimum is achieved, there exist \(c^*\in\mathbb{R}\) and \(\vartheta^*\in L^2(M)\) such that \( \mathbb{E}\Bigl[\bigl(H-c^*-\int_{0}^T\vartheta_t^*\diff M_t\bigr)^2\Bigr] = \min_{c\in\mathbb{R},\,\vartheta\in L^2(M)}\mathbb{E}\Bigl[\bigl(H-c-\int_{0}^T\vartheta_t\diff M_t\bigr)^2\Bigr]\). Furthermore, it is clear that \( \min_{c\in\mathbb{R},\,\vartheta\in L^2(M)}\mathbb{E}\Bigl[\bigl(H-c-\int_{0}^T\vartheta_t\diff M_t\bigr)^2\Bigr]\) is a lower bound for the real subset \(\left\{\mathbb{E}\Bigl[\bigl(H-c-\int_{0}^T\vartheta_t\diff M_t\bigr)^2\Bigr]\;:\; c\in\mathbb{R},\vartheta\in \mathcal{E}(\psi)\right\}\). Indeed, \(\mathcal{E}(\psi)\subset L^2(M)\); hence, the first infimum is over a larger set. Hence, the second claim of \cref{thm:cont-deep-ms-hedging-martingale} reduces to showing that for every \(\delta>0\) there exists \(\vartheta\in\mathcal{E}(\psi)\) such that \(\mathbb{E}\Bigl[\bigl(H-c^*-\int_{0}^T\vartheta_t^*\diff M_t\bigr)^2\Bigr]+\delta>\mathbb{E}\Bigl[\bigl(H-c^*-\int_{0}^T\vartheta_t\diff M_t\bigr)^2\Bigr]\), or equivalently, there exists \((\vartheta^n)_{n\in\mathbb{N}}\subset\mathcal{E}(\psi)\) such that \(\lim_{n\to\infty}\mathbb{E}\Bigl[\bigl(H-c^*-\int_{0}^T\vartheta^n_t\diff M_t\bigr)^2\Bigr]=\mathbb{E}\Bigl[\bigl(H-c^*-\int_{0}^T\vartheta_t^*\diff M_t\bigr)^2\Bigr]\)\footnote{If this \(\delta\)-property holds, the sequence can be constructed by first choosing a null sequence \((\delta_n)_{n\geq 1}\), and for every \(n\in\mathbb{N}\), taking the corresponding \(\delta_n\)-close \(\vartheta^n\in\mathcal{E}(\psi)\). Conversely, assume there is such a sequence \((\vartheta^n)_{n\in\mathbb{N}}\subset\mathcal{E}(\psi)\); then \(\mathbb{E}\Bigl[\bigl(H-c^*-\int_{0}^T\vartheta^n_t\diff M_t\bigr)^2\Bigr]\) approximates \(\mathbb{E}\Bigl[\bigl(H-c^*-\int_{0}^T\vartheta_t^*\diff M_t\bigr)^2\Bigr]\) from above, as the latter is always smaller than \(\mathbb{E}\Bigl[\bigl(H-c^*-\int_{0}^T(\vartheta)_t\diff M_t\bigr)^2\Bigr]\) for any \(\vartheta\in\mathcal{E}(\psi)\). Hence, for every \(\delta\), there exists an \(N\in\mathbb{N}\) for which the delta property holds for all \(\vartheta^n\) with \(n\geq N\)} as a real sequence.
\end{remark}
Moreover, recall that the first claim has been shown in \cref{sec:Mean Variance Hedging in square integrable setting}. It remains to show the second part. 
\begin{proof}
    The first part of the claim has been shown in \cref{sec:Mean Variance Hedging in square integrable setting} and is a consequence of the Hilbert projection theorem. Denote by \((c^*,\vartheta^*)\in\mathbb{R}\times L^2(M)\) the unique strategy that achieves the minimum in \cref{eq:martingale setup, continuous deep variance hedging}. It remains to show the second part. \\\\
    By \cref{prop:Lp(M) density of E(psi)} for every \(\delta>0\) there exists \(U\in\mathcal{E}(\psi)\) such that \(\lVert \vartheta
    ^*\bullet M-\vartheta\bullet M\rVert_{\mathcal{H}^2}<\delta\). Hence
    \begin{align}
        \lVert H-c^*-(\vartheta\bullet M)_{\infty}\rVert_{2}&\leq\lVert H-c^*-(\vartheta^*\bullet M)_{\infty}\rVert_{2}+\lVert (\vartheta^*\bullet M)_{\infty}-(\vartheta\bullet M)_{\infty}\rVert_{2}\\
        &\leq\lVert H-c^*-(\vartheta^*\bullet M)_{\infty}\rVert_{2}+\lVert (\vartheta^*\bullet M)-(\vartheta\bullet M)\rVert_{\mathcal{H}^2}
        \\&\leq\lVert H-c^*-(\vartheta^*\bullet M)_{\infty}\rVert_{2}+\delta\;.
    \end{align}
    By \cref{rem:inf and sequences} this concludes the proof. 
\end{proof}
Notice that what made the proof of the second part work was essentially the continuity of the map \(\vartheta\mapsto(\vartheta\bullet M)_{\infty}\) is a continuous linear operator from \(L^2(M)\to  L^2(\mathcal{F}_{\infty},\mathbb{P})\).
\section{Extension for semimartingales}
Regarding \cref{thm:martingale setup, continuous deep variance hedging}, one is led to the natural question of whether one can obtain the claim for a more general setup. In this section, we will investigate a first approach to this. In particular, we want to allow for more general integrands. More precisely, we will allow a subclass of semimartingales as price processes \(M\). For this, we will have to adapt the space of integrable processes in a way that density results of the type \cref{cor:p-ucp density of S} and \cref{cor:ucp and p-ucp density of E}. We will first outline the general idea before giving concrete conditions and the corresponding extension.
\bigbreak
Consider the decomposition of a semimartingale
\begin{equation}
    X=X_0+M+A\;,
\end{equation}
where \(M\in\mathcal{H}^2_{0,loc}\): local martingale starting at \(0\) and \(A\in FV_{0}\): finite variation process starting at \(0\).
\subsection{Proof Attempts and Notes}
\cref{thm:cont-deep-ms-hedging-martingale} was in the realm of classical martingale stochastic calculus, with integrators \(M\in\mathcal{H}^{2}\)(even \(\mathcal{H}^{2+\epsilon}\) and integrands \(\vartheta\in L^2(M)\). In light of \cref{thm:Approx of stoch. Integrals}, we aim to prove a similar result as in \cref{thm:cont-deep-ms-hedging-martingale} for the more general case where \(M\) is a semimartingale. Following \parencite{Schweizer2001GuidedTourQuadraticHedging} we restrict the class of possible strategies to \(\vartheta_{GLP}\) such that \(G_T(\vartheta_{GLP})\) is a closed subspace in \(L^2(\mathbb{P})\). In this case, we know that the minimization problem has a solution, as \(\mathcal{A}=\mathbb{R}+G_T(\Theta_{GLP})\) is a closed subspace in \(L^2(\mathbb{P})\).; this means there exists an optimal trading strategy \((V_0^*,\vartheta^*)\in\mathcal{A}=\mathbb{R}+G_T(\Theta_{GLP})\) (\textcolor{blue}{I guess we don't even need strong conditions to get \(\mathcal{A}\) closed as we don't need explicit expression of the optimal trading strategy at first, hence we can directly project on \(\bar{\mathcal{A}}\) which is closed in \(L^2(\mathbb{P})\). In that line of thought, any \(\Theta\subset\Theta_2\) is sufficient for the existence of a solution in \(\bar{\mathcal{A}}\)}.) Then we have the following consequence of the Hilbert Projection theorem
\begin{theorem}[Hilbert Projection for Mean Sqaure Hedging]\label{thm:Hilbert Projection for Mean Sqaure Hedging}
    Let \(\Theta\subset\Theta_2\) and \(\mathcal{A}=\mathbb{R}+G_T(\Theta)\subset L^2(\mathbb{P})\). Then the \(L^2\)-closure \(\bar{\mathcal{A}}\) is a closed subspace of \(L^2(\mathbb{P})\). Hence, there exists \((V_0^*,\vartheta^*)\) such that 
    \begin{equation}
   \lVert H-V_0^*-\int_{0}^T\vartheta^*_s\diff X_s\rVert_{L^2}= \inf_{(c,\vartheta)\in\bar{\mathcal{A}}}\lVert H-c-\int_{0}^T\vartheta_s\diff X_s\rVert_{L^2}
    \end{equation}
    and the solution is unique. 
\end{theorem}
\begin{remark}
    Notice, we can study the related problem where we fix \(V_0\) (corresponds to the situation where the price \(V_0\) of the liability \(H\) might be determined by the market) and study
    \begin{equation}
        \inf_{\vartheta\in\Theta}\lVert H-V_0-\int_{0}^T\vartheta_s\diff X_s\rVert_{L^2}
    \end{equation}
    Again we can consider the the \(L^2\) closure of \(G_T(\Theta)=\left\{\int_0^T \vartheta_s\diff X_s:\vartheta\in\Theta\right\}\) and use the Hilbert Projection theorem on this closed subspace of \(L^2\).
\end{remark}
Then, crucially, notice that from \cref{thm:Approx of stoch. Integrals} we obtain the convergence in probability of the final values of the respective integral processes. Indeed, the \(\mathbb{P}\)-\mathbb{P}-ucp convergence of the integral processes, restricted to a finite time horizon T, reads as:
\begin{equation}
    \sup_{0\leq t\leq T}\left|\int_{0}^t\vartheta_s^n\diff X_s-\int_{0}^t\vartheta^*_s\diff X_s\right|\to 0\quad\text{in Probability}
\end{equation}
where \(\{\vartheta^n\}_{n\in\mathbb{N}}\in\mathcal{S}(\psi)\) is a minimizing sequence from \cref{thm:Approx of stoch. Integrals}
\\\\
In particular, this implies 
\begin{equation}
    \int_{0}^T\vartheta_s^n\diff X_s\xrightarrow{\mathbb{P}}\int_{0}^T\vartheta^*_s\diff X_s\;,
\end{equation}
where convergence holds in probability. Imposing the additional assumption that \(\left\{\left|\int_{0}^T\vartheta_s^n\diff X_s\right|^2\right\}_{n\in\mathbb{N}}\) is a uniformly integrable family of random variables  leads to the \(L^2\) convergence\footnote{In particular, the condition also implies that the random variable \(\int_{0}^T\vartheta_s^n\diff X_s\in L^2(\mathbb{P)}\quad\forall n\in\mathbb{N}\)}: 
\begin{equation}
\int_{0}^T\vartheta_s^n\diff X_s \xrightarrow{L^2(\mathbb{P})} \int_{0}^T\vartheta^*_s\diff X_s\;.
\end{equation}
We do NOT want the uniform integrability of \(\left\{\left|\int_{0}^T\vartheta_s^n\diff X_s\right|^2\right\}_{n\in\mathbb{N}}\) to follow from conditions on a sequence of trading strategies. Indeed, \cref{thm:Approx of stoch. Integrals} asserts the existence of a sequence \(\{\vartheta^n\}_{n\in\mathbb{N}}\subset\mathcal{S}(\psi)\). We don't want a theorem that says something like: 
\\\\
"If that sequence happens to be square uniformly integrable, then it follows...."
\\\\
Hence, if necessary, we shall impose conditions only on the integrator \(X\). Recall that for \(\vartheta_n\in\mathcal{S}(\psi)\) that is \(\vartheta^n_t=\varphi_0\mathbbm{1}_{\{0\}}+\sum_{i=1}^l\varphi_i^n\mathbbm{1}_{\tau_i,\tau_{i+1}}\) and \(\varphi_i^{(n)}=f^{(n)}(Y^{\tau_i}_{t_1},\dots,Y^{\tau_i}_{t_r})\) where \(f\in\mathcal{NN}((\mathbb{R}^{d\times r})^*,\psi)\)\footnote{the parameter \(r\in\mathbb{N}\) is not further specified. It relies to the existence result of \parencite[Thm.~5.23.]{Arandjelovic2025}}, the stochastic integral takes the form \(\int_0^T\vartheta^n_s\diff X_s = \sum_{i=1}^k \varphi^{(n)}_i(X_{\tau_i\wedge T}-X_{\tau_{i+1}\wedge T})\). Note that \(\left|\int_0^T\vartheta^n_s\diff X_s\right|^2 \leq \sum_{i=1}^k \left|\varphi^{(n)}_i\right|^2\left|(X_{\tau_i\wedge T}-X_{\tau_{i+1}\wedge T})\right|^2\). But there can't exist a uniform bound (in n) on \(|\varphi_i^n|\) as this would imply that neural networks can't approximate constants above this bound which is absurd. Also, notice for \(\tau<\sigma\) two stopping times, that \(f_n = c_n\mathbbm{1}_{[\tau,\sigma)}\in\mathcal{S}(\psi)\) for all \(c_n\in\mathbb{R}\). Take \(c_n\to\infty\) as \(n\to\infty\) a real deterministic sequence, then \(\int_0^T(f_n)_s\diff X_s=c_n(X_{\tau\wedge T}-X_{\sigma\wedge T})\), and hence \(\mathbb{E}[\left|\int_0^T(f_n)_s\diff X_s\right|^2]=\mathbb{E}[\left|c_n\right|^2\left|(X_{\tau\wedge T}-X_{\sigma\wedge T})\right|^2]\). If this is to be UI, then in particular, it is bounded in \(L^1\). But \(\sup_{n\in\mathbb{N}}\mathbb{E}[c_n^2(X_{\tau\wedge T}-X_{\sigma\wedge T})^2]=(\sup_{n\in\mathbb{N}}c_n^2)\mathbb{E}[(X_{\tau\wedge T}-X_{\sigma\wedge T})^2]<\infty\) by assumption. But \(\sup_{n\in\mathbb{N}}c_n^2=\infty\) hence we must have \(\mathbb{E}[(X_{\tau\wedge T}-X_{\sigma\wedge T})^2]=0\) which implies \(X_{\tau\wedge T}-X_{\sigma\wedge T}=0\quad\mathbb{P}\)-a.s. As this holds for all \(\mathbb{F}\)-stopping times \(\tau,\sigma\), in particular deterministic ones, we have \(X\equiv 0\) almost surely. Hence, there can't be a (non-trivial) condition on \(X\) making the family \(\left\{\left|\int_{0}^T\vartheta_s^n\diff X_s\right|^2\right\}_{n\in\mathbb{N}}\) UI, as the latter being UI for any sequence \((\vartheta^n)_{n\in\mathbb{N}}\) implies that \(X\equiv 0\).
\bigbreak
\textcolor{blue}{As an initial observation, recall that \(\int_{0}^T\vartheta_s^n\diff X_s\xrightarrow{\mathbb{P}}\int_{0}^T\vartheta^*_s\diff X_s\;\) implies that there exists a subsequence \(\left\{\int_{0}^T\vartheta_s^{n(k)}\diff X_s\right\}_{k\in\mathbb{N}}\) that converges \(\mathbb{P}\)-a.s. to \(\int_{0}^T\vartheta_s^{*}\diff X_s\). Hence, on a full measure set, it is (uniformly) bounded as a real converging sequence (assuming that \(\int_{0}^T\vartheta_s^{*}\diff X_s<\infty\) on that full measure set). If it is only \(\mathbb{P}\)-a.s. finite, take the intersection of the respective full measure sets to reach the same conclusion.) Hence, there exists a positive constant \(M\in\mathbb{R}_+\) such that \(\left|\int_{0}^T\vartheta_s^{n(k)}\diff X_s\right|^2<M^2\quad \mathbb{P}\)-a.s. for all \(k\in\mathbb{N}\). But \(M^2\) is integrable; hence, the family \(\left\{\left|\int_{0}^T\vartheta_s^{n(k)}\diff X_s\right|^2\right\}_{k\in\mathbb{N}}\) is uniformly integrable.\footnote{If \(H\geq |X_n|\) for all \(n\in\mathbb{N}\), we have \(\mathbb{E}[|X|\mathbbm{1}_{|X|\geq K}]\leq\mathbb{E}[H\mathbbm{1}_{H\geq K}]\) for any \(K>0\). As the left hand side going to 0 as \(K\to\infty\) it gives us the uniform integrability of \(\{X_n\}_{n\in\mathbb{N}}\)}. Thus, we get }
\\\\
\textcolor{red}{This conclusion is WRONG. Notice that the bound \(M = M(\omega)\) and hence doesn't define a random variable \(H\in L^2\)(at least not because it is constant almost surely). Indeed the argument used, goes as follows: For every \(\omega\) the sequence \(\left(\int_{0}^T\vartheta_s^{n(k)}\diff X_s\right)(\omega)\) is bounded as a real converging sequence, hence \(\left|\left(\int_{0}^T\vartheta_s^{n(k)}\diff X_s\right)(\omega)\right|\leq M(\omega)\).}
\\\\
\begin{equation}
\int_{0}^T\vartheta_s^{n(k)}\diff X_s \xrightarrow{L^2(\mathbb{P})} \int_{0}^T\vartheta^*_s\diff X_s\;\text{ as } k\to\infty
\end{equation}
Hence, restricting ourselves to this subsequence (which obviously still converges in probability to the same limit), we have 
\begin{align*}
    &\inf_{k\in\mathbb{N}}\Vert[H-V_0^*-\int_{0}^T\vartheta_s^{n(k)}\diff X_s]\rVert_{L^2}\leq \liminf_{k\to\infty}\lVert H-V_0^*-\int_{0}^T\vartheta_s^{n(k)}\diff X_s]\rVert_{L^2}\\&\leq \liminf_{k\to\infty}\lVert H-V_0^*-\int_{0}^T\vartheta_s^{*}\diff X_s\rVert_{L^2}+\lVert \int_{0}^T\vartheta_s^{*}\diff X_s -\int_{0}^T\vartheta_s^{n(k)}\diff X_s\rVert_{L^2}
    \\&=\lVert H-V_0^*-\int_{0}^T\vartheta_s^{*}\diff X_s\rVert_{L^2}
\end{align*}
where \(\lVert H-V_0^*-\int_{0}^T\vartheta_s^{*}\diff X_s\rVert_{L^2}=\inf_{\vartheta\in\Theta}\lVert H-V_0^*-\int_{0}^T\vartheta_s\diff X_s\rVert_{L^2}\).
By \cref{rem:inf and sequences}, the existence of this sequence is equivalent to a more general statement:
\begin{theorem}[Semimartingale Deep Hedging]\label{thm:general deep hedging}
    Let \(X\) be a semimartingale and \(\Theta\subset\Theta_2\). Let \((V_0^*,\vartheta^*)\) be the solution to \cref{thm:Hilbert Projection for Mean Sqaure Hedging}. Then there exists \((\vartheta_n)_{n\in\mathbb{N}}\subset\mathcal{S}(\mathcal{\psi})\) such that:
    \begin{equation}
        \min_{\vartheta\in \Theta}\lVert H-V_0^*-\int_{0}^T\vartheta_s\diff X_s\rVert_{L^2}=\inf_{\vartheta\in\mathcal{S}(\psi)}\lVert H-V_0^*-\int_{0}^T\vartheta_s\diff X_s\rVert_{L^2}\;,
    \end{equation}
\end{theorem}
where \(\lVert H-V_0^*-\int_{0}^T\vartheta^*_s\diff X_s\rVert_{L^2}=\min_{\vartheta\in \Theta}\lVert H-V_0^*-\int_{0}^T\vartheta_s\diff X_s\rVert_{L^2}\).
\\\\
\newpage
\underline{Possible strategies}:
In order to get deep-hedging result similar to \cref{thm:general deep hedging} we need 
\begin{equation}
    \int_0^T\vartheta^n_s\diff X_s \overset{L^2}{\to} \int_0^T\vartheta_s\diff X_s
\end{equation}
For an approximating sequence \((\vartheta_n)_{n\in\mathbb{N}}\subset\mathcal{S}(\psi)\) of \(\vartheta\in\tilde{\Theta}\), where the sense of approximation is to be determinded and the set \(\tilde{\Theta}\) such that \(\mathcal{S}(\psi)\subset\tilde{\Theta}\). Assuming this is done, we have
\begin{align*}
  &\inf_{\vartheta\in\mathcal{S}(\psi)}\lVert H-V_0^*-\int_0^T \vartheta_s^n\diff X_s\rVert_{L^2}\\&\leq
    \inf_{n\in\mathbb{\mathbb{N}}}\lVert H-V_0^*-\int_0^T \vartheta_s^n\diff X_s\rVert_{L^2}\leq \liminf_{n\to\infty}\lVert H-V_0^*-\int_0^T \vartheta_s^n\diff X_s\rVert_{L^2}
    \\&\leq \liminf_{n\to\infty}\lVert H-V_0^*-\int_0^T \vartheta_s^*\diff X_s\rVert_{L^2}+\lVert \int_0^T \vartheta_s^n\diff X_s-\int_0^T \vartheta_s^*\diff X_s\rVert_{L^2}\\&=\lVert H-V_0^*-\int_0^T \vartheta_s^*\diff X_s\rVert_{L^2}\;,
\end{align*}
where \(\lVert H-V_0^*-\int_0^T \vartheta_s^*\diff X_s\rVert_{L^2}=\argmin_{\vartheta\in\Theta}\lVert H-V_0^*-\int_0^T \vartheta_s^*\diff X_s\rVert_{L^2}\) for some set \(\Theta\subset\Theta_2\). On the other hand we have \(\inf_{\vartheta\in\tilde{\Theta}}\lVert H-V_0^*-\int_0^T \vartheta_s^n\diff X_s\rVert_{L^2}\leq\inf_{\vartheta\in\mathcal{S}(\psi)}\lVert H-V_0^*-\int_0^T \vartheta_s^n\diff X_s\rVert_{L^2}\), because \(\mathcal{S}(\psi)\subset\tilde{\Theta}\). Thus we get equality. \\
We consider two possible strategies.
\bigbreak
(1): From \parencite[Thm~5.29.]{Arandjelovic2025} we have \(\mathcal{S}(\psi)\) is dense in \(\mathcal{S}\) which is dense in \(\mathcal{L}\), the space of càglàd processes (see Protter for this). Then for \(\vartheta_n\overset{\mathbb{P}-ucp}{\to}\vartheta\) where \(\vartheta\in \mathcal{L}\), we have 
\begin{align*}
    \int \vartheta_s^n\diff X_s\overset{\mathbb{P}-ucp}{\to}\int \vartheta_s\diff X_s\implies\int_0^T \vartheta_s^n\diff X_s\overset{\mathbb{P}}{\to}\int_0^T\vartheta_s\diff X_s
\end{align*}
and 
\begin{align*}
    &\left(\vartheta_s^n\diff X_s\overset{\mathbb{P}}{\to}\int_0^T\vartheta_s\diff X_s \right) + \left(\left\{\left|\int_{0}^T\vartheta_s^{n}\diff X_s\right|^2\right\}_{n\in\mathbb{N}}\text{uniformly integrable}\right)\\&\implies \left(\int_0^T\vartheta_s^n\diff X_s\overset{L^2}{\to}\int_0^T\vartheta_s\diff X_s \right)
\end{align*}
WARNING: Notice that for sequences of the form \(\vartheta_n=c_n\cdot\mathbbm{1}_{[\tau,\sigma)}\in\mathcal{S}(\psi)\) for \(\lim_{n\to\infty}c_n=\infty\), requiring that \(\left\{\left|\int_{0}^T\vartheta_s^{n}\diff X_s\right|^2\right\}_{n\in\mathbb{N}}\) is uniformly integrable implies \(X=0\) as a process. This is a pathological example as it doesn't converge to an interesting limiting process \(\vartheta\), which we assume. This just shows that  \(\left\{\left|\int_{0}^T\vartheta_s^{n}\diff X_s\right|^2\right\}_{n\in\mathbb{N}}\) can't be UI for all \((\vartheta_n)_{n\in\mathbb{N}}\subset\mathcal{S}(\psi)\). 
\bigbreak
(2): Notice that \(Y_n\overset{L^2-ucp}{\to}Y\) i.e. \(\lim_{n\to\infty}\mathbb{E}[\sup_{0\leq t\leq T}|Y_t^n-Y_t|^2]=0\) implies \(Y^n_T\overset{L^2}{\to}Y_T\). Hence, we want a version of \parencite[Thm~5.29]{Arandjelovic2025} with \(L^2-ucp\) convergence of the integrals. For this, we might want to use the fact that for \(\Phi(\cdot)=(\cdot)^2\), \parencite[Thm~5.28]{Arandjelovic2025} implies that \(\mathcal{S}(\psi)\) is \(L^2-ucp\) dense in \(\mathcal{S}\) and use some kind of \(L^2-ucp\) continuity of the stochastic integral \(H\to H\bullet X\) for well chosen classes for \(H\) and \(X\). An clearly weaker assumption would be to ask for convergence in \(\mathcal{H}^2\) which implies the \(L^2\) convergence of the end-points (Doob). A well known case where we have the continuity of the integral map w.r.t. to the \(L^2\)-convergence of the end-points is the isometry property (Itô isomtrey) of the map 
\begin{align}
    \begin{cases}
        H\to H\bullet M\\
        L^2(M)\to\mathcal{H}^2_0
    \end{cases}
\end{align}
for \(M\in\mathcal{H}^2_{0,loc}\)
Indeed in this case, by \parencite[Thm~5.32]{Arandjelovic2025} we even have that \(\mathcal{E}(\psi)\) is dense in \(L^p(M)\) for \(p\in[1,\infty)\) if \(M\in\mathcal{H}_0^{p+\epsilon}\) for some \(\epsilon>0\). (Then we also have that \(\mathcal{S}(\psi)\) is dense in \(L^p(M)\)). Hence we have 
\begin{equation}
    \lVert H\bullet M\rVert_{\mathcal{H}^2}=\lVert H\rVert_{L^2(M)}
\end{equation}
Using the sequence \(H_n=\vartheta_n-\vartheta\) and the fact that convergence in \(\mathcal{H}^2\) implies convergence in \(L^2-ucp\) (by Doob's maximal inequality, which is applicable because \(H\bullet M\) is a square integrable martingale in this case) we recover \parencite[Thm~5.35]{Arandjelovic2025}.
\\\\
To get a slight extension of \parencite[Thm~5.35]{Arandjelovic2025}, we notice that we don't need the full isometry property; rather, we need the boundedness (continuity) of the map \(H\to H \bullet M\) and the corresponding density of \(\mathcal{S}(\psi)\). Hence, for a semimartingale 
\begin{equation}
    X = X_0 + A + M
\end{equation}
for \(A\) a finite-variation process and \(M\in\mathcal{H}^2_{0,loc}\) (the addition of \(A\) destroys the Itô-isometry property) we try to bound 
\begin{equation}
     \lVert \left(H\bullet X\right)_T\rVert_{L^2}\leq K_X\lVert H\rVert_{L^2(M)}
\end{equation}
for some constant \(K_X = K(X)\). By \(|a+b|^2\leq 2|a|^2+2|b|^2\) it is enough to bound \(\lVert H\bullet A\rVert_{\mathcal{H}^2}\) and \(\lVert H\bullet M\rVert_{\mathcal{H}^2}\) both by \(\lVert H\rVert_{L^2(M)}\) and some multiplicative constant. Notice that for \(M\in\mathcal{H}^2_{0,loc}\) and \( H \in L^2(M)\), \(H\bullet M\) is a square integrable martingale; hence, we can bound \(\lVert H\bullet M\rVert_{\mathcal{H}^2}=\lVert (H\bullet M)_T\rVert_{L^2}\leq 4\lVert H\rVert_{L^2}\) by Doob. For the finite variation part, we might need assumptions on \(A\) (and hence on \(X\)).
\\
Based on this oberservations we consider the following approach. Consider a new space \(L^2(M,A)\) for a semimartingale \(X=X_0+M+A\) where \(A\) is a finite-variation process and \(M\in\mathcal{H}^2_{0,loc}\). Define it by \footnote{Note that in our case \([X]=[M]\) and \(L^2(M,A)\subset L^2(M)\), so \(\int_0^T\vartheta_s\diff M_s\in\mathcal{H}_0^2\); hence, \(\left(\int \vartheta_s\diff M_s\right)^2-[\int \vartheta_s\diff M_s]\) is a martingale starting at 0, and \([\int\vartheta_s\diff M_s]=\int \vartheta_s^2\diff [M]_s\). Hence, \(\mathbb{E}[\left(\int_0^T \vartheta_s\diff M_s\right)^2]=\mathbb{E}[\int_0^T \vartheta_s^2\diff [M]_s]\)}
\begin{align*}
H\in L^2(M,A) \iff \left(\mathbb{E}[\int_0^T H_s^2\diff [X]_s]+\mathbb{E}[\left(\int_0^T |H|_s\diff |A|_s\right)^2]\right)^{\frac{1}{2}}\coloneqq\lVert H\rVert_{L^2(M,A)}<\infty
\end{align*}\footnote{Inspired by the fact that for \(\int H_s\diff X_s = \int H_s\diff M_s+\int H_s\diff A_s\) and \(\mathbb{E}[\left(\int H_s\diff M_s+\int H_s\diff A_s\right)^2]\leq 2\mathbb{E}[\left(\int H_s \diff M_s\right)^2]+2\mathbb{E}[\left(\int H_s \diff A_s\right)^2]\leq2\cdot 4\mathbb{E}[\left(\int (H_s)^2 \diff [M]_s\right)^2]+2\mathbb{E}[\left(\int H_s \diff A_s\right)^2]\) from \(|a+b|^2\leq 2|a|^2+2|b|^2\) and BDG with \(C_2 = 4\).}\textcolor{blue}{One would still need to check that \(\lVert H\rVert_{L^2(M,A)}\) defines a norm}
where in our case \([X_0 + A + M]=[M]\). With this we have 
\begin{equation}
    \lVert\int_0^T\vartheta_s\diff X_s\rVert_{L^2}\leq K_X\lVert \vartheta \rVert_{L^2(M,A)}
\end{equation}
Indeed,
\begin{align*}\label{eq:bounding the L^2 norm of the stochastic integral}
    &\mathbb{E}[\left(\int_0^T H_s\diff X_s\right)^2]=\mathbb{E}[\left(\int_0^T H_s\diff M_s+\int_0^T H_s\diff A_s\right)^2]
    \\&
    \leq 2\mathbb{E}[\left(\int_0^T H_s \diff M_s\right)^2]+2\mathbb{E}[\left(\int_0^T H_s \diff A_s\right)^2]
    \\&
    \leq2\cdot 4\mathbb{E}[\int_0^T (H_s)^2 \diff [M]_s]+2\mathbb{E}[\left(\int_0^T H_s \diff A_s\right)^2]\leq
    \\&
     \leq2\cdot 4\mathbb{E}[\int_0^T (H_s)^2 \diff [M]_s]+2\mathbb{E}[\left(\int_0^T |H|_s \diff |A|_s\right)^2]\leq
    \lVert H\rVert_{L^2(M,A)}
\end{align*}
from \(|a+b|^2\leq 2|a|^2+2|b|^2\) and BDG with \(C_2 = 4\)\footnote{\(\mathbb{E}[\left(\int_0^TH_s\diff M_s\right)^2]\leq\mathbb{E}[\sup_{0\leq t\leq T}\left(\int_0^t H_s\diff M_s\right)^2]\leq C_2\mathbb{E}[\int_0^T(H_s)^2\diff [M]_s]\)}. Then we have
\begin{align}
    &\vartheta_n\overset{L^2(M,A)}{\to}\vartheta\implies\int_0^T\vartheta_s^n\diff X_s \overset{L^2}{\to}\int_0^T\vartheta_s\diff X_s
\end{align}
In this approach, we need to show that \(\lVert\cdot \rVert_{L^2(M,A)}\) is indeed a norm (or a semi-norm), that \(\mathcal{S}(\psi)\) is dense in \(L^2(M,A)\), and whether \(L^2(M,A)\) is rich enough to be interesting.
\section{Proof Attempt}
We present the first version of a generalization of \parencite[Thm~5.35]{Arandjelovic2025}. From \parencite[Thm.~5.32.]{Arandjelovic2025} we know that \(\mathcal{E}(\psi)\) is dense in \(L^p(M)\) for \(p\in[1,\infty)\) and \(M\in\mathcal{H}_0^{p+\epsilon}\) for some \(\epsilon>0\). 
Now consider a semimartingale \(X=X_0+M+A\), where \(M\in\mathcal{H}_0^{p+\epsilon}\) and \(A\) are adapted processes with RCLL trajectories, \(A_.(\omega)\) of finite variation and \(A_0=0\). Additionally, assume the following:
\begin{itemize}
    \item \(A<<\langle M\rangle\) path-wise (i.e. for every \(\omega\) as Lebesgue Stieltjes measures, with density process \(\lambda\).
    \item \(\operatorname*{ess\sup}\left(\langle M\rangle_{\infty}(\lambda^2+1)^*_{\infty}\right)<\infty\) where \((\lambda^2+1)^*_{\infty}=\sup_{t\geq 0}(\lambda^2+1)\).\footnote{Replace \(\infty\) by \(T\) in the finite horizon setting.}
\end{itemize}
Then we have, by density, of \(\mathcal{E}(\psi)\) in \(L^2(M)\) that
\begin{equation}
    \forall H\in L^2(M), \exists (H_n)_{n\in\mathbb{N}}\subset\mathcal{E}(\psi):\lim_{n\to\infty}\lVert H-H_n\rVert_{L^2(M)}=0\;.
\end{equation}
We can bound the \(L^2\) difference of the integrals by 
\begin{align}
    &\mathbb{E}[(H\bullet X)_{\infty}^2]\leq 2\left((\mathbb{E}[(H\bullet \langle M\rangle)_{\infty}^2]+\mathbb{E}[(H\bullet A)_{\infty}^2]\right)\\
    &\leq 2\left((\mathbb{E}[\langle M\rangle_{\infty}(H^2\bullet \langle M\rangle)_{\infty}]+\mathbb{E}[(H\lambda\bullet \langle M\rangle)_{\infty}^2]\right)\\
    &\leq2\left((\mathbb{E}[\langle M\rangle_{\infty}(H^2\bullet \langle M\rangle)_{\infty}]+\mathbb{E}[\langle M\rangle_{\infty}(H^2\lambda^2\bullet \langle M\rangle)_{\infty}]\right)\\
    &=2\mathbb{E}[\langle M\rangle_{\infty}(H^2(1+\lambda^2)\bullet\langle M\rangle)_{\infty}]\\
    &\leq 2\mathbb{E}[\langle M\rangle_{\infty}(1+\lambda^2)^{*}_{\infty}(H^2\bullet\langle M\rangle)_{\infty}]\\
    &\leq2\operatorname*{ess\sup}(\langle M\rangle_{\infty}(1+\lambda^2)^{*}_{\infty})\cdot\mathbb{E}[H^2\bullet \langle M\rangle]\\
    &=C(X)\lVert H\rVert_{L^2(M)}^2\;.
\end{align}
We used Cauchy Schwarz in step 2 and 3, along with the absolute continuity of \(A<<\langle M\rangle\). In step 5, we used the Hölder inequality with \(p=1,q=\infty\).\\
By assumption \(C(X)<\infty\), hence substituting for \(H\), \(Z^n=H-H^n\), we get that 
\begin{equation}
    \lim_{n\to\infty}\mathbb{E}[(Z^n\bullet X)_{\infty}^2]\leq\lim_{n\to\infty}C(X)\lVert H\rVert_{L^2(M)}^2=0
\end{equation}
By the argument given earlier, this allows for a deep mean-variance hedging result for a subclass of semimartingales.
\section{Different Approach}
We present a slightly different approach based on the observation \cref{eq:bounding the L^2 norm of the stochastic integral}. Notice that under the assumption that \(A<<\langle M\rangle\) path-wise (i.e. for every \(\omega\) as Lebesgue Stieltjes measures, with density process \(\lambda\) we have:
\begin{align*}
&\lVert H\rVert_{L^2(M,A)}^2=\mathbb{E}[\int_0^\infty (H_s)^2\diff [M]_s]+\mathbb{E}[\left(\int_0^\infty |H_s||\lambda_s|\diff[M]_s\right)^2]\\&\leq\mathbb{E}[\int_0^\infty (H_s)^2\diff [M]_s+\mathbb{E}[\int_0^{\infty}(H_s)^2\diff [M]_s\cdot\int_0^{\infty}(\lambda_s)^2\diff [M]_s]
\\&
\leq\mathbb{E}[\int_0^\infty (H_s)^2\diff [M]_s+\mathbb{E}[\left(\int_0^{\infty}(H_s)^2\diff [M]_s\right)^{p}]^{1/p}\cdot\mathbb{E}[\left(\int_0^{\infty}(\lambda_s)^2\diff [M]_s\right)^{q}]^{1/q}
\;,
\end{align*}
where we used the fact that if \(A<<[M]\) with density \(\lambda\) then \(|A|<<[M]\) with density \(|\lambda|\). Furthermore we used the Kunita-Watanabe inequality in the second step and Hölder's inequality in the third step. In the case \(p=1,q=\infty\) the above reads
\begin{equation}
    \lVert H\rVert_{L^2(M,A)}\leq\lVert H\rVert_{L^2(M)}\left(1+\operatorname*{ess\sup}\left(\int_0^{\infty}(\lambda_s)^2\diff [M]_s\right)\right)\;.
\end{equation}
Notice \(\int_0^{\infty}(\lambda_s)^2\diff [M]_s=\int_0^{\infty}|\lambda_s|\diff A_s\leq \int_0^{\infty}|\lambda_s|\diff |A|_s\). We impose the additional assumption
\begin{equation}
    \operatorname*{ess\sup}\int_0^{\infty}|\lambda_s|\diff |A|_s<\infty\;.
\end{equation}
Then, because we have the trivial inequality
\begin{align}\label{eq:essential boundedness assumption}
    \mathbb{E}[\int_0^{\infty}H_s^2\diff[M]_s]\leq  \mathbb{E}[\int_0^{\infty}H_s^2\diff[M]_s]+  \mathbb{E}[\left(\int_0^{\infty}|H_s||\lambda_s|\diff[M]_s\right)^2]=\lVert H\rVert_{L^2(M,A)}
\end{align}
we have that the norms on \(L^2(M)\) and \(L^2(M,A)\),under the condition that \(X = X_0+M+A\) with \(A<<[M]\) and \(M\in\mathcal{H}^{p+\epsilon}_0\) for some \(p\in[0,\infty)\) and \(\epsilon>0\), along with the additional assumption \cref{eq:essential boundedness assumption}, are equivalent. Because both spaces \(L^2(M)\) and \(L^2(M,A)\) are subspaces of the space of \(\mathcal
P\)-measurable functions in, where \(\mathcal{P}\) is the predictable \(\sigma\)-algebra on \(\Omega\times[0,\infty)\), they are defined through their norms; the spaces are equal:
\begin{equation}
    L^2(M) =L^2(M,A)\;.
\end{equation}
Hence, density of \(\mathcal{E}(\psi)\) in \(L^2(M)\) is obviously equivalent to density of \(\mathcal{E}(\psi)\) in \(L^2(M,A)\) under the stated assumptions on \(M\) and \(X\). Hence, for every \(H\in L^2(M,A)\) there exists \((H_n)_{n\in\mathbb{N}}\subset \mathcal{E}(\psi)\) such that \(\lim_{n\to\infty}\lVert H-H_n\rVert_{L^2(M,A)}=0\). Then we have, by the observation \cref{eq:bounding the L^2 norm of the stochastic integral}, with \(T=\infty\) and \(H_s\coloneqq (H-H_n)_s\)
\begin{equation}
    \lVert (H-H^n)\bullet X\rVert_{L^2}
    \lesssim\lVert H-H^n\rVert_{L^2(M,A)}\;,
\end{equation}
thus
\begin{equation}
    \lim_{n\to\infty}\lVert (H-H^n)\bullet X\rVert_{L^2}=0
\end{equation}
\bigbreak
The two approaches differ in their additional assumption
\begin{enumerate}
    \item \(\operatorname*{ess\sup}(\langle M\rangle_{\infty}(1+\lambda^2)^{*}_{\infty})<\infty\)
    \item  \(\operatorname*{ess\sup}\int_0^{\infty}|\lambda_s|\diff |A|_s<\infty\)
\end{enumerate}
Notice that
\begin{equation}
    \int_0^\infty(\lambda_s)^2\diff [M]_s\leq\int_0^{\infty}(1+\lambda_s^2)\diff[M]_s\leq(1+\lambda_s^2)_{\infty}^*[M]_{\infty}\,
\end{equation}
where we used the non-negativity and monotonicity of the process \([M]\). Hence, (2.) is a weaker assumption than (1.).
\textcolor{blue}{potential problem with this approach. The norms are equivalent but not really norms. First one needs to check that \(\lVert \cdot\rVert_{L^2(M,A)}\) defines at least a subadditiv, homogenous map and then that the process which are identified by \(\lVert X-Y\rVert_{L^2(M,A)}=0\) are also identified in \(L^2(M)\), that is \(\lVert X-Y\rVert_{L^2(M)}=0\). So that the norms are norms on the same quotient space under some equivalence relation identifying processes}
First, notice that \(\lVert H\rVert_{L^2(M,A)}=0\) is subadditive
Under the assumption that \(A<<[M]\) has density \(\lambda\), we have that \(\lVert X-Y\rVert_{L^2(M)}=0\implies\lVert X-Y\rVert_{L^2(M,A)}=0\). This follows from the fact that \(c\lVert X\rVert_{L^2(M)}\leq\lVert X\rVert_{L^2(M,A)}\leq C\lVert X\rVert_{L^2(M)}\) for \(c=1,C=(1+\operatorname*{ess\sup}\int_0^{\infty}|\lambda_s|\diff |A|_s)\). By assumption, \(c,C\) are finite. This implies that under the equivalence relation \(X\sim Y\) for \(X\sim Y\) for \(\lVert X-Y\rVert_{L^2(M,A)}=0\). Hence, they are norms on the same space \(L^2(\Omega\times [0,T],\mathcal{P})/\sim\), and in particular, it makes sense to say that they are equivalent. 
\bigbreak
This proves the following result:
\begin{theorem}[Semimartingale Deep Hedging]\label{thm:general deep hedging}
    Let \(X=X_0+M+A\) be a semimartingale with \(M\in \mathcal{H}_{0}^{2+\epsilon}\) for some \(\epsilon>0\), \(A<<[M]\) with density \(\lambda=(\lambda_s)_{s\in\mathbb{R}_+}\) and with \(\operatorname*{ess\sup}\int_0^{\infty}|\lambda_s|\diff |A|_s<\infty\) \. Let \(H\in L^2(\mathbb{P})\) and \((V_0^*,\vartheta^*)\) be the solution to \cref{thm:Hilbert Projection for Mean Sqaure Hedging}. Then there exists \((\vartheta_n)_{n\in\mathbb{N}}\subset\mathcal{E}(\mathcal{\psi})\) such that:
    \begin{equation}
        \min_{\vartheta\in \Theta}\lVert H-V_0^*-\int_{0}^T\vartheta_s\diff X_s\rVert_{L^2}=\inf_{\vartheta\in\mathcal{E}(\psi)}\lVert H-V_0^*-\int_{0}^T\vartheta_s\diff X_s\rVert_{L^2}\;,
    \end{equation}
\end{theorem}
where \(\lVert H-V_0^*-\int_{0}^T\vartheta^*_s\diff X_s\rVert_{L^2}=\min_{\vartheta\in \Theta}\lVert H-V_0^*-\int_{0}^T\vartheta_s\diff X_s\rVert_{L^2}\).
\begin{proof}
    The proof is given above and relies on the fact that in this case \(L^2(M)=L^2(M,A)\). By \parencite[Thm.~5.32]{Arandjelovic2025} \(\mathcal{E}(\psi)\) is dense in \(L^2(M)\) and hence also trivially in \(L^2(M,A)\) as well. The only difference is that the optimal \(V_0^*,H^*\) are for the more general problem with \(X= X_0+M+A\). The proof is the same as in \parencite[THm.~5.35]{Arandjelovic2025} as for any \(\delta>0\) there exists \(U\in\mathcal{E}(\psi)\) such that \(\lVert(U^*\bullet X)_{\infty}-(U\bullet X)_{\infty}\rVert_{L^2}\leq \lVert U^*-U\rVert_{L^2(M,A)}<\delta\) by the calculation in \cref{eq:bounding the L^2 norm of the stochastic integral}.
\end{proof}
One should give examples, if they exist, of finite variation processes \(A\) satisfying the assmuptions:
\begin{itemize}
    \item \(A<<[M]\) with density \(\lambda\)
    \item \(\operatorname*{ess\sup}\int_0^{\infty}|\lambda_s|\diff |A|_s<\infty\)
\end{itemize}

Furthermore we will think of a more general approach generalizing \parencite[Thm.~5.35]{Arandjelovic2025} again with \(X\) a semimartingale and approximants from \(\mathcal{S}(\psi)\). For this we might think again about strategy (1) and if there are sensible conditions on \(\vartheta\) that transfer to conditions on \(\vartheta_n\) such that the respective family of stochastic integrals (more precisely the square processes) become a uniformly integrable family. 
\textblue{Possible, non-trivial process satisfying assumptions: Let 
\begin{equation}
    S_t = \lambda B_t + \mu t
\end{equation}
Then \(\langle S\rangle_t= \lambda^2 t\) and \(A_t=\mu t\)
\chapter{Conclusion and Outlook}
\section{Summary of contributions}
\section{outlook}

\section{Possible extensions if time permits}
Extension of \parencite[Prop.~5.15.]{Arandjelovic2025} to 
bounded measurable non-constant activation functions
based on \parencite[Thm.~5]{Hornik1991}, using tools from \parencite{Rudin1990FourierGroups}.


% =================================================
% Bibliography
% =================================================
\cleardoublepage
\printbibliography

% =================================================
% Appendix
% =================================================
\begin{appendices}
\chapter{Mathematical tools}
Chebychev inequality, Doob's \(L^p\)-inequality, Burkholder-Davis-Gundy inequality (BDG), functional monotone class theorem, Bitcheler-Dellacherie Thm
\chapter{Proofs deferred from the main text}
\section{Equality of Infima under closure and continuous mapping}\label{Appendix:Infima over closure}
\begin{proposition}\label{prop:minimization of closure}
     Let X be a normed space, \(Y\subset X\) a subset and \(d:X\to\mathbb{R}\) continous. Then 
     \begin{equation}
         \inf d(Y)=\inf d(\overline{Y})\;.
     \end{equation}
\end{proposition}
\begin{proof}
    Because \(Y\subset\overline{Y}\) we have \(\inf d(\overline{Y})\leq\inf d(Y)\). For the other direction we need continuity of \(d:X\to\mathbb{R}\). Let \(y\in\overline{Y}\). There exists \((y_n)_{n\in\mathbb{N}}\subset Y\) such that \(\lim_{n\to\infty}y_n = y\) in \(X\). By continuity of \(d\) we thus have \(\lim_{n\to\infty}d(y_n) = d(y)\). Furthermore, by definition,
    \begin{equation}
        d(y)=\lim_{n\to\infty}d(y_n)\geq\inf d(Y)\;.
    \end{equation}
    Taking the infimum over \(\overline{Y}\) on both sides we have,
    \begin{equation}
        \inf d(\overline{Y}) \geq \inf d(Y);.
    \end{equation}
\end{proof}
Notice that the above proof also works if \(d\) is only upper-semi continuous. Then we have \(d(y)\geq\limsup_{n\to\infty}{d(y_n)}\) and we get the same conclusion. Furthermore, we can generalize the above to a general topological space \(Y\) by replacing sequences by nets, as all we have used is that for \(y\in\overline{Y}\) there exists \((y_n)_{n\in\mathbb{N}}\subset Y\) such that \(\lim_{n\to\infty}y_n = y\) in \(X\)
\section{Direct proof of Finite Deep mean variance Hedging}
\begin{proof}
    Recall that 
    \begin{align}
        &\Theta=\left\{\vartheta=(\vartheta_k)_{k=0}^{n-1}:\vartheta_k \text{ is }\mathcal{F}_{k-1}\text{-measurable}\right\}\\
        &\Theta_M\left\{\vartheta=(\vartheta_k)_{k=0}^{n-1}:\vartheta_k=F_k(Y_0,\dots Y_{k-1}),F_k\in\Psi_M\right\}\\
        &\Theta_{\infty}=\bigcup_{M\in\mathbb{N}}\Theta_M\;.
    \end{align}
    By the Doob-representation Lemma we have for \(\vartheta\in\Theta\) that there exists \(f_k:\mathbb{R}^{rk}\to\mathbb{R}\),  (Borel-)measurable such that
    \begin{equation}
        \vartheta_k=f_k(Y_0,\dots Y_{k-1})\;.
    \end{equation}
    Because \(\Omega\) is finit we have that \(f_k\in L^2(\mathbb{R}^{rk},\mu)\), where \(\mu=\mathcal{L}(Y_0,\dots Y_{k-1})\). Hence, we see that 
    \begin{equation}
        \Theta=\bigoplus_{k=1}^{n}L^2(\Omega,\mathcal{F}_{k-1},\mathbb{P})\;.
    \end{equation}
    Furthermore let
    \begin{equation}
        \Xi_k=\left\{X:\mathcal{L}^0(\Omega,\mathcal{F}):X=F_k(Y_0,\dots,Y_{k-1})\text{ where } F_k\in\Psi\right\}
    \end{equation}
    Then,
    \begin{equation}
        \Theta_{\infty}=\bigoplus_{k=1}^{n}\Xi_k\;.
    \end{equation}
    By density of of \(\Psi\) in \(L^2\) we have \(\overline{\Xi_k}^{L^2}=L^2(\Omega,\mathcal{F}_k,\mathbb{P})\).
\end{proof}
We define the following norm on \(\Theta\):
\begin{equation}
    \lVert(\vartheta_1,\dots,\vartheta_{n})\rVert\coloneqq\sum_{k=1}^{n}\lVert \vartheta_k\rVert_{L^2}
\end{equation}
Then we have that 
\begin{equation}
    \overline{\Theta}_{\infty}=\bigoplus_{k=1}^{n}\overline{\Xi_k}^{L^2}=\bigoplus_{k=1}^{n}L^2(\Omega,\mathcal{F}_k,\mathbb{P})=\Theta\;,\footnotemark
\end{equation}\footnotetext{This can be checked by hand. The inclusion \(\overline{\bigoplus_{k=1}^{n}\Xi_k}\subset\bigoplus_{k=1}^{n}\overline{\Xi_k}^{L^2}\) follows from the fact that \(\bigoplus_{k=1}^{n}\overline{\Xi_k}^{L^2}\) is closed w.r.t. to the product topology which is exactly induced by the norm \(\lVert(\vartheta_1,\dots,\vartheta_{n})\rVert\coloneqq\sum_{k=1}^{n}\lVert \vartheta_k\rVert_{L^2}\). Recalling that \(\overline{\bigoplus_{k=1}^{n}\Xi_k}\)is by definition the smallest closed set containing \(\bigoplus_{k=1}^{n}\Xi_k\), the inclusion follows. The reverse inclusion can be checked by hand.}
in particular \(\overline{\Theta}_{\infty}=\overline{\Theta}\), where the closures are taken with respect to the previously introduced norm.
If we show that the map 
\begin{equation}
    J:\vartheta\mapsto\mathbb{E}[\left(H-V_0-(\vartheta\bullet X)_T\right)^2]
\end{equation}
is continuous with respect to the \(\lVert\cdot\rVert\)-induced topology, the 
\cref{prop:finite deep mean variance hedging} follows from \cref{prop:minimization of closure}. The continuity follows from the repetitive application of Cauchy-Schwarz and the observation that in the finite case 
\begin{equation}
    (\vartheta\bullet X)_T=\sum_{k=1}^{n}\vartheta_k(X_{k}-X_{k-1})\leq2\max_{0\leq k\leq n}\max_{\omega\in\Omega}|X_k(\omega)|\sum_{k=1}^n\vartheta_k\;,
\end{equation}
where \(C(X)\coloneqq\max_{0\leq k\leq n}\max_{\omega\in\Omega}|X_k(\omega)|<\infty\) as this \(X_k\) is a finite random variable on the finite sample space \(\Omega\).
Denoting \(Z\coloneqq H-V_0\) for some fixed \(H,V_0\), we find 
\begin{align*}
    &\left|\mathbb{E}[(Z-(\vartheta\bullet X)_T)^2]-\mathbb{E}[(Z-(\vartheta'\bullet X)_T)^2]\right|\\
    &\leq\lVert Z\rVert_{L^2}\lVert(\vartheta-\vartheta')\bullet X\rVert_{L^2}+2C(X)n\lVert \vartheta-\vartheta'\rVert\\
    &\leq C(X)\lVert \vartheta-\vartheta'\rVert +2nC(X)\lVert \vartheta-\vartheta'\rVert\\
    &C(X)\left(1+2n\right)\lVert \vartheta-\vartheta'\rVert\;.
\end{align*}
Hence the map \(J\) is 1-Lipschitz and thus (uniformly) continuous.
\end{appendices}
\end{document}

